% Generated mechanically from frozen input p12/ko/build/ch110/examples_full.tex.
% Do not edit this generated file; edit the bound locale source instead.

% OK, start here.
%

\setcounter{chapter}{109}
\chapter{예시}



\phantomsection
\label{examples-section-phantom}


\section{서론}
\label{examples-section-introduction}

\noindent
이 장에서는 이론을 밝혀 주는 예들을 제시한다.






\section{공집합인 극한}
\label{examples-section-empty-limit}

\noindent
이 예는 Waterhouse가 제시한 것이다. \cite{Waterhouse}를 보라.
$S$를 비가산 집합이라 하자. 각 유한 부분집합
$T \subset S$에 대하여 단사 사상 $T \to \mathbf{N}$들의 집합을 $M_T$라 하자.
$T \subset T' \subset S$가 유한이면 제한 사상 $M_{T'} \to M_T$는
전사이다. 따라서 $S$의 유한 부분집합들로 이루어진 유향 부분순서 집합 위에서
전이 사상이 전사인 역계를 얻는다.
그러나 $\lim M_T = \emptyset$이다. 극한의 원소가 있다면
단사 사상 $S \to \mathbf{N}$을 정의하기 때문이다.




\section{영인 극한}
\label{examples-section-zero-limit}

\noindent
$(S_i)_{i \in I}$를 전이 사상이 전사이고
$\lim S_i = \emptyset$인 비어 있지 않은 집합들의 유향 역계라 하자.
절 \ref{examples-section-empty-limit}을 보라.
$K$를 체라 하고
$$
V_i = \bigoplus\nolimits_{s \in S_i} K
$$
라 놓자. 그러면 $i \geq j$일 때 전이 사상 $V_i \to V_j$는 전사이다.
그러나 $\lim V_i = 0$이다. 실제로 $v = (v_i)$가 극한의 원소이면,
$v_i$의 지지집합은 $\lim T_i \not = \emptyset$을 만족하는 유한 부분집합
$T_i \subset S_i$가 된다. 범주론, 보조정리
\ref{categories-lemma-nonempty-limit}을 보라.

\medskip\noindent
각 $i$에 대하여 모든 기저 벡터 $s \in S_i$를 $1$로 보내는 유일한
$K$-선형 사상 $V_i \to K$를 생각하자. 그 핵을 $W_i \subset V_i$라 하자.
그러면
$$
0 \to (W_i) \to (V_i) \to (K) \to 0
$$
은 유향 집합 $I$ 위의 벡터 공간 역계들의 분할되지 않는 짧은 완전열이다.
따라서 $W_i$는 전이 사상이 전사이고, 극한은 영이며,
$R^1\lim$은 영이 아닌 $K$-벡터 공간들의 유향계이다.




\section{준콤팩트 공간들의 비준콤팩트한 역극한}
\label{examples-section-lim-not-quasi-compact}

\noindent
$\mathbf{N}$은 자연수의 집합을 나타낸다고 하자.
각 정수 $n$에 대하여 $I_n$을 $> n$인 모든 자연수의 집합이라 하자.
$T_n$을 $\mathbf{N}$ 위에서 $\{1\}, \ldots , \{n\}, I_n$을 기저로 갖는
유일한 위상이라 정의하자. $X_n$으로 위상 공간 $(\mathbf{N}, T_n)$을 나타내자.
각 $m < n$에 대하여 항등 사상
$$
f_{n, m} : X_n \longrightarrow X_m
$$
은 연속이다. 분명히 $m < n < p$이면 합성
$f_{p, n} \circ f_{n, m}$은 $f_{p, m}$과 같다. 따라서
$((X_n), (f_{n,m}))$은 준콤팩트 위상 공간들의 유향 역계이다.

\medskip\noindent
$T$를 $\mathbf{N}$ 위의 이산 위상이라 하고 $X$를 $(\mathbf{N}, T)$라 하자.
그러면 모든 정수 $n$에 대하여 항등 사상
$$
f_n : X \longrightarrow X_n
$$
은 연속이다. 이것이 위 유향계의 역극한임을 보이자. $(Y, S)$를 임의의
위상 공간이라 하자. 모든 정수 $n$에 대하여
$$
g_n : (Y, S) \longrightarrow (\mathbf{N}, T_n)
$$
을 연속 사상이라 하자. 모든 $m < n$에 대해
$f_{n,m} \circ g_n = g_m$, 즉 계 $(g_n)$이 위 유향계와 양립한다고 가정하자.
특히 집합 사상 $g_n$들은 모두 하나의 집합 사상
$$
g : Y \longrightarrow \mathbf{N}.
$$
과 같다. 더욱이 모든 정수 $n$에 대하여 $\{n\}$은 $X_n$에서 열려 있으므로
$g^{-1}(\{n\}) = g_n^{-1}(\{n\})$도 $Y$에서 열려 있다.
따라서 집합 사상 $g$는 위상 $S$가 주어진 $Y$에서 위상 $T$가 주어진
$\mathbf{N}$으로의 연속 사상이다. 그러므로 $(X, (f_n))$은 위 유향계의 역극한이다.

\medskip\noindent
그러나 한원소 집합들로 이루어진 무한 열린 피복은 역극한을 갖지 않으므로,
$X$는 명백히 준콤팩트가 아니다.

\begin{lemma}
\label{examples-lemma-lim-not-quasi-compact}
$\mathbf{N}$ 위의 준콤팩트 위상 공간들의 역계 중에서 그 극한이
준콤팩트가 아닌 것이 존재한다.
\end{lemma}

\begin{proof}
위의 논의를 보라.
\end{proof}





\section{섬유곱 위의 구조층}
\label{examples-section-silly}

\noindent
$X, Y, S, a, b, p, q, f$는 『스킴의 유도 범주』의 서론,
절 \ref{perfect-section-kunneth}에서와 같다고 하자. 그림으로 나타내면
$$
\xymatrix{
& X \times_S Y \ar[ld]^p \ar[rd]_q \ar[dd]^f \\
X \ar[rd]_a & & Y \ar[ld]^b \\
& S
}
$$
이다. 그러면 표준 사상
$$
can :
p^{-1}\mathcal{O}_X \otimes_{f^{-1}\mathcal{O}_S} q^{-1}\mathcal{O}_Y
\longrightarrow
\mathcal{O}_{X \times_S Y}
$$
이 존재하지만 일반적으로 동형사상은 아니다.

\medskip\noindent
예를 들어 $S = \Spec(\mathbf{R})$, $X = \Spec(\mathbf{C})$,
$Y = \Spec(\mathbf{C})$라 하자. 그러면
$X \times_S Y = \Spec(\mathbf{C}) \amalg \Spec(\mathbf{C})$는 두 점을 가진
이산 공간이고, 층 $p^{-1}\mathcal{O}_X$, $q^{-1}\mathcal{O}_Y$,
$f^{-1}\mathcal{O}_S$는 각각 $\mathbf{C}$, $\mathbf{C}$, $\mathbf{R}$을
값으로 갖는 상수층이다. 따라서 $can$의 원천은 두 점을 가진 이산 공간 위에서
$\mathbf{C} \otimes_\mathbf{R} \mathbf{C}$를 값으로 갖는 상수층이다.
그러므로 그 대역 단면은 차원이 $8$인 $\mathbf{R}$-벡터 공간인 반면,
$can$의 표적의 대역 단면을 취하면 $\mathbf{C} \times \mathbf{C}$를 얻고,
이는 차원이 $4$인 $\mathbf{R}$-벡터 공간이다.

\medskip\noindent
또 다른 예는 다음과 같다. $k$를 대수적으로 닫힌 체라 하자.
$S = \Spec(k)$, $X = \mathbf{A}^1_k$, $Y = \mathbf{A}^1_k$를 생각하자.
그러면 비어 있지 않은 열린집합
$U \subset X \times_S Y = \mathbf{A}^2_k$에 대하여 그 상
$p(U) \subset X = \mathbf{A}^1_k$와 $q(U) \subset \mathbf{A}^1_k$는 열려 있고,
독자는 다음을 보일 수 있다.
$$
\left(
p^{-1}\mathcal{O}_X \otimes_{f^{-1}\mathcal{O}_S} q^{-1}\mathcal{O}_Y
\right)(U) = \mathcal{O}_X(p(U)) \otimes_k \mathcal{O}_Y(q(U))
$$
이것은 $\mathcal{O}_{X \times_S Y}(U)$와 같지 않다. 여기서 $U$는 기약 곡선
$C$의 여집합이며, 이 곡선은 $X \times_S Y = \mathbf{A}^2_k$ 안에 있고
$p|_C$와 $q|_C$는 모두 상수가 아니다.

\medskip\noindent
일반적인 경우로 돌아가자. $z = (x, y, s, \mathfrak p)$를 $X \times_S Y$의 점이라 하자.
이 점은 『스킴』의 보조정리 \ref{schemes-lemma-points-fibre-product}에서와 같은 점이다.
그러면 $z$에서 줄기를 취했을 때 사상 $can$은 사상
$$
can_z :
\mathcal{O}_{X, x} \otimes_{\mathcal{O}_{S, s}} \mathcal{O}_{Y, y}
\longrightarrow
\mathcal{O}_{X \times_S Y, z}
$$
을 준다. 표적은 원천의 국소화이므로 이는 평탄한 환 준동형이다
(세부 사항은 생략한다. $X$, $Y$, $S$가 아핀인 경우로 환원하라).
원천은 일반적으로 국소환이 아니며, 이 사실은 $can$이 일반적으로
동형사상이 아님을 보는 또 다른 방법을 준다.

\medskip\noindent
더 일반적으로 $\mathcal{O}_X$-가군 $\mathcal{F}$와
$\mathcal{O}_Y$-가군 $\mathcal{G}$가 주어졌다고 하자. 그러면 표준 사상
\begin{align*}
& p^{-1}\mathcal{F} \otimes_{f^{-1}\mathcal{O}_S} q^{-1}\mathcal{G} \\
& =
p^{-1}(\mathcal{F} \otimes_{\mathcal{O}_X} \mathcal{O}_X)
\otimes_{f^{-1}\mathcal{O}_S}
q^{-1}(\mathcal{O}_Y \otimes_{\mathcal{O}_Y} \mathcal{G}) \\
& =
p^{-1}\mathcal{F} \otimes_{p^{-1}\mathcal{O}_X} p^{-1}\mathcal{O}_X
\otimes_{f^{-1}\mathcal{O}_S}
q^{-1}\mathcal{O}_Y \otimes_{q^{-1}\mathcal{O}_Y} q^{-1}\mathcal{G} \\
& \xrightarrow{can}
p^{-1}\mathcal{F} \otimes_{q^{-1}\mathcal{O}_X}
\mathcal{O}_{X \times_S Y}
\otimes_{q^{-1}\mathcal{O}_Y} q^{-1}\mathcal{G}  \\
& =
p^{-1}\mathcal{F} \otimes_{q^{-1}\mathcal{O}_X}
\mathcal{O}_{X \times_S Y}
\otimes_{\mathcal{O}_{X \times_S Y}}
\mathcal{O}_{X \times_S Y}
\otimes_{q^{-1}\mathcal{O}_Y} q^{-1}\mathcal{G}  \\
& =
p^*\mathcal{F} \otimes_{\mathcal{O}_{X \times_S Y}} q^*\mathcal{G}
\end{align*}
이 존재하지만 동형사상인 경우는 드물다.




\section{국소환은 정역이지만 정수적이지 않은 연결 스킴}
\label{examples-section-connected-locally-integral-not-integral}

\noindent
연결이고 모든 국소환은 정역이지만 정수적이지는 않은 아핀 스킴
$X = \Spec(A)$의 예를 들겠다. $X$가 연결이라는 것은 $A$에 자명하지 않은
멱등원이 없다는 뜻이다. 『대수학』의 보조정리
\ref{algebra-lemma-disjoint-decomposition}을 보라.
$X$의 국소환들이 정역이라는 것은 $fg = 0$이 $A$에서 성립할 때마다 $X$의 모든 점이
$f$ 또는 $g$ 중 하나가 영이 되는 근방을 갖는다는 뜻이다.
$A$가 정역이 아닌 한 $X$는 정수적이지 않다
(『성질』의 정의 \ref{properties-definition-integral}).

\medskip\noindent
대략적으로 구성은 다음과 같다. $X_0$를 아핀 평면 위의 십자,
즉 두 좌표축의 합집합이라 하자. 이제 $X_1$을 평면의 부풀리기에서 $X_0$의
(축소된) 완전 역상이라 하자($X_1$은 사슬을 이루는 세 개의
유리 성분을 갖는다). 그런 다음 얻어진 곡면을 $X_1$의 두 특이점에서 부풀리기하고,
$X_2$를 $X_1$의 축소된 역상이라 하며(이 스킴은 다섯 개의 유리 성분을 갖는다),
이를 계속한다. $X$를 이들의 역극한으로 취한다. 이 구성에서 유일한 문제는
부풀리기가 사영 직선을 붙이므로 $X_1$이 아핀이 아니라는 점이다.
대신 아핀 직선을 붙여 이 문제를 고치자(따라서 우리의 스킴은 위에서 설명한
것의 열린 부분집합이 된다).

\medskip\noindent
완전히 대수적인 구성은 다음과 같다. 각 $k \ge 0$에 대하여 $A_k$를
다음 환이라 하자. 그 원소는 다항식 $p_i \in \mathbf{C}[x]$의 모음으로서,
여기서 $i = 0, \ldots, 2^k$이며 $p_i(1) = p_{i + 1}(0)$을 만족하는 것이다.
$X_k = \Spec(A_k)$라 놓자. $X_k$는 서로 횡단적으로 만나 사슬을 이루는
$2^k + 1$개의 아핀 직선의 합집합이다. 환 준동형 $A_k \to A_{k + 1}$을
$$
(p_0, \ldots, p_{2^k})
\longmapsto
(p_0, p_0(1), p_1, p_1(1), \ldots, p_{2^k}),
$$
로 정의하자. 달리 말해 하나 걸러 하나의 다항식은 상수이다. 이로써 $A_k$는
$A_{k + 1}$의 부분환과 동일시된다. $A$를 $A_k$들의 직극한(본질적으로는
그들의 합집합)이라 하고 $X = \Spec(A)$라 놓자. 모든 $k$에 대해 자연스러운
매장 $A_k \to A$, 즉 사상 $X\to X_k$가 있다. 각 $A_k$는 연결이지만
정수적이지 않으므로 $A$도 연결이지만 정수적이지 않다. 이제 $A$의 국소환이
정역임을 보이면 된다.

\medskip\noindent
$f, g \in A$가 $fg = 0$을 만족하고 $x \in X$라 하자. $x$의 근방으로서 $f$나 $g$ 중
하나가 영이 되는 것을 구성하자. $k$를 $f, g \in A_{k - 1}$이 되도록 택한다
($k - 1$이라는 첨자에 주의하라). $y$를 $x$의 $X_k$에서의 상이라 하자.
$y$가, $f$ 또는 $g$를 $\mathcal{O}_{X_k}$의 단면으로 보았을 때 영이 되는 근방을
갖는다고 보이면 충분하다. $y$가 $X_k$의 매끄러운 점, 즉 $2^k + 1$개의
직선 중 오직 하나 위에 놓이면 이는 명백하다. 따라서 $y$가 $2^k$개의
특이점 중 하나여서 $X_k$의 두 성분이 $y$를 지난다고 가정해도 된다.
그러나 이 두 성분 중 하나(홀수 첨자를 가진 성분)에서는 $f$와 $g$가
$X_{k - 1}$ 위의 함수에서 당겨졌으므로 모두 상수이다. 모든 곳에서 $fg = 0$이므로
$f$나 $g$ 중 하나, 이를테면 $f$가 다른 성분에서 영이다.
따라서 필요한 대로 $f$는 두 성분 모두에서 영이다.


\section{완비화가 완비가 아닌 경우}
\label{examples-section-noncomplete-completion}

\noindent
$R$을 환, $\mathfrak m$을 극대 아이디얼이라 하자. 다음 완비화를 생각하자.
$$
R^\wedge = \lim R/\mathfrak m^n.
$$
$R^\wedge$는 극대 아이디얼
$\mathfrak m' = \Ker(R^\wedge \to R/\mathfrak m)$을 갖는 국소환이다.
실제로 $x = (x_n) \in R^\wedge$가 $\mathfrak m'$에 속하지 않으면
$y = (x_n^{-1}) \in R^\wedge$가 $xy = 1$을 만족하므로 $R^\wedge$는 국소환이다.
이는 『대수학』의 보조정리 \ref{algebra-lemma-characterize-local-ring}에 따른다.
$R^\wedge$가 그 극한 위상에 관해서 완비라는 것은 항상 참이다
(『대수학 더 보기』의 절 \ref{more-algebra-section-topological-ring}의 논의를 보라).
그러나 그 밖에는 다음 질문들이 남는다.
\begin{enumerate}
\item $\mathfrak m R^\wedge = \mathfrak m'$인가?
\item $R^\wedge$를 $R^\wedge$-가군으로 볼 때 $\mathfrak m'$-진 완비인가?
\item $R^\wedge$를 $R$-가군으로 볼 때 $\mathfrak m$-진 완비인가?
\end{enumerate}
이 질문들의 답은 모두 부정적이다. 아래 예는 Bart de Smit와
Hendrik Lenstra의 미출판 노트에서 가져왔다. 또한
\cite[Exercise III.2.12]{Bourbaki-CA}와 \cite[Example 1.8]{Yekutieli}를 보라.

\medskip\noindent
$k$를 체라 하고 $R = k[x_1, x_2, x_3, \ldots]$,
$\mathfrak m = (x_1, x_2, x_3, \ldots)$라 하자. 원소 $f$를 $R^\wedge$의 원소로 보면
(다중 첨자 표기법을 써서)
$$
f = \sum a_I x^I
$$
와 같은 무한일 수도 있는 합으로 생각할 수 있다. 여기서 각 $d \geq 0$에 대해
영이 아닌 $a_I$ 중 $|I| = d$인 것은 유한 개뿐이다. 극대 아이디얼
$\mathfrak m' \subset R^\wedge$는 상수항이 영인 $f$들로 이루어진다.
특히 원소
$$
f = x_1 + x_2^2 + x_3^3 + \ldots
$$
는 $\mathfrak m'$에 속하지만 $\mathfrak m R^\wedge$에는 속하지 않는다.
따라서 이 예에서 (1)은 거짓이다. 그런데 (1)이 거짓이면 (3)도 반드시 거짓이다.
$\mathfrak m' = \Ker(R^\wedge \to R/\mathfrak m)$이고 『대수학』의 보조정리
\ref{algebra-lemma-hathat}을 $n = 1$로 적용할 수 있기 때문이다.

\medskip\noindent
마지막으로 $R^\wedge$가 $\mathfrak m'$-진 완비가 아님을 보이자.
$n \geq 1$에 대해 $K_n = \Ker(R^\wedge \to R/\mathfrak m^n)$이라 하자.
그러면 짧은 완전열
$$
0 \to K_n/(\mathfrak m')^n \to R^\wedge/(\mathfrak m')^n \to
R/\mathfrak m^n \to 0
$$
이 있다. 사영 사상 $R^\wedge \to R/\mathfrak m^{n + 1}$은
$(\mathfrak m')^n$을 $\mathfrak m^n/\mathfrak m^{n + 1}$ 위로 보낸다.
따라서 $K_{n + 1} \to K_n/(\mathfrak m')^n$은 전사이다. 그러므로 역계
$\left(K_n/(\mathfrak m')^n\right)$은 전이 사상이 전사이고,
역극한을 취하면 완전열
$$
0 \to \lim K_n/(\mathfrak m')^n \to
\lim R^\wedge/(\mathfrak m')^n \to
\lim R/\mathfrak m^n \to 0
$$
을 얻는다. 이는 『대수학』의 보조정리
\ref{algebra-lemma-Mittag-Leffler}에 따른다.
따라서 $R^\wedge$가 $\mathfrak m'$에 관해 완비일 필요충분조건은
$K_n = (\mathfrak m')^n$이 모든 $n \geq 1$에 대해 성립하는 것이다.

\medskip\noindent
우리 예에서 $R^\wedge$가 $\mathfrak m'$-진 완비가 아님을 보이기 위해
$K_2 = \Ker(R^\wedge \to R/\mathfrak m^2)$가 $(\mathfrak m')^2$와
같지 않음을 보이자. $(\mathfrak m')^2$의 원소는 유한합
\begin{equation}
\label{examples-equation-sum}
\sum\nolimits_{i = 1, \ldots, t} f_i g_i
\end{equation}
으로 나타낼 수 있다. 여기서 $f_i, g_i \in R^\wedge$는 상수항이 영이다.
예를 얻기 위해 다음 성질을 갖는 $z \in K_2$를
다음 꼴로 택하겠다.
$$
z = z_1 + z_2 + z_3 + \ldots
$$
\begin{enumerate}
\item 수열 $1 < d_1 < d_2 < d_3 < \ldots $와
$0 < n_1 < n_2 < n_3 < \ldots$가 존재하여
$z_i \in k[x_{n_i}, x_{n_i + 1}, \ldots, x_{n_{i + 1} - 1}]$이고
이 원소는 차수 $d_i$의 동차식이다.
\item 환 $k[[x_{n_i}, x_{n_i + 1}, \ldots, x_{n_{i + 1} - 1}]]$에서 $z_i$는
합 (\ref{examples-equation-sum}) 중 $t \leq i$인 것으로 나타낼 수 없다.
\end{enumerate}
이는 분명히 $z$가 $(\mathfrak m')^2$에 속하지 않는다는 뜻이다.
관계 (\ref{examples-equation-sum})의 상을 환
$k[[x_{n_i}, x_{n_i + 1}, \ldots, x_{n_{i + 1} - 1}]]$에서 $i$가 충분히 클 때
취하면 모순이 생기기 때문이다.
따라서 모든 $t > 0$에 대해 $d \gg 0$과 정수 $n$이 존재하여,
동차 원소 $z \in k[x_1, \ldots, x_n]$를 그 차수가 $d$가 되도록 택하여, 이 원소를
합 (\ref{examples-equation-sum})으로 나타내는 일이 주어진 $t$와 환 $k[[x_1, \ldots, x_n]]$에서는
불가능하도록 할 수 있음을 보이면 충분하다. $n > 2t$로 잡고 $d > 1$인 수로서
$k$의 표수와 서로소인 임의의 것을 택하여 $z = \sum_{i = 1, \ldots, n} x_i^d$라 놓자.
그러면 아이디얼
$$
(\frac{\partial z}{\partial x_1}, \ldots, \frac{\partial z}{\partial x_n})
=
(dx_1^{d - 1}, \ldots, dx_n^{d - 1})
$$
의 영점 자취는 한 점으로 이루어진다. 한편 Leibniz 법칙에 의해
$$
\frac{\partial ( \sum\nolimits_{i = 1, \ldots, t} f_i g_i ) }{\partial x_j}
\in (f_1, \ldots, f_t, g_1, \ldots, g_t)
$$
이므로 이 편도함수들의 영점 자취는 적어도 다음 집합을 포함한다.
$$
V(f_1, \ldots, f_t, g_1, \ldots, g_t) \subset
\Spec(k[[x_1, \ldots, x_n]]).
$$
그런데 $V(f_1, \ldots, f_t, g_1, \ldots, g_t)$의 차원은
적어도 $n - 2t \geq 1$이므로 모순이다.

\begin{lemma}
\label{examples-lemma-noncomplete-completion}
국소환 $R$과 극대 아이디얼 $\mathfrak m$이 존재하여, $R^\wedge$가
$R$의 $\mathfrak m$에 관한 완비화이고 다음 성질들을 갖는다.
\begin{enumerate}
\item $R^\wedge$는 국소환이지만 그 극대 아이디얼은
$\mathfrak m R^\wedge$와 같지 않다.
\item $R^\wedge$는 완비 국소환이 아니다.
\item $R^\wedge$는 $\mathfrak m$-진으로 완비인 $R$-가군이 아니다.
\end{enumerate}
\end{lemma}

\begin{proof}
위의 논의에서 따른다. 실제로 $R = k[x_1, x_2, x_3, \ldots]$일 때
국소화 $R_{\mathfrak m}$의 완비화는 $R$의 완비화와 같다.
\end{proof}

\section{완비가 아닌 몫}
\label{examples-section-noncomplete-quotient}

\noindent
$k$를 체라 하자. 다음과 같이 놓는다.
$$
R = k[t, z_1, z_2, z_3, \ldots, w_1, w_2, w_3, \ldots, x]/
(z_it - x^iw_i, z_i w_j)
$$
특히 이 환에서 $z_iz_jt = 0$임에 주의하라. 원소 $f$를 $R$에서
임의로 택하면 다음과 같은 유한합으로 유일하게 쓸 수 있다.
$$
f = \sum\nolimits_{i = 0, \ldots, d} f_i x^i
$$
여기서 각 $f_i \in k[t, z_i, w_j]$에는 곱 $z_it$ 또는 $z_iw_j$가
들어간 항이 없다. 더욱이 $f$를 이와 같이 썼을 때 $f \in (x^n)$일
필요충분조건은 $f_i = 0$이 모든 $i < n$에 대해 성립하는 것이다.
따라서 $x$는 영인자가 아니고 $\bigcap (x^n) = 0$이다.
$R^\wedge$를 $R$의 아이디얼 $(x)$에 관한 완비화라 하자.
$R^\wedge$는 $(x)$-진 완비임에 주의하라. 『대수학』의 보조정리
\ref{algebra-lemma-hathat-finitely-generated}를 보라.
위에서 보았듯이 $R^\wedge$의 모든 원소는 무한합
$$
f = \sum\nolimits_{i = 0}^\infty f_i x^i
$$
으로 유일하게 쓸 수 있으며, 각 $f_i \in k[t, z_i, w_j]$에는 곱
$z_it$ 또는 $z_iw_j$가 들어간 항이 없다. 다음 원소를 생각하자.
$$
f = \sum\nolimits_{i = 1}^\infty x^i w_i =
xw_1 + x^2w_2 + x^3w_3 + \ldots
$$
즉 $f_n = w_n$이다. $f \in (t , x^n)$이 모든 $n$에 대해 성립함에
주의하라. 이는 $x^mw_m \in (t)$이 모든 $m$에 대해 성립하기 때문이다.
$f \not \in (t)$임을 주장한다. 이를 보이려고 $tg = f$라고 가정하자.
여기서 $g = \sum g_lx^l$은 위와 같은 표준형으로 쓴 것이다.
$tz_iz_j = 0$이므로 어떤 $g_l$에도 곱 $z_iz_j$가 들어간 항이 없다고
가정해도 된다. $tg$를 표준형으로 만드는 과정을 살피면 다음을 알 수 있다.
임의의 항 $c m$을 $g_l$에서 택하자. 여기서 $c \in k$이고 $m$은
$t, z_i, w_j$에 관한 단항식이다. 그러면 다음과 같이 치환한다.
\begin{enumerate}
\item 단항식 $m$에 어떤 $z_i$도 나타나지 않으면 $ctm$은
$f_l$의 항이다.
\item 단항식 $m$에 어떤 $z_i$가 나타나면 이는 $m = z_i$이고,
$cw_i$는 $f_{l + i}$의 항이다.
\end{enumerate}
$g_0$는 다항식이므로 그 안에는 유한 개의 변수 $z_i$만 나타난다.
$n$을 $z_n$이 $g_0$에 나타나지 않도록 택하자. 그러면 위 규칙에 의해
$w_n$은 $f_n$에 나타나지 않는데, 이는 모순이다. 따라서
$R^\wedge/(t)$는 완비가 아니다. 『대수학』의 보조정리
\ref{algebra-lemma-quotient-complete}를 보라.

\begin{lemma}
\label{examples-lemma-noncomplete-quotient}
환 $R$과 주 아이디얼 $I$, 주 아이디얼 $J$가 존재하며, 그 환은 첫 번째
아이디얼에 관해 완비이지만 $R/J$는 $I$-진 완비가 아니다.
\end{lemma}

\begin{proof}
위의 논의를 보라.
\end{proof}


\section{완비화는 완전하지 않다}
\label{examples-section-completion-not-exact}

\noindent
다음은 간단한 예이다. $R = k[t]$라고 하자. $M^\wedge$는 $R$-가군의
$I = (t)$에 관한 완비화를 나타낸다. $P = K = \bigoplus_{n \in \mathbf{N}} R$이고
$M = \bigoplus_{n \in \mathbf{N}} R/(t^n)$이라 하자. 그러면 짧은 완전열
$0 \to K \to P \to M \to 0$이 있다. 첫 번째 사상은 $t^n$을 곱하는
사상이며, 이는 직합의 $n$번째 성분마다 적용된다. 우리는
$0 \to K^\wedge \to P^\wedge \to M^\wedge \to 0$이 가운데에서 완전하지 않다고
주장한다. 실제로 $\xi = (t^2, t^3, t^4, \ldots) \in P^\wedge$는
$M^\wedge$에서 영으로 가지만 $K^\wedge \to P^\wedge$의 상에는 속하지 않는다.
그 상에 속한다면 $(t, t, t, \ldots)$의 상이어야 하는데, 이는
$K^\wedge$의 원소가 아니기 때문이다.

\medskip\noindent
다음은 더 ``작은'' 예이다. 보조정리 \ref{examples-lemma-noncomplete-quotient}의
상황에서 짧은 완전열 $0 \to J \to R \to R/J \to 0$은 $I$에 관해 완비화한 뒤
완전하지 않다. 실제로 $f \in J$가 생성원이면 $f : R \to J$는 전사이고,
따라서 $R \to J^\wedge$도 전사이다. 그러므로 $J^\wedge \to R$의 상은
$(f) = J$이다. 그러나 $R/J$가 완비가 아니라는 것은 전사
$R \to (R/J)^\wedge$의 핵이 $J$보다 엄밀히 크다는 뜻이다. 『대수학』의
보조정리 \ref{algebra-lemma-completion-generalities}와
\ref{algebra-lemma-quotient-complete}를 보라.
같은 이유로 열 $R \to R \to R/(f) \to 0$도 완비화한 뒤 완전하지 않다.

\begin{lemma}
\label{examples-lemma-completion-not-exact}
\begin{slogan}
일반적으로 완비화는 왼쪽 완전하지도 오른쪽 완전하지도 않다.
\end{slogan}
일반적으로 완비화는 완전 함자가 아니며, 심지어 오른쪽 완전하지도 않다.
이는 $I$가 유한 생성일 때에도 유한 표시 가군들의 범주에서 성립한다.
\end{lemma}

\begin{proof}
위의 논의를 보라.
\end{proof}



\section{완비 가군들의 범주는 아벨 범주가 아니다}
\label{examples-section-non-abelian}

\noindent
$R$을 환이라 하고 $I \subset R$을 유한 생성 아이디얼이라 하자.
$\mathcal{A}$를 $I$-진 완비 $R$-가군들의 범주라 하자. 『대수학』의 정의
\ref{algebra-definition-complete}를 보라. $\varphi : M \to N$을
$\mathcal{A}$의 사상이라 하자. $\varphi$의 $\mathcal{A}$에서의 여핵은
통상적인 여핵의 $I$-진 완비화 $(\Coker(\varphi))^\wedge$이다
($I$가 유한 생성이므로 이 완비화는 완비이다. 『대수학』의 보조정리
\ref{algebra-lemma-hathat-finitely-generated}를 보라).
$K = \Ker(\varphi)$라 하자. $\varphi$가 연속이므로 이는 $M$의 닫힌
부분가군이다. 아래 보조정리 \ref{examples-lemma-closed-is-complete}에 의해 $K$는
$I$-진 완비이고, 따라서 $\varphi$의 $\mathcal{A}$에서의 핵이다.

\begin{lemma}
\label{examples-lemma-closed-is-complete}
$I \subset R$을 환의 아이디얼이라 하자. $M$을 $I$-진 완비 $R$-가군이라
하고 $N \subset M$을 부분가군이라 하자. 다음은 서로 동치이다.
\begin{enumerate}
\item $N$은 $M$에서 닫혀 있다.
\item $N = \bigcap_{n \geq 1} (N + I^n M)$이다.
\item $M/N$은 $I$-진 완비이다.
\end{enumerate}
$I$가 유한 생성이면 이 조건들로부터 $N$도 $I$-진 완비임이 따른다.
\end{lemma}

\begin{proof}
(1)과 (2)의 동치는 『형식 공간』의 보조정리
\ref{formal-spaces-lemma-closed}에서 따르고, (2)와 (3)의 동치는
『대수학』의 보조정리 \ref{algebra-lemma-quotient-complete}이다.
$I$가 유한 생성이고 (1)이 성립한다고 하자. $N^\wedge$를 $I$-진 완비화라
하자. 여기서 완비화되는 가군은 $N$이다. $M$이 $I$-진 완비이므로 포함 사상은
$$
N \to N^\wedge \to M
$$
으로 분해된다. $N^\wedge \to M$은 연속이고, $N$은 $N^\wedge$에서 조밀하며
$N$은 $M$에서 닫혀 있으므로 $N^\wedge \to M$은 $N$을 거쳐 분해된다.
따라서 $N^\wedge = N \oplus C$인 $A$-가군 $C$가 존재한다. 그러므로 $N$은
$I$-진 완비 가군 $N^\wedge$의 직합 인자이다. 『대수학』의 보조정리
\ref{algebra-lemma-hathat-finitely-generated}를 보라
(여기서 $I$가 유한 생성임을 쓴다). 따라서 $N$은 $I$-진 완비이다.
\end{proof}

\noindent
이제 일반적으로 $\Im \not = \Coim$인 현상이 범주 $\mathcal{A}$에서
일어남을 보이는 예를 들자.
$R = \mathbf{Z}_p = \lim_n \mathbf{Z}/p^n\mathbf{Z}$를 $p$-진 정수환으로
잡고 $I = (p)$로 잡는다. 다음 사상을 생각하자.
$$
\text{diag}(1, p, p^2, \ldots) :
\left(\bigoplus\nolimits_{n \geq 1} \mathbf{Z}_p\right)^\wedge
\longrightarrow
\prod\nolimits_{n \geq 1} \mathbf{Z}_p
$$
왼쪽은 직합의 $p$-진 완비화이다. 따라서 왼쪽의 원소는 벡터
$(x_1, x_2, x_3, \ldots)$이며, 그 성분은 $x_i \in \mathbf{Z}_p$이다. 그
$p$-진 값매김 $v_p(x_i) \to \infty$는 $i \to \infty$일 때 성립한다.
이 벡터는 $(x_1, px_2, p^2x_3, \ldots)$로 간다.
따라서 $(1, p, p^2, \ldots)$는 상의 폐포에는 속하지만 상에는 속하지 않는다.
위에서 기술한 핵과 여핵에 의해 이 사상에 대해 $\Im \not = \Coim$임이
분명하다.

\begin{lemma}
\label{examples-lemma-complete-modules-not-abelian}
$R$을 환이라 하고 $I \subset R$을 유한 생성 아이디얼이라 하자.
$I$-진 완비 $R$-가군들의 범주는 핵과 여핵을 갖지만, $R$이 뇌터 환이어도
일반적으로 아벨 범주는 아니다.
\end{lemma}

\begin{proof}
위를 보라.
\end{proof}



\section{유도 완비 가군들의 범주}
\label{examples-section-derived-complete-modules}

\noindent
이 절을 읽기 전에 『대수학 더 보기』의 절
\ref{more-algebra-section-derived-complete-modules}을 읽기 바란다.

\medskip\noindent
$A$를 환이라 하고, $I$를 $A$의 아이디얼이라 하자. 유도 완비 가군들의
범주를 $\mathcal{C}$라 쓰자. 이는 『대수학 더 보기』의 정의
\ref{more-algebra-definition-derived-complete}에서 정의한 범주이다.

\medskip\noindent
$T$를 집합이라 하고, $M_t$, $t \in T$를 유도 완비 가군들의 족이라 하자.
일반적으로 $\bigoplus M_t$는 유도 완비 가군이 아니라고 주장한다.
구체적인 예로 $A = \mathbf{Z}_p$와 $I = (p)$를 잡고
$\bigoplus_{n \in \mathbf{N}} \mathbf{Z}_p$를 생각하자.
$\bigoplus_{n \in \mathbf{N}} \mathbf{Z}_p$에서 그 $p$-진 완비화로 가는
사상은 전사가 아니다. 만일 $\bigoplus_{n \in \mathbf{N}} \mathbf{Z}_p$가
유도 완비라면 이 사상이 전사여야 하므로 이는 유도 완비일 수 없다.
『대수학 더 보기』의 보조정리
\ref{more-algebra-lemma-complete-derived-complete}를 보라.
따라서 포함 함자 $\mathcal{C} \to \text{Mod}_A$는 직합과 (여과) 쌍대극한
어느 것과도 가환하지 않는다.

\medskip\noindent
$I$가 유한 생성이라고 가정하자. 『대수학 더 보기』의 절
\ref{more-algebra-section-derived-complete-modules}의 논의에 의해 범주
$\mathcal{C}$는 임의의 쌍대극한을 갖는다. 그러나 $\mathcal{C}$에서
여과 쌍대극한은 완전하지 않다고 주장한다. 실제로
$A = \mathbf{Z}_p$와 $I = (p)$를 잡자. 다음과 같은 포함들이 있다.
$f_n : \mathbf{Z}_p/p\mathbf{Z}_p \to \mathbf{Z}_p/p^n\mathbf{Z}_p$는
$p$-진 완비 $A$-가군 사이의 사상이며 $p^{n - 1}$을 곱하여 주어진다.
다음 가환 도식들이 존재한다.
$$
\xymatrix{
\mathbf{Z}_p/p\mathbf{Z}_p \ar[r]_{f_n} \ar[d]^1 &
\mathbf{Z}_p/p^n\mathbf{Z}_p \ar[d]_p \\
\mathbf{Z}_p/p\mathbf{Z}_p \ar[r]^{f_{n + 1}} &
\mathbf{Z}_p/p^{n + 1}\mathbf{Z}_p
}
$$
이 포함들의 범주 $\mathcal{C}$에서의 쌍대극한은 사상
$\mathbf{Z}_p/p\mathbf{Z}_p \to 0$을 준다고 주장한다. 실제로 오른쪽 계의
$\text{Mod}_A$에서의 쌍대극한은 $\mathbf{Q}_p/\mathbf{Z}_p$이다.
따라서 $\mathcal{C}$에서의 쌍대극한은
$$
H^0((\mathbf{Q}_p/\mathbf{Z}_p)^\wedge) =
H^0(\mathbf{Z}_p[1]) = 0
$$
이다. 『대수학 더 보기』의 절
\ref{more-algebra-section-derived-complete-modules}에 의한 것이며,
여기서 ${}^\wedge$는 유도 완비화이다. 이로써 주장이 증명되었다.

\begin{lemma}
\label{examples-lemma-derived-complete-modules}
$A$를 환이라 하고 $I \subset A$를 아이디얼이라 하자. 유도 완비 가군들의
범주 $\mathcal{C}$는 아벨 범주이고, 포함 함자
$F : \mathcal{C} \to \text{Mod}_A$는 완전하며 임의의 극한과 가환한다.
$I$가 유한 생성이면 $\mathcal{C}$는 임의의 직합과 쌍대극한을 갖지만
$F$는 일반적으로 이들과 가환하지 않는다. 끝으로 여과 쌍대극한은 일반적으로
$\mathcal{C}$에서 완전하지 않으므로 $\mathcal{C}$는 Grothendieck 아벨 범주가
아니다.
\end{lemma}

\begin{proof}
『대수학 더 보기』의 보조정리
\ref{more-algebra-lemma-derived-complete-modules}와 위의 논의를 보라.
\end{proof}




\section{평탄하지 않은 완비화}
\label{examples-section-nonflat}

\noindent
뇌터인 경우와 달리 환의 아이디얼에 관한 완비화는 항상 평탄하지는 않다.
『대수학 더 보기』의 예 \ref{more-algebra-example-not-glueing-pair}에서
이미 이 현상의 두 예를 보았다. 이 절에서는 두 예를 더 제시한다.

\begin{lemma}
\label{examples-lemma-countable-fg-tensor}
$R$을 환이라 하자. $M$을 가산인 $R$-가군이라 하자. 그러면 $M$이 유한
$R$-가군일 필요충분조건은
$M \otimes_R R^\mathbf{N} \to M^\mathbf{N}$이 전사인 것이다.
\end{lemma}

\begin{proof}
$M$이 유한 가군이면 『대수학』의 명제
\ref{algebra-proposition-fg-tensor}에 의해 이 사상은 전사이다.
역으로 사상이 전사라고 하자. $m_1, m_2, m_3, \ldots$를 $M$의 원소들의
열거라 하자. 텐서곱의 원소
$\sum_{j = 1, \ldots, m} x_j \otimes a_j$가
$(m_n) \in M^\mathbf{N}$으로 간다고 하자. 그러면
$x_1, \ldots, x_m$이 $M$을 $R$ 위에서 생성함을 알 수 있다. 이는
『대수학』의 명제 \ref{algebra-proposition-fg-tensor}의 증명에서와 같다.
\end{proof}

\begin{lemma}
\label{examples-lemma-countable-fp-tensor}
$R$을 가산 환이라 하고 $M$을 가산 $R$-가군이라 하자. $M$이 유한 표시일
필요충분조건은 표준 사상
$M \otimes_R R^\mathbf{N} \to M^\mathbf{N}$이 동형사상인 것이다.
\end{lemma}

\begin{proof}
$M$이 유한 표시 가군이면 『대수학』의 명제
\ref{algebra-proposition-fp-tensor}에 의해 사상은 동형사상이다.
역으로 사상이 동형사상이라 하자. 보조정리
\ref{examples-lemma-countable-fg-tensor}에 의해 $M$은 유한 가군이다.
전사 $R^{\oplus m} \to M$을 택하고 그 핵을 $K$라 하자.
$K$는 $R^{\oplus m}$의 부분가군이므로 가산이다. 『대수학』의 명제
\ref{algebra-proposition-fp-tensor}의 증명처럼 논증하면
$K \otimes_R R^\mathbf{N} \to K^\mathbf{N}$이 전사임을 알 수 있다.
보조정리 \ref{examples-lemma-countable-fg-tensor}에 의해 $K$는 유한 $R$-가군이고,
따라서 $M$은 유한 표시이다.
\end{proof}

\begin{lemma}
\label{examples-lemma-countable-coherent}
$R$을 가산 환이라 하자. $R$이 연접 환일 필요충분조건은
$R^\mathbf{N}$이 평탄 $R$-가군인 것이다.
\end{lemma}

\begin{proof}
$R$이 연접이면 $R^\mathbf{N}$은 평탄 가군이다(『대수학』의 명제
\ref{algebra-proposition-characterize-coherent}). $R^\mathbf{N}$이
평탄이라고 하자. 유한 생성 아이디얼
$I \subset R$을 택하자. $I$가 $R$-가군으로서 유한 표시임을 보이면 된다.
구체적으로 사상 $I \otimes_R R^\mathbf{N} \to R^\mathbf{N}$은
$R^\mathbf{N}$이 평탄이므로 단사이고, 그 상은 $I^\mathbf{N}$이다
(보조정리 \ref{examples-lemma-countable-fg-tensor}).
그러므로 보조정리 \ref{examples-lemma-countable-fp-tensor}에서 결론이 따른다.
\end{proof}

\noindent
환 $R$이 가산이라고 하자. $R[[x]]$는 $R^\mathbf{N}$과 $R$-가군으로서
동형이다. 따라서 보조정리 \ref{examples-lemma-countable-coherent}에 의해
$R \to R[[x]]$가 평탄일 필요충분조건은 $R$이 연접인 것이다.
연접이 아닌 가산 환은 많이 있는데, 예를 들면 다음과 같다.
$$
R = k[y, z, a_1, b_1, a_2, b_2, a_3, b_3, \ldots]/
(a_1 y + b_1 z, a_2 y + b_2 z, a_3 y + b_3 z, \ldots)
$$
여기서 $k$는 가산 체이다. 이 환은 아이디얼 $(y, z)$가 환 $R$에서
유한 표시 $R$-가군이 아니므로 연접이 아니다. 또한 $R[[x]]$는
$R[x]$의 주 아이디얼 $(x)$에 관한 완비화이다.

\begin{lemma}
\label{examples-lemma-completion-polynomial-ring-not-flat}
완비화 $R[[x]]$가 $R[x]$의 $(x)$에서의 완비화이고 $R$ 위에서 평탄하지
않아, 더더욱 $R[x]$ 위에서도 평탄하지 않은 환이 존재한다.
\end{lemma}

\begin{proof}
위의 논의를 보라.
\end{proof}

\noindent
어떤 환 $R$에 대해서는 $R[[x]]$가 $R$ 위에서 평탄하지만 $R[[x]]$가
$R[x]$ 위에서는 평탄하지 않다. Badam Baplan의
\href{https://math.stackexchange.com/users/164860/badam-baplan}{이 글}을 보라.
구체적으로 $R$을 값매김환이라 하자. 그러면 $R$은 연접이고
(『대수학』의 예 \ref{algebra-example-valuation-ring-coherent}), 따라서
『대수학』의 명제 \ref{algebra-proposition-characterize-coherent}에 의해
$R[[x]]$는 $R$ 위에서 평탄이다. 한편 다음 보조정리가 있다.

\begin{lemma}
\label{examples-lemma-almost-integral-when-powerseries-flat}
$R$을 분수체 $K$를 갖는 정역이라 하자. $R[[x]]$가 $R[x]$ 위에서 평탄이면,
$R$이 정규일 필요충분조건은 $R$이 완전 정규인 것이다
(『대수학』의 정의 \ref{algebra-definition-almost-integral}).
\end{lemma}

\begin{proof}
$\alpha \in K$와 영이 아닌 $r \in R$가 $r \alpha^n \in R$을 모든
$n \geq 1$에 대해 만족한다고 하자. $f = \sum r \alpha^{n - 1} x^n$을
$R[[x]]$의 원소로 생각하자. $\alpha = a/b$로 쓰자. 여기서
$a, b \in R$이고 $b$는 영이 아니다. 그러면 $(a x - b)f = -rb$이다.
따라서 $rb$는 $(ax - b)R[[x]]$의 원소이다.
$S = \{h \in R[x] : h(0) = 1\}$라 하자. 이는 곱셈적 부분집합이고
$R[x] \to R[[x]]$의 평탄성으로부터
$S^{-1}R[x] \to R[[x]]$가 충실히 평탄임이 따른다(세부 사항은 생략한다.
『대수학』의 보조정리 \ref{algebra-lemma-ff-rings}을 사용하라). 따라서
$$
S^{-1}R[x]/(ax - b)S^{-1}R[x] \to R[[x]]/(ax - b)R[[x]]
$$
는 단사이다. 그러므로 등식 $h rb = (ax - b) g$을 만족하는 어떤
$h \in S$와 $g \in R[x]$가 존재한다.
$h = 1 + h_1 x + \ldots + h_d x^d$로 쓰면
$$
1 + h_1 x + \ldots + h_d x^d = (1/r)(\alpha x - 1)g
$$
를 얻는다. $K[x]$에서의 이 인수분해는 $K[x^{-1}]$에서 대응하는
인수분해를 주므로, 원하는 대로 $\alpha$는 $R$에 계수를 갖는 일계수
다항식의 근이다.
\end{proof}

\begin{lemma}
\label{examples-lemma-completion-polynomial-ring-not-flat-bis}
$R$이 차원 $> 1$인 값매김환이면 $R[[x]]$는 $R$ 위에서는 평탄하지만
$R[x]$ 위에서는 평탄하지 않다.
\end{lemma}

\begin{proof}
$R$이 완전 정규가 아님을 보이면 위 논증으로 충분하다(값매김환은 정규이다.
『대수학』의 보조정리 \ref{algebra-lemma-valuation-ring-normal}을 보라).
소 아이디얼들의 사슬 $\mathfrak p \subset \mathfrak m \subset R$을 택하고,
영이 아닌 $x \in \mathfrak p$와
$y \in \mathfrak m \setminus \mathfrak p$를 택하자.
$x y^{-n} \in R$은 모든 $n \geq 1$에 대해 성립한다. 그렇지 않으면
$y^n/x \in R$인데, $y \not \in \mathfrak p$이므로 이는 불가능하다.
따라서 $1/y$는 $R$ 위에서 거의 정수적이지만 $R$에는 속하지 않는다.
\end{proof}

\noindent
다음으로 국소화의 완비화가 평탄하지 않은 예를 구성하자. 환
$$
R = k[y, z, a_1, a_2, a_3, \ldots]/(ya_i, a_i a_j)
$$
을 생각하자. $f \in R$을 $z$의 잉여류라 하자. 환 준동형
\begin{equation}
\label{examples-equation-nonflat}
R[[x]] \longrightarrow R_f[[x]]
\end{equation}
은 평탄하지 않다. $I$를 $y : R[[x]] \to R[[x]]$의 핵이라 하자.
$g$가 $I$의 전형적인 원소이면 $g = \sum g_{n, m} a_mx^n$ 꼴이다.
여기서 $g_{n, m} \in k[z]$이고 고정된 $n$마다 영이 아닌 $g_{n, m}$은
유한 개뿐이다. $J$를 $y : R_f[[x]] \to R_f[[x]]$의 핵이라 하자.
$J \not = I R_f[[x]]$이다. 그렇지 않다면 등식
$$
\sum z^{-n} a_n x^n = \sum\nolimits_{i = 1, \ldots, m} h_i g_i
$$
이 어떤 $m \geq 1$, $g_i \in I$, $h_i \in R_f[[x]]$에 대해 성립할 것이다.
$h_i = \bar h_i \bmod (y, a_1, a_2, a_3, \ldots)$라 쓰자. 여기서
$\bar h_i \in k[z, 1/z][[x]]$이다. $a_n$의 계수를 보고 위에서 기술한
원소 $g_i$를 사용하면
$$
z^{-n} x^n = \sum \bar h_i \bar g_{i, n}
$$
을 얻는데, 여기서 $\bar g_{i, n} \in k[z][[x]]$이다. 이는 모든
$z^{-n}x^n$이 유한 $k[z][[x]]$-가군, 곧 원소 $\bar h_i$들이 생성하는
가군에 포함됨을 뜻한다. $k[z][[x]]$가 뇌터 환이므로
$R[z][[x]]$-가군으로 보았을 때 $k[z, 1/z][[x]]$ 안에서
$1, z^{-1}x, z^{-2}x^2, \ldots$가 생성하는 부분가군은 유한 생성이다.
『대수학』의 보조정리 \ref{algebra-lemma-characterize-integral-element}에 의해
$z^{-1}x$가 $k[z][[x]]$ 위에서 정수적이라는 결론이 나오는데, 이는 불가능하다.
한편 (\ref{examples-equation-nonflat})이 평탄이면 $J = IR_f[[x]]$를 얻게 된다.
실제로 완전열 $0 \to I \to R[[x]] \xrightarrow{y} R[[x]]$을
$R_f[[x]]$와 텐서하면 된다.

\begin{lemma}
\label{examples-lemma-nonflat-completion-localization}
환 $A$가 주 아이디얼 $I$에 관해 완비이고, 원소 $f \in A$가 주어져서
$I$-진 완비화 $A_f^\wedge$가 $A_f$의 완비화이면서 $A$ 위에서 평탄하지
않은 경우가 존재한다.
\end{lemma}

\begin{proof}
$A = R[[x]]$, $I = (x)$로 놓고 $R_f[[x]]$가 $R[[x]]_f$의 완비화임을 관찰하라.
\end{proof}

\section{준연접 가군들의 비아벨 범주}
\label{examples-section-nonabelian-QCoh}

\noindent
『스택 위의 층』의 절 \ref{stacks-sheaves-section-quasi-coherent}에서
$\Sch$ 위의 준군 범주에서 올이 있는 범주 위의 준연접 가군들의 범주를
정의했다. 『스택 위의 층』의 절
\ref{stacks-sheaves-section-quasi-coherent-algebraic-stacks}에서
대수적 스택에 대해서는 이 범주가 아벨 범주임을 보이지만,
이 절에서는 형식적 대수공간에 대해서는 그렇지 않음을 보인다.

\medskip\noindent
$\mathbf{Z}_p$를 $p$-진 위상을 갖는 위상환으로 보자.
$X = \text{Spf}(\mathbf{Z}_p)$라 하자. 『형식 공간』의 정의
\ref{formal-spaces-definition-affine-formal-spectrum}를 보라.
그러면 $X$는 $(\Sch/\mathbf{Z})_{fppf}$ 위의 집합층이고, 준집합들의
스택 $\mathcal{X}$를 준다. 『스택』의 보조정리
\ref{stacks-lemma-when-stack-in-sets}를 보라. 따라서 『스택 위의 층』의 절
\ref{stacks-sheaves-section-quasi-coherent-algebraic-stacks}의 논의를 적용할 수 있다.

\medskip\noindent
$\mathcal{F}$를 $\mathcal{X}$ 위의 준연접 가군이라 하자.
$X = \colim \Spec(\mathbf{Z}/p^n\mathbf{Z})$이므로 『스택 위의 층』의
보조정리 \ref{stacks-sheaves-lemma-quasi-coherent}에 의해 $\mathcal{F}$는
다음 조건의 열 $(\mathcal{F}_n)$로 주어진다.
\begin{enumerate}
\item $\mathcal{F}_n$은 $\Spec(\mathbf{Z}/p^n\mathbf{Z})$ 위의 준연접 가군이다.
\item 전이 사상은 동형사상
$\mathcal{F}_n = \mathcal{F}_{n + 1}/p^n\mathcal{F}_{n + 1}$을 준다.
\end{enumerate}
가군으로 바꾸면 $\mathcal{F}$는 계 $(M_n)$에 대응한다. 각 $M_n$은
$p^n$에 의해 소멸하는 아벨 군이고 전이 사상은 동형사상
$M_n = M_{n + 1}/p^n M_{n + 1}$을 유도한다. 이때
$M = \lim M_n$은 $p$-진 완비 가군이고 $M_n = M/p^n M$이다. 『대수학』의
보조정리 \ref{algebra-lemma-limit-complete}를 보라. 따라서 $X$ 위의
준연접 가군 범주는 $p$-진 완비 아벨 군들의 범주와 동치이다.
이 범주는 아벨 범주가 아니다. 절 \ref{examples-section-non-abelian}을 보라.

\begin{lemma}
\label{examples-lemma-quasi-coherent-not-abelian}
준연접\footnote{여기서는 위에서 정의한 의미의 준연접 가군을 뜻한다.
Stacks project의 설정 방식 때문에 이것이 실제로 올바른 정의이다. 위에서
보았듯이 우리의 정의는 $\text{Spf}(A)$ 위의 준연접 가군을 완비 $A$-가군으로
정의하는 소박한 정의와 일치한다.} 가군의 범주는 형식적 대수공간 $X$에
대해 일반적으로 아벨 범주가 아니며, 이는 $X$가 뇌터 아핀 형식적
대수공간일 때에도 그러하다.
\end{lemma}

\begin{proof}
위의 논의를 보라.
\end{proof}






\section{국소적으로 분할되지만 분할되지 않는 열}
\label{examples-section-nonsplit-locally-split}

\noindent
짧은 완전열
$$
0 \to M \to
\bigoplus\nolimits_{p\text{ prime}} \mathbf{Z}_{(p)} \to \mathbf{Q} \to 0
$$
을 생각하자. 여기서 $M$은 오른쪽 사상의 핵이고, 이 사상은 각 직합
성분에서 포함 $\mathbf{Z}_{(p)} \to \mathbf{Q}$를 사용한다. 이
$\mathbf{Z}$-가군의 열은 영이 아닌 준동형
$\mathbf{Q} \to \mathbf{Z}_{(p)}$이 존재하지 않으므로 분할되지 않는다.
한편 임의의 소수 $p$에서 국소화하면 이 열은 분할된다.

\begin{lemma}
\label{examples-lemma-nonsplit-locally-split}
어떤 환 $R$과 분할되지 않는 가군열이 존재하여, 그 열은 Zariski 국소적으로
분할된다.
\end{lemma}

\begin{proof}
위의 논의를 보라.
\end{proof}





\section{정칙열과 밑변환}
\label{examples-section-regular-base-change}

\noindent
환 $R$과 그 안의 정칙열 $(x, y, z)$를 구성하여, 영이 아닌 원소
$\delta \in R/zR$가 존재하고 $x\delta = y\delta = 0$이 되게 하겠다.

\medskip\noindent
이 예를 구성하기 위해 먼저 특이한 가군 $E$를 환 $k[x, y, z]$ 위에서 구성한다.
여기서 $k$는 임의의 체이다. 즉 $E$는 다음 도식과 같은 밀어내기이다.
$$
\xymatrix{
\frac{xk[x, y, z, y^{-1}]}{xyk[x, y, z]} \ar[r] \ar[d]^{z/x} &
\frac{k[x, y, z, x^{-1}, y^{-1}]}{yk[x, y, z, x^{-1}]} \ar[r] \ar[d] &
\frac{k[x, y, z, x^{-1}, y^{-1}]}{yk[x, y, z, x^{-1}] + xk[x, y, z, y^{-1}]}
\ar[d] \\
\frac{k[x, y, z, y^{-1}]}{yzk[x, y, z]} \ar[r] &
E \ar[r] &
\frac{k[x, y, z, x^{-1}, y^{-1}]}{yk[x, y, z, x^{-1}] + xk[x, y, z, y^{-1}]}
}
$$
두 행은 짧은 완전열이다(조판 문제 때문에 양 끝의 영을 생략했다).
$E$를 달리 기술하면
$$
E = \{(f, g) \mid f \in k[x, y, z, x^{-1}, y^{-1}],
g \in k[x, y, z, y^{-1}] \}/\sim
$$
이고, $(f, g) \sim (f', g')$일 필요충분조건은 어떤
$h \in k[x, y, z, y^{-1}]$가 존재하여
$$
f = f' + xh \bmod yk[x, y, z, x^{-1}], \quad
g = g' - zh \bmod yzk[x, y, z]
$$
인 것이다. 다음을 주장한다. (a) $x : E \to E$는 단사이다. (b)
$y : E/xE \to E/xE$는 단사이다. (c) $E/(x, y)E = 0$이다. (d) 영이 아닌
원소 $\delta \in E/zE$가 존재하여 $x\delta = y\delta = 0$이다.

\medskip\noindent
(a)를 보이자. $(f, g)$가 $E$의 원소를 주고 $(xf, xg) \sim 0$이라고 하자.
그러면 어떤 $h \in k[x, y, z, y^{-1}]$가 존재하여
$xf + xh \in yk[x, y, z, x^{-1}]$이고
$xg - zh \in yzk[x, y, z]$이다. $h = \sum a_{i, j, k}x^iy^jz^k$에서
$j \leq 0$인 단항식만 나타난다고 가정해도 된다.
$xg - zh \in yzk[x, y, z]$이므로 $i > 0$인 항만 나타난다. 즉
$h = xh'$이고, 여기서 $h' \in k[x, y, z, y^{-1}]$이다.
$(f, g) \sim (f + xh', g - zh')$이므로 $g = 0$, $h = 0$이라 가정할 수 있다.
이 경우 $xf \in yk[x, y, z, x^{-1}]$에서
$f \in yk[x, y, z, x^{-1}]$이 따르므로 $(f, g) \sim 0$이다.
따라서 $x : E \to E$는 단사이다.

\medskip\noindent
$x$를 곱하는 사상은
$\frac{k[x, y, z, x^{-1}, y^{-1}]}{yk[x, y, z, x^{-1}]}$ 위에서
동형사상이므로 $E/xE$는
$$
\frac{k[x, y, z, y^{-1}]}{
yzk[x, y, z] + xk[x, y, z, y^{-1}] + zk[x, y, z, y^{-1}]}
=
\frac{k[x, y, z, y^{-1}]}{xk[x, y, z, y^{-1}] + zk[x, y, z, y^{-1}]}
$$
와 동형이다. 따라서 $y$를 곱하는 사상은 $E/xE$ 위에서 동형사상이다.
이로부터 (b)와 (c)가 분명히 따른다.

\medskip\noindent
$e \in E$를 $(1, 0)$의 동치류라 하자. $e \in zE$라고 가정하면
$f \in k[x, y, z, x^{-1}, y^{-1}]$,
$g \in k[x, y, z, y^{-1}]$, 그리고 $h \in k[x, y, z, y^{-1}]$가 존재하여
$$
1 + zf + xh \in yk[x, y, z, x^{-1}], \quad
0 + zg - zh \in yzk[x, y, z].
$$
이는 불가능하다. 단항식 $1$은 $zf$에도 $xh$에도 나타날 수 없다.
한편 $ye = 0$이고 $xe = (x, 0) \sim (0, -z) = z(0, -1)$이다. 따라서
$\delta$를 $e$의 $E/zE$에서의 합동류로 두면 (d)를 얻는다.

\begin{lemma}
\label{examples-lemma-strange-regular-sequence}
국소환 $R$과 그 극대 아이디얼 안의 정칙열 $x, y, z$가 존재하여,
영이 아닌 원소 $\delta \in R/zR$가 존재하고 $x\delta = y\delta = 0$이다.
\end{lemma}

\begin{proof}
$R = k[x, y, z] \oplus E$라 놓자. 여기서 $E$는 위에서 구성한 가군을 제곱 영
아이디얼로 본 것이다. 그러면 $x, y, z$가 $R$의 정칙열이고,
$\delta \in E/zE \subset R/zR$가 원하는 성질의 원소를 줌은 분명하다.
국소 예를 얻으려면 $R$을 극대 아이디얼 $\mathfrak m = (x, y, z, E)$에서
국소화한다. $x, y, z$는 국소화가 완전하므로 여전히 정칙열이고,
$\delta$는 $\mathfrak m$에서 지지되므로 여전히 영이 아니다.
\end{proof}

\begin{lemma}
\label{examples-lemma-base-change-regular-sequence}
국소환의 국소 준동형 $A \to B$와 정칙열 $x, y$가 존재한다. 이 정칙열은
$B$의 극대 아이디얼 안에 있고, $B/(x, y)$는 $A$ 위에서 평탄이지만 그 상
$\overline{x}, \overline{y}$, 즉 $x, y$의 $B/\mathfrak m_AB$에서의 상들은 정칙열은커녕
Koszul-정칙열도 아니다.
\end{lemma}

\begin{proof}
$A = k[z]_{(z)}$라 놓고 $B = (k[x, y, z] \oplus E)_{(x, y, z, E)}$라 하자.
보조정리 \ref{examples-lemma-strange-regular-sequence}의 증명에 의해 $x, y, z$는
$B$의 정칙열이다. 따라서 $x, y$는 $B$의 정칙열이고 $B/(x, y)$는
꼬임 없는 $A$-가군이므로 평탄이다. 한편 영이 아닌 원소
$\delta \in B/\mathfrak m_AB = B/zB$가 존재하고
$\overline{x}, \overline{y}$에 의해 소멸한다. 따라서
$H_2(K_\bullet(B/\mathfrak m_AB, \overline{x}, \overline{y})) \not = 0$이다.
그러므로 $\overline{x}, \overline{y}$는 Koszul-정칙열이 아니며, 특히
정칙열도 아니다. 『대수학 더 보기』의 보조정리
\ref{more-algebra-lemma-regular-koszul-regular}를 보라.
\end{proof}


\section{무한 차원 뇌터 환}
\label{examples-section-Noetherian-infinite-dimension}

\noindent
『대수학』의 명제 \ref{algebra-proposition-dimension}에서 보았듯이 뇌터
국소환은 유한 차원이다. 그러나 무한 차원 뇌터 환도 존재한다.
\cite[Appendix, Example 1]{Nagata}를 보라.

\medskip\noindent
$k$를 체라 하고 다음 환을 생각하자.
$$
R = k[x_1, x_2, x_3, \ldots ].
$$
$\mathfrak p_i = (x_{2^{i - 1}}, x_{2^{i - 1} + 1}, \ldots, x_{2^i - 1})$라
하자. 여기서 $i = 1, 2, \ldots$이고, 이들은 $R$의 소 아이디얼이다.
$S$를 다음 곱셈적 부분집합이라 하자.
$$
S = \bigcap\nolimits_{i \geq 1} (R \setminus \mathfrak p_i).
$$
환 $A = S^{-1}R$을 생각하자. 다음을 주장한다.
\begin{enumerate}
\item $A$의 극대 아이디얼들은 $\mathfrak m_i = \mathfrak p_iA$이다.
\item $A_{\mathfrak m_i} = R_{\mathfrak p_i}$이고, 이는 차원 $2^i$인
뇌터 국소환이다.
\item 환 $A$는 뇌터 환이다.
\end{enumerate}
따라서 이것이 찾던 예임은 분명하다. 세부 사항은 생략한다.




\section{완비화가 비축소인 국소환}
\label{examples-section-local-completion-nonreduced}

\noindent
『대수학』의 예 \ref{algebra-example-bad-dvr-char-p}에서 표수 $p$의 뇌터 국소
정역 $R$로서 차원이 $1$이고 완비화가 비축소인 예를 주었다. 이 절에서는
표수 영에서 비슷한 환을 주는 \cite[Proposition 3.1]{Ferrand-Raynaud}의
예를 제시한다.

\medskip\noindent
$\mathbf{C}\{x\}$를 복소수체 $\mathbf{C}$ 위의 수렴 멱급수환이라 하자.
모든 멱급수의 환 $\mathbf{C}[[x]]$은 그 완비화이다.
$K = \mathbf{C}\{x\}[1/x]$를 수렴 로랑 급수체라 하자. 대수적 미분들의
$K$-가군 $\Omega_{K/\mathbf{C}}$는 $K$의 $\mathbf{C}$ 위 미분들로 이루어져
있으며 무한 차원 $K$-벡터 공간이다(증명은 생략한다).
$f_n \in x\mathbf{C}\{x\}$, $n \geq 1$을 택하여
$
\text{d}x, \text{d}f_1, \text{d}f_2, \ldots
$
가 $\Omega_{K/\mathbf{C}}$의 한 기저의 일부가 되게 할 수 있다. 따라서
$\mathbf{C}$-미분
$$
D : \mathbf{C}\{x\} \longrightarrow \mathbf{C}((x))
$$
이 존재하여 $D(x) = 0$이고 $D(f_i) = x^{-n}$이다. 다음과 같이 놓자.
$$
A = \{f \in \mathbf{C}\{x\} \mid D(f) \in \mathbf{C}[[x]]\}
$$
다음을 주장한다.
\begin{enumerate}
\item $\mathbf{C}\{x\}$는 $A$ 위에서 정수적이다.
\item $A$는 국소 정역이다.
\item $\dim(A) = 1$이다.
\item $A$의 극대 아이디얼은 $x$와 $xf_1$로 생성된다.
\item $A$는 뇌터 환이다.
\item $A$의 완비화는 $\mathbf{C}[[x]]$ 위의 쌍대수환과 같다.
\end{enumerate}
쌍대수환은 비축소이므로 $A$가 원하는 예를 준다.

\medskip\noindent
$0 \not = f \in x\mathbf{C}\{x\}$이면 $D(f) = h/f^n$으로 쓸 수 있다.
여기서 어떤 $n \geq 0$과 $h \in \mathbf{C}[[x]]$를 취한다. 따라서
$D(f^{n + 1}/(n + 1)) \in \mathbf{C}[[x]]$이고
$D(f^{n + 2}/(n + 2)) \in \mathbf{C}[[x]]$이다. 그러므로
$f^{n + 1}, f^{n + 2} \in A$이다! 특히 (1)이 성립한다. 또한 $A$의 분수체는
$\mathbf{C}\{x\}$의 분수체와 같다. 더 나아가 $A \cap x\mathbf{C}\{x\}$가
$A$의 비가역원들의 집합이므로 $A$는 차원 $1$인 국소 정역이다. (4)를 보이면
$A$가 뇌터 환임이 따른다(증명은 생략한다).
$f \in A \cap x\mathbf{C}\{x\}$라 하고 $D(f) = h$,
$h \in \mathbf{C}[[x]]$라 쓰자. $h = c + xh'$로 쓰되
$c \in \mathbf{C}$, $h' \in \mathbf{C}[[x]]$라 하자. 그러면
$D(f - cxf_1) = c + xh' - c = xh'$이다. 한편
$f - cxf_1 = xg$인 $g \in \mathbf{C}\{x\}$가 있고, 위 계산에 의해
$D(g) = h' \in \mathbf{C}[[x]]$이므로 $g \in A$이다. 따라서 원하는 대로
$f = cxf_1 + xg \in (x, xf_1)$이다.

\medskip\noindent
마지막으로 왜 $A$의 완비화가 비축소인지 보자. $\hat A$로 $A$의 완비화를
나타내자. 이는 $\mathbf{C}[[x]]$, 즉 $\mathbf{C}\{x\}$의 완비화로 전사하는데,
$x \in A$이기 때문이다. 이 사상을 $\psi : \hat A \to \mathbf{C}[[x]]$라 하자.
위에서 $\mathfrak m_A = (x, xf_1)$임을 보았고, 간단한 계산으로
$D(\mathfrak m_A^n) \subset (x^{n - 1})$이다. 따라서
$D : A \to \mathbf{C}[[x]]$는 연속이고, 연속 미분
$\hat D : \hat A \to \mathbf{C}[[x]]$를 $\psi$ 위에서 유도한다. 그러므로 환 준동형
$$
\psi + \epsilon \hat D :
\hat A
\longrightarrow
\mathbf{C}[[x]][\epsilon].
$$
을 얻는다. $\hat A$는 일차원 뇌터 완비 국소환이므로 이 화살표가 전사임을
보이면 $\hat A$가 비축소임이 따른다. 실제로 이 사상은 동형사상이지만
그 확인은 생략한다. 부분환 $\mathbf{C}[x]_{(x)} \subset A$는 완비화 위에서
$i : \mathbf{C}[[x]] \to \hat A$를 유도하며 $i \circ \psi = \text{id}$이고
$D \circ i = 0$이다(구성상 $D(x) = 0$이다). 원소 $x^nf_n \in A$를 생각하면
$$
(\psi + \epsilon D)(x^nf_n) = x^n f_n + \epsilon
$$
이고, 이는 모든 $n \geq 1$에 대해 성립한다. 이들 관찰에서 전사성이 쉽게 따른다.





\section{완비화가 비축소인 또 다른 국소환}
\label{examples-section-another-local-completion-nonreduced}

\noindent
이 절에서는 주 아이디얼에 관해 완비인 $2$차원 뇌터 국소 정역으로서,
그 국소화를 다시 완비화하면 비축소가 되는 예를 만든다.

\medskip\noindent
$p$를 소수라 하자. $k$를 표수 $p$인 체라 하자. 이때 $k$는 그 부분체
$k^p$, 즉 $p$제곱들의 부분체 위에서 무한 차수라고 하자. 예를 들어
$k = \mathbf{F}_p(t_1, t_2, t_3, \ldots)$로 잡을 수 있다. 환
$$
A =
\left\{
\begin{matrix}
\sum a_{i, j} x^iy^j \in k[[x, y]] \text{ 가운데 모든 }n \geq 0
\text{에 대해 다음이 성립하는 것들} \\
[k^p(a_{n, n}, a_{n, n + 1}, a_{n + 1, n},
a_{n, n + 2}, a_{n + 2, n}, \ldots) : k^p] < \infty
\end{matrix}
\right\}
$$
을 생각하자. 집합으로는
$$
k^p[[x, y]] \subset A  \subset k[[x, y]]
$$
이다. $f$를 $A$의 임의의 원소라 하자. 그러면 이 원소를 급수
$$
f = f_0 + f_1 xy + f_2 (xy)^2 + f_3 (xy)^3 + \ldots
$$
로 유일하게 쓸 수 있으며
$$
f_n = a_{n, n} + a_{n, n + 1} y + a_{n + 1, n} x +
a_{n, n + 2} y^2 + a_{n + 2, n} x^2 + \ldots
$$
이다. $A$를 정의한 식의 조건은 $f_n$의 계수들이 $k^p$의 유한 확대를
생성한다는 뜻이다. 이 표현에서 $A$는 $k^p[[x, y]]$-부분대수로서
$k[[x, y]]$ 안에 있고, 아이디얼 $xy$에 관해 완비임이 분명하다. 또한
$$
A/xy A = C \times_k D
$$
이다. 여기서 $k^p[[x]] \subset C \subset k[[x]]$와
$k^p[[y]] \subset D \subset k[[y]]$는 『대수학』의 예
\ref{algebra-example-bad-dvr-char-p}에 나오는 멱급수의 부분환이다.
따라서 $C$와 $D$는 이산 값매김환이고 $A/ xy A$는 뇌터 환이다.
『대수학』의 보조정리 \ref{algebra-lemma-completion-Noetherian}에 의해
$A$도 뇌터 환이다. $\dim(k[[x, y]]) = 2$이므로 『대수학』의 보조정리
\ref{algebra-lemma-integral-sub-dim-equal}에 의해 $\dim(A) = 2$이다.

\medskip\noindent
멱급수 $f = \sum a_i x^i$를 택하여 $k^p(a_0, a_1, a_2, \ldots)$가
$k^p$ 위에서 무한 차수가 되게 하자. 그러면 $f \not \in A$이지만
$f^p \in A$이다. 다음과 같이 놓자.
$$
B = A[f] \subset k[[x, y]]
$$
$B$는 $A$ 위에서 유한이므로 $B$가 뇌터 환임을 알 수 있다. 또한 $B$는
$xy$가 생성하는 아이디얼에 관해 완비이다. 『대수학』의 보조정리
\ref{algebra-lemma-completion-tensor}를 보라. 실제로 $B$는 $A$ 위에서
$1, f, f^2, \ldots, f^{p - 1}$을 기저로 하는 자유 가군이다. 증명은 생략한다.

\medskip\noindent
환
$$
(B_y)^\wedge = (B[1/y])^\wedge = \lim B[1/y]/(xy)^n B[1/y] =
\lim B[1/y] / x^n B[1/y]
$$
이 비축소임을 주장한다. 이 환은
$$
(A_y)^\wedge = (A[1/y])^\wedge = \lim A[1/y]/(xy)^n A[1/y] =
\lim A[1/y] / x^n A[1/y]
$$
위에서 $1, f, \ldots, f^{p - 1}$을 기저로 하는 자유 가군이다. 그러나 원소
$g \in (A_y)^\wedge$가 존재하여 $f^p = g^p$이다. 실제로
$g = \sum a_i x^i$를 택하면 된다. 이는 $f$에 쓴 것과 같은 식이며
$(A_y)^\wedge$에서 의미가 있다. 그러므로
$$
(B_y)^\wedge = (A_y)^\wedge[f]/(f^p - g^p) \cong
(A_y)^\wedge[\tau]/(\tau^p)
$$
는 비축소이다. 사실 이 예는 조금 더 말해 준다. $(A_y)^\wedge$는
균등화원 $x$를 갖고 잉여체가 위 이산 값매김환 $D$의 분수체인 이산
값매김환이다. 따라서 심지어
$$
(B_y)^\wedge[1/(xy)] = ((B_y)^\wedge)_{xy}
$$
도 비축소이다. 이는 다음 종류의 예를 준다.

\begin{lemma}
\label{examples-lemma-nonreduced-recompletion}
$2$차원 뇌터 국소 정역 $(B, \mathfrak m)$과 주 아이디얼 $I = (b)$가 존재한다.
이 정역은 그 아이디얼에 관해 완비이고, $f \in \mathfrak m$,
$f \not \in I$인 원소가 존재하여 $I$-진 완비화
$C = (B_f)^\wedge$, 곧 주 국소화 $B_f$의 이 완비화가 비축소이고, 심지어
$C_b = C[1/b] = (B_f)^\wedge[1/b]$도 비축소이다.
\end{lemma}

\begin{proof}
위의 논의를 보라.
\end{proof}








\section{사슬 조건을 만족하지 않는 뇌터 국소환}
\label{examples-section-non-catenary-Noetherian-local}

\noindent
뇌터 국소환에는 성공적인 차원론이 있지만 사슬 조건을 만족하지 않는 뇌터
국소환도 존재한다. 한 예는 \cite[Appendix, Example 2]{Nagata}에 있다.
여기서는 가장 단순한 경우의 이 예를 제시한다. 즉 뇌터 국소 정역 $A$를
구성하는데, 그 차원은 $2$이고 보편적으로 사슬 조건을 만족하지 않는다.
($A$ 자체는 자동으로 사슬 조건을 만족한다. 『연습문제』의 문제
\ref{exercises-exercise-Noetherian-local-domain-dim-2-catenary}를 보라.)
보편적으로 사슬 조건을 만족하지 않는 뇌터 국소환이 존재하면 사슬 조건을
만족하지 않는 뇌터 국소환도 존재한다. 이 절 끝에서 지금의 구체적 예로 이를
자세히 설명한다.

\medskip\noindent
$k$를 체라 하고 형식적 멱급수환 $k[[x]]$, 즉 $k$ 위의 한 변수 형식적
멱급수환을 생각하자.
$$
z = \sum\nolimits_{i = 1}^\infty a_i x^i
$$
를 형식적 멱급수라 하자. $z$가 로랑 급수체
$k((x)) = k[[x]][1/x]$의 원소로서 $k(x)$ 위에서 초월적이라고 가정한다.
다음과 같이 놓자.
$$
z_j
=
x^{-j}(z - \sum\nolimits_{i = 1, \ldots, j - 1} a_i x^i)
=
\sum\nolimits_{i = j}^\infty a_i x^{i - j}
\in k[[x]].
$$
$z = xz_1$임에 주의하라. $R$을 $k[[x]]$의 부분환으로서 $x$, $z$와 모든
$z_j$가 생성하는 환이라 하자. 즉
$$
R = k[x, z_1, z_2, z_3, \ldots ] \subset k[[x]].
$$
$\mathfrak m = (x)$과
$\mathfrak n = (x - 1, z_1, z_2 + a_1, z_3 + a_1 + a_2, \ldots)$를
$R$의 아이디얼로 생각하자.

\medskip\noindent
$xz_{j + 1} + a_j = z_j$이다. 따라서 $R/\mathfrak m = k$이고
$\mathfrak m$은 극대 아이디얼이다. 더욱이 $R$의 원소가
$\mathfrak m$에 속하지 않으면 $k[[x]]$에서 가역원으로 사상되므로
$R_{\mathfrak m} \subset k[[x]]$이다. 사실 $R_{\mathfrak m}$이 잉여체
$k$인 이산 값매김환임을 쉽게 이끌어낼 수 있다.

\medskip\noindent
다음을 주장한다.
$$
R/(x - 1) =
k[x, z_1, z_2, z_3, \ldots ]/(x - 1)
\cong
k[z].
$$
실제로 $z_{j + 1} = z_j - a_j - (x - 1)z_{j + 1}$이고, 따라서
$z_{j + 1}$의 류를 $z_j$의 류로 몫 $R/(x - 1)$에서 나타낼 수 있다. 구성상
$R$의 분수체는 초월차수가 $2$인 $k$의 확대체이므로 $z$는 $k$ 위에서
$R/(x - 1)$ 안에서 초월적이고, 원하는 동형사상이 따른다. 따라서
$\mathfrak n = (x - 1, z)$는 극대 아이디얼이다. 사실 사상
$$
k[x, x^{-1}, z]_{(x - 1, z)} \longrightarrow R_{\mathfrak n}
$$
은 동형사상이다($x^{-1}$은 $R_{\mathfrak n}$에서 가역이고
$z_{j + 1} = x^{-1}z_j - a_j = \ldots = f_j(x, x^{-1}, z)$이기 때문이다).
따라서 $R_{\mathfrak n}$은 차원 $2$이고 잉여체가 $k$인 정칙 국소환이다.

\medskip\noindent
$S$를 다음 곱셈적 부분집합이라 하자.
$$
S =
(R \setminus \mathfrak m) \cap (R \setminus \mathfrak n) =
R \setminus (\mathfrak m \cup \mathfrak n)
$$
그리고 $B = S^{-1}R$이라 놓자. 다음을 주장한다.
\begin{enumerate}
\item $B$는 $k$-대수이다.
\item $B$의 극대 아이디얼은 $\mathfrak mB$와 $\mathfrak nB$ 두 개이다.
\item 이 극대 아이디얼들에서의 잉여체는 $k$이다.
\item $B_{\mathfrak mB} = R_{\mathfrak m}$과
$B_{\mathfrak nB} = R_{\mathfrak n}$은 각각 차원 $1$, $2$인 뇌터 정칙 국소환이다.
\item $B$는 뇌터 환이다.
\end{enumerate}
확인의 세부 사항은 생략한다.

\medskip\noindent
위 성질들을 갖는 $k$-대수 $B$가 주어질 때마다 다음과 같이 예를 얻는다.
$A = k + \text{rad}(B) \subset B$라 놓되
$\text{rad}(B) = \mathfrak mB \cap \mathfrak nB$는 Jacobson 근기라 하자.
$B$는 $A$ 위에서 유한이고, 따라서 Eakin의 정리(\cite{Eakin},
\cite[Appendix A1]{Nagata}, 또는 나중에 넣을 참고문헌을 보라)에 의해
$A$는 뇌터 환임을 쉽게 알 수 있다. 또한 $A$는 $B$와 같은 분수체와 잉여체
$k$를 갖는 국소 정역이다. $B$의 차원이 $2$이므로 『대수학』의 보조정리
\ref{algebra-lemma-integral-sub-dim-equal}에 의해 $A$의 차원도 $2$이다.

\medskip\noindent
$A$가 보편적으로 사슬 조건을 만족한다면 『대수학』의 보조정리
\ref{algebra-lemma-dimension-formula}의 차원 공식이
$\dim(B_{\mathfrak mB}) = 2$를 주는데, 이는 모순이다.

\medskip\noindent
$B$는 $A$ 위에서 한 원소로 생성된다. 따라서
$B = A[x]/\mathfrak p$이고, 여기서 $\mathfrak p$는 $A[x]$의 소 아이디얼이다.
$\mathfrak m' \subset A[x]$를 $\mathfrak mB$에 대응하는 극대 아이디얼이라 하자.
한편으로 $\dim(A[x]_{\mathfrak m'}) = 3$이고, 다른 한편
$$
(0)
\subset \mathfrak pA[x]_{\mathfrak m'}
\subset \mathfrak m'A[x]_{\mathfrak m'}
$$
는 소 아이디얼의 극대 사슬이다. 따라서 $A[x]_{\mathfrak m'}$은 사슬 조건을
만족하지 않는 뇌터 국소환의 예이다.





\section{나쁜 뇌터 국소환의 존재}
\label{examples-section-bad}

\noindent
$(A, \mathfrak m, \kappa)$를 뇌터 완비 국소환이라 하자.
\cite{Lech}에서는 $A$가 뇌터 국소 정역의 완비화가 되는 충분조건으로
$\text{depth}(A) \geq 1$이고 $A$가 부분환으로
$\mathbf{Q}$ 또는 $\mathbf{F}_p$를 포함하거나, $\mathbf{Z}$를
부분환으로 포함하면서 $A$가 $\mathbf{Z}$-가군으로서 꼬임이 없는 경우를 보였다.
이는 ``기묘한'' 성질을 가진 뇌터 국소 정역의 많은 예를 만들어 준다.

\medskip\noindent
예를 들어 이를 $A = \mathbf{C}[[x, y]]/(y^2)$에 적용하면 완비화가
비축소인 뇌터 국소 정역을 얻는다.
절 \ref{examples-section-local-completion-nonreduced}와 비교하라.

\medskip\noindent
\cite{LLPY}에서는 $A$가 축소 국소 뇌터 환의 완비화가 될 조건을 특징지었다.

\medskip\noindent
\cite{Heitmann-completion-UFD}에서는 $A$가 국소 뇌터 유일 인수분해 정역
$R$의 완비화가 되는 충분조건으로 $\text{depth}(A) \geq 2$이고 $A$가
부분환으로 $\mathbf{Q}$ 또는 $\mathbf{F}_p$를 포함하거나, $\mathbf{Z}$를
부분환으로 포함하면서 $A$가 $\mathbf{Z}$-가군으로서 꼬임이 없는 경우를 보였다.
특히 $R$은 정규이고(『대수학』의 보조정리 \ref{algebra-lemma-UFD-normal}),
따라서 $R$의 헨젤화도 정규 정역이다
(『대수학 더 보기』의 보조정리 \ref{more-algebra-lemma-henselization-normal}).
그러므로 위와 같은 $A$는 헨젤 뇌터 국소 정규 정역의 완비화이다
($R$과 그 헨젤화의 완비화는 같다. 『대수학 더 보기』의 보조정리
\ref{more-algebra-lemma-henselization-noetherian}을 보라).

\medskip\noindent
이를 적용하면 뇌터 국소 유일 인수분해 정역 $R$으로서
$R^\wedge \cong \mathbf{C}[[x, y, z, w]]/(wx, wy)$인 것을 얻는다.
$\Spec(R^\wedge)$는 정칙 $2$차원 성분과 정칙 $3$차원 성분의 합집합이다.
환 $R$은 보편적으로 사슬 조건을 만족할 수 없다. 극대 아이디얼을 부풀리기한
$$
X \longrightarrow \Spec(R)
$$
을 생각하자. 그러면 $X$는 정수적 스킴이다.
닫힌점 $x \in X$로서 $\dim(\mathcal{O}_{X, x}) = 2$인 것이 존재한다.
구체적으로 완비 국소환의 수준에서 $x$를
$2$차원 성분의 순변환 위에 놓이고 $3$차원 성분의 순변환 위에는 놓이지 않게
택한다. 『사상』의 보조정리 \ref{morphisms-lemma-dimension-formula}에 의해
$R$은 보편적으로 사슬 조건을 만족하지 않는다.
절 \ref{examples-section-non-catenary-Noetherian-local}과 비교하라.

\medskip\noindent
위 환은 $3$차원 국소 뇌터 유일 인수분해 정역이므로 사슬 조건을 만족한다.
그러나 \cite{Ogoma-example}에서는 정규 국소 뇌터 정역 $R$으로서
$R^\wedge \cong \mathbf{C}[[x, y, z, w]]/(wx, wy)$이고
$R$이 사슬 조건을 만족하지 않는 것을 구성한다.
\cite{Heitmann-Ogoma}와 \cite{Lech-YAPO}도 보라.

\medskip\noindent
\cite{Heitmann-isolated}에서는 $A$가 고립 특이점을 갖는 국소 뇌터 환
$R$의 완비화가 되기 위한 조건으로, $A$가 부분환으로
$\mathbf{Q}$ 또는 $\mathbf{F}_p$를 포함하거나
$A$의 잉여 표수가 $p > 0$이고 $p$가 $A$의 어떤 진 국소화에서도
영이 아닌 영인자로 보내지지 않는 경우를 보였다.
여기서 뇌터 국소환 $R$이 고립 특이점을 갖는다는 것은
$R_\mathfrak p$가 정칙 국소환이고, 이것이 모든 비극대 소 아이디얼
$\mathfrak p \subset R$에 대해 성립한다는 뜻이다.

\medskip\noindent
논문 \cite{Nishimura-few}와 \cite{Nishimura-few-II}에는 주어진 완비화를
갖는 ``나쁜'' 뇌터 국소환들의 긴 목록이 있다. 특히 $2$차원 Nagata
국소 정규 정역으로서 완비화가 $\mathbf{C}[[x, y, z]]/(yz)$인 예와
완비화가 $\mathbf{C}[[x, y, z]]/(y^2 - z^3)$인 또 하나의 예를 구성한다.

\medskip\noindent
여담으로 \cite{Loepp}에서는 $A$가 축소이고 등차원이며
$A$의 어떤 정수도 영인자가 아니면 $A$가 우수 뇌터 국소 정역의
완비화임을 보였다. 그러나 이 경우에는 우수한 환을 얻으므로
``나쁜'' 뇌터 국소환을 주지 않는다!




\section{뇌터 Jacobson 환에서의 차원}
\label{examples-section-noetherian-jacobson}

\noindent
$k$를 유한체의 대수적 폐포라 하자.
$A = k[x, y]$이고 $X = \Spec(A)$라 하자.
$C = V(x)$를 $y$축이라 하자(이는 임의의 다른 $1$차원 정수적
닫힌 부분스킴이어도 되며, 그 부분스킴은 $X$ 안에 있다).
$C_1, C_2, C_3, \ldots$를 $X$의 다른 $1$차원 정수적 닫힌 부분스킴들의
열거라 하자. $p_1, p_2, p_3, \ldots$를
$C$의 닫힌점들의 열거라 하자.

\medskip\noindent
주장: 모든 $n$에 대해 기약 닫힌집합 $Z_n \subset X$가 존재하는데,
그 차원은 $1$이고 집합론적으로
$$
\{p_n\} = Z_n \cap (C \cup C_1 \cup \ldots \cup C_n)
$$
이다. 이를 위해 $Y = C \cup C_1 \cup C_2 \cup \ldots \cup C_n$이라 놓자.
이는 축소 아핀 $1$차원 대수적 스킴이며 $k$ 위에 놓여 있다.
어떤 $f \in k[x, y]$에 대해 집합론적으로
$V(f) \cap  Y = \{p_n\}$가 되게 찾으면 충분하다.
그러면 $Z_n$을 $V(f)$의 적당한 기약 성분으로 택할 수 있기 때문이다.
제한 사상
$$
k[x, y] \longrightarrow \Gamma(Y, \mathcal{O}_Y)
$$
이 전사이므로 정칙함수 $g$를 $Y$ 위에서 찾되, 그 영점집합이
집합론적으로 $\{p_n\}$가 되게 하면 충분하다.
이를 위해 유효 Cartier 인자 $D \subset Y$로서 지지가
$p_n$인 것을 택한다(『다양체』의 보조정리
\ref{varieties-lemma-complement-codim-1-closed-points}에 의해 가능하다).
따라서 $\mathcal{O}_X(ND) \cong \mathcal{O}_X$가 어떤 $N > 0$에 대해
성립함을 보이면 충분하다. 유한체의 대수적 폐포 위의 아핀 $1$차원 대수적 스킴의
Picard 군은 꼬임 군이므로(여기에 나중의 참고문헌을 넣을 것) 주장이 따른다.

\medskip\noindent
위와 같은 $Z_n$을 모든 $n$에 대해 택하자. $k[x, y]$는 유일 인수분해
정역이므로 $Z_n = V(f_n)$이라 쓸 수 있으며, 여기서 $f_n \in A$는 기약원이다.
$S \subset k[x, y]$를 $f_1, f_2, f_3, \ldots$가 생성하는 곱셈적
부분집합이라 하자. 뇌터 환 $B = S^{-1}A$를 생각하자.

\medskip\noindent
분명히 환 준동형 $A \to B$는 국소환들을 동일시하고
단사 $\Spec(B) \to \Spec(A)$를 유도한다.
더욱이 곡선 $C_1$을 보면 $C \cap C_1$의 점들만
$\Spec(A)$에서 $\Spec(B)$로 갈 때 제거된다.
특히 $\Spec(B)$에는 $A$의 극대 아이디얼에 대응하는 극대 아이디얼이
무한히 많이 있다. 한편 $xB$는 극대 아이디얼이다. 실제로
$B/xB$의 스펙트럼은 유일한 소 아이디얼로 이루어진다. 이는
$C = V(x)$의 모든 닫힌점을
제거했지만 일반점은 제거하지 않았기 때문이다.
끝으로 $i \geq 1$에 대해 곡선 $C_i$를 생각하자.
$C_i = V(g_i)$라 쓰자. 여기서 $g_i \in A$는 기약원이다.
$C_i = Z_n$인 어떤 $n$이 있으면 $g_iB$는 단위 아이디얼이다.
그렇지 않으면 $C_i$의 닫힌점 중 유한 개만 제외하고 모두
$A$에서 $B$로 갈 때 남는다. 구체적으로 제거되는 점은
$(Z_1 \cup \ldots \cup Z_{i - 1} \cup C) \cap C_i$의 점들이며,
이 점들만 $C_i$에서 제거된다.

\medskip\noindent
위에서 기술한 $B$의 소 스펙트럼 구조와
『대수학』의 보조정리 \ref{algebra-lemma-noetherian-dim-1-Jacobson}에 의해
$B$는 Jacobson 환이다. 극대 아이디얼은 $A$의 극대 아이디얼들 중
$\Spec(B)$에 속하는 것들(이들은 무한히 많다)과 극대 아이디얼 $xB$이다.
따라서 차원 $1$과 $2$인 국소환들이 있다.

\begin{lemma}
\label{examples-lemma-Noetherian-Jacobson}
보편적으로 사슬 조건을 만족하는 Jacobson 뇌터 정역 $B$와
극대 아이디얼 $\mathfrak m_1, \mathfrak m_2$가 존재하여
$\dim(B_{\mathfrak m_1}) = 1$이고 $\dim(B_{\mathfrak m_2}) = 2$이다.
\end{lemma}

\begin{proof}
$B$의 구성은 위에 주어졌다. 『대수학』의 보조정리
\ref{algebra-lemma-localization-catenary}와
『사상』의 보조정리 \ref{morphisms-lemma-ubiquity-uc}에 의해
$B$가 보편적으로 사슬 조건을 만족한다는 점만 지적한다.
\end{proof}





\section{밑공간은 뇌터이지만 뇌터 스킴은 아닌 경우}
\label{examples-section-noetherian-not-noetherian}

\noindent
밑 위상 공간이 뇌터인 스킴이 뇌터 스킴이 아닐 수 있음을 보이는 두 예를 든다.

\begin{example}
\label{examples-example-many-variables}
$k$를 체라 하고 $A = k[x_1, x_2, x_3, \dots] /
(x_1^2, x_2^2, x_3^2, \dots)$라 하자. $A$의 모든 소 아이디얼은
멱영원 $x_1, x_2, x_3, \dots$를 포함하므로
$\mathfrak p = (x_1, x_2, x_3, \dots)$가 $A$의 유일한 소 아이디얼이다.
따라서 $\operatorname{Spec} A$의 밑 위상 공간은 한 점이고, 특히 뇌터이다.
그러나 $\mathfrak p$는 명백히 유한 생성이 아니다.
\end{example}

\begin{example}
\label{examples-example-many-mononials}
$k$를 체라 하고 $A \subseteq k[x, y]$를
$k$와 단항식 $\{xy^i\}_{i \ge 0}$가 생성하는 부분환이라 하자.
$A$의 소 아이디얼 중 $x$를 포함하지 않는 것은
$A_x \cong k[x, x^{-1}, y]$의 소 아이디얼과 일대일로 대응한다.
소 아이디얼 $\mathfrak p$가 $x$를 포함한다면 모든 $xy^i$, $i \ge 0$도
포함한다. 실제로 $(xy^i)^2 = x(xy^{2i}) \in \mathfrak p$이고
$\mathfrak p$는 근기 아이디얼이기 때문이다.
따라서 $\mathfrak p = (\{xy^i\}_{i \ge 0})$이다.
그러므로 $\operatorname{Spec} A$의 밑 위상 공간은 뇌터이다. 이 공간은
뇌터 스킴 $\Spec(A[x, x^{-1}, y])$의 점들과 소 아이디얼 $\mathfrak p$로
이루어지기 때문이다. 그러나 환 $A$는 $\mathfrak p$가 유한 생성이 아니므로
비뇌터이다. 이 예에서 $A$는 정역이기도 하다.
\end{example}





\section{정규화는 준아핀이지만 자신은 준아핀이 아닌 다양체}
\label{examples-section-nonquasi-affine}

\noindent
이런 종류의 예가 존재한다는 언급은
\cite[II Remark 6.6.13]{EGA}에 있다. 거기서는 그러한 예로 EGA 제5권을
인용하지만 제5권은 출판되지 않았다.

\medskip\noindent
$k$를 체라 하자.
$Y = \mathbf{A}^2_k \setminus \{(0, 0)\}$라 하자.
유한 전사 쌍유리 사상
$\pi : Y \longrightarrow X$를 구성하되,
$X$는 $k$ 위의 다양체이고 $X$는 준아핀이 아니게 하겠다.
구체적으로 $Y$ 안의 다음 곡선들을 생각하자.
$$
\begin{matrix}
C_1 & : & x = 0 \\
C_2 & : & y = 0
\end{matrix}
$$
$C_1 \cap C_2 = \emptyset$임에 주의하라. 동형사상
$\varphi : C_1 \to C_2$, $(0, y) \mapsto (y^{-1}, 0)$을 택하자.
다음 조건을 만족하는 위와 같은 사상 $\pi : Y \to X$가 유일하게
존재한다고 주장한다.
$$
\xymatrix{
C_1
\ar@<1ex>[rr]^{\text{id}} \ar@<-1ex>[rr]_{\varphi}
& &
Y \ar[r]^\pi & X
}
$$
이 도식은 다양체의 범주에서, 심지어 스킴의 범주에서도 쌍대등화자 도식이다.
이를 잠시 받아들이고 그러한 $X$가 준아핀일 수 없음을 보이자.
실제로 분명히
$$
\Gamma(X, \mathcal{O}_X) =
\{ f \in k[x, y] \mid f(0, y) = f(y^{-1}, 0)\} =
k \oplus (xy) \subset k[x, y].
$$
를 얻는다. 특히 이 함수들은 점 $(1, 0)$과
$(-1, 0)$을 분리하지 못하지만, 그 두 점의 $X$에서의 상은 서로 다르다
(아래에서 이를 보일 것이다). 단, 이는 $k$의 표수가 $2$가 아닐 때의 말이다.

\medskip\noindent
$X$가 존재함을 보이기 위해 Zariski 열린집합
$D(x + y) \subset Y$를 $Y$에서 생각하자. 이는 환
$k[x, y, 1/(x + y)]$의 스펙트럼이고,
곡선 $C_1$, $C_2$는 모두 $D(x + y)$에 들어간다. 더욱이 사상
$$
C_1 \amalg C_2
\longrightarrow
D(x + y) \cap Y = \Spec(k[x, y, 1/(x + y)])
$$
은 닫힌 몰입이다. 『대수학 더 보기』의 보조정리
\ref{more-algebra-lemma-fibre-product-finite-type}에 의해 환
$$
A =
\{f \in k[x, y, 1/(x + y)] \mid f(0, y) = f(y^{-1}, 0)\}
$$
은 $k$ 위의 유한형이다. 한편 열린집합
$D(xy) \subset Y$가 $Y$에 있고, 이는 곡선 $C_1$과
$C_2$ 모두와 만나지 않는다. 이는 환
$$
B = k[x, y, 1/xy].
$$
의 스펙트럼이다. $A_{xy} \cong B_{x + y}$임에 주의하라
($A$는 원소 $xyP(x, y)$를 포함하는데 이는 임의의 다항식 $P$에 대해
성립하며, 또 원소 $xy/(x + y)$도 분명히 포함하기 때문이다).
스킴 $X$는 아핀 스킴 $\Spec(A)$와 $\Spec(B)$를 동형사상
$A_{xy} \cong B_{x + y}$로 붙여 얻으며, 따라서 분명히
$k$ 위의 유한형이다. 이것이 분리임을 보이려면 환 준동형
$A \otimes_k B \to B_{x + y}$가 전사임을 보이면 된다.
이를 위해 $A \otimes_k B$가 원소
$xy/(x + y) \otimes 1/xy$를 포함하고 이것이 $1/(x + y)$로 감을 이용하라.
사상 $Y \to X$는 자연스러운 사상
$D(x + y) \to \Spec(A)$와 $D(xy) \to \Spec(B)$로 주어진다.
이 둘은 모두 유한이므로 $Y \to X$도 원하는 대로 유한이다.
$X$가 실제로 위 도식의 쌍대등화자라는 확인은 생략한다. 다만
(여기에 스킴의 범주에서의 밂아웃에 관한 나중의 참고문헌을 넣을 것)을 보라.
사상 $\pi : Y \to X$가 점 $(1, 0)$과
$(-1, 0)$을 $X$의 서로 다른 점으로 보내는 이유는 함수
$(x + y^3)/(x + y)^2 \in A$의 첫째 값은
$1/1$이고, 둘째 값은\ $-1/(-1)^2 = -1$이며 이들은 언제나 다르기 때문이다
(표수가 $2$이면 독자가 표수 $2$에 맞는 점들을 직접 찾아 보라).
이 논의를 보조정리로 요약한다.

\begin{lemma}
\label{examples-lemma-quasi-affine-normalization-not-quasi-affine}
$k$를 체라 하자.
정규화는 준아핀이지만 자신은 준아핀이 아닌 다양체 $X$가 존재한다.
\end{lemma}

\begin{proof}
위의 논의와 (여기에 정규화에 관한 나중의 참고문헌을 넣을 것)을 보라.
\end{proof}



\section{스킴론적 상 취하기}
\label{examples-section-scheme-theoretic-image}

\noindent
$k$를 체라 하자. $t$를 변수라 하자. $Y = \Spec(k[t])$이고
$X = \coprod_{n \geq 1} \Spec(k[t]/(t^n))$라 하자.
사상 $f : X \to Y$는 닫힌 몰입
$\Spec(k[t]/(t^n)) \to \Spec(k[t])$을 사용해 정의하되,
이를 각 $n \geq 1$에 대해 취한다. 이 경우 다음이 성립한다.
\begin{enumerate}
\item $f$의 스킴론적 상
(『사상』의 정의 \ref{morphisms-definition-scheme-theoretic-image})은
$Y$이다. 한편 $f$의 상은 닫힌점
$t = 0$이며, 이 점은 $Y$에 있다. 따라서 $f$의 스킴론적 상의 밑 닫힌집합은
$f$의 상의 폐포와 같지 않다.
\item 스킴론적 상의 구성은 열린 부분스킴
$V = \Spec(k[t, 1/t]) \subset Y$로의 제한과 가환하지 않는다.
실제로 $V$의 $X$에서의 역상은 공집합이므로
$f|_{f^{-1}(V)} : f^{-1}(V) \to V$의 스킴론적 상은 공 스킴이다.
이는 $Y \cap V$와 같지 않다.
\end{enumerate}




\section{국소 닫힌 부분집합의 상}
\label{examples-section-non-chevalley}

\noindent
Chevalley의 정리는 아핀 스킴 사이의 유한 표시 사상에 의한
구성가능집합의 상이 구성가능하다고 말한다. 『대수학』의 정리
\ref{algebra-theorem-chevalley}와
『사상』의 절 \ref{morphisms-section-constructible}을 보라.
국소 닫힌 부분집합의 상에는 같은 명제가 성립하지 않음을 보이겠다.

\medskip\noindent
$k$를 표수가 $0$인 체라 하자. 사영 사상
$$
f :
X = \Spec(k[t, x_1, x_2, \ldots, y_1, y_2, \ldots])
\longrightarrow
\Spec(k[x_1, x_2, \ldots, y_1, y_2, \ldots]) = Y
$$
을 생각하자. 이는 유한 표시 사상이다.
닫힌 부분집합 $Z$를 $X$ 안에서
$$
x_1(t - 1) = 0,\quad
x_2(t - 1)(t - 2) = 0,\quad
x_3(t - 1)(t - 2)(t - 3) = 0,\quad \ldots
$$
로 정의하자. $U = \bigcup_{j \geq 1} U_j$를 다음과 같이 정의되는
$X$의 열린집합이라 하자.
$$
U_j = \text{다음 식 }y_j(t - 1)(t - 2) ... (t - j)\text{이 영이 아닌 점들}
$$
그러면
$$
f(Z \cap U_j) =
\text{다음 조건을 만족하는 점들: }x_1, \ldots, x_j\text{는 영이고 }y_j\text{는 영이 아님}
$$
이다. $B = f(Z \cap U) = \bigcup_{j \geq 1} f(Z \cap U_j)$는
$Y$의 국소 닫힌 부분집합들의 유한 합집합이 아니라고 주장한다.

\medskip\noindent
주장을 증명하자. $B = A_1 \cup \cdots \cup A_m$이
$B$의 유한 피복이고 그 구성원은 $Y$의 국소 닫힌 부분집합들이라고 가정하자.
$n$에 관한 귀납법으로 $m \geq n$임을 보이겠다.
$n = 1$인 경우는 $B$가 비어 있지 않으므로 성립한다.
$n > 1$이고 $n - 1$에 대한 귀납 가정이 성립한다고 하자.
$B$의 폐포는 $(x_1 = 0)$이므로 어떤 $A_i$는
$(x_1 = 0)$의 비어 있지 않은 열린 부분집합을 포함해야 한다.
그러면 $A_i$는 $(x_1 = 0)$에서 열려 있어야 한다.
그러나 그러한 열린 부분집합은 $y_1 = 0$인 점을 포함할 수 없다.
실제로 $B$의 점에서 $y_1 = 0$이면 $x_2 = 0$이고, 따라서
$B$는 $(x, y)$의 어떤 근방도 $(x_1 = 0)$ 안에서 포함하지 않는다.
그러므로 나머지 $m - 1$개의 집합을 제한하면
$B \cap (y_1 = 0)$의 구성가능 피복을 얻는다.
그런데 오른쪽 이동 사상
$x_i \mapsto x_{i + 1}$, $y_i \mapsto y_{i + 1}$은 $B$를
$B \cap (y_1 = 0)$과 동일시한다! 따라서 귀납 가정에 의해
$m - 1 \geq n - 1$이고, 따라서 $m \geq n$이다.
이로써 귀납 단계와 주장이 증명되었다.

\begin{lemma}
\label{examples-lemma-no-chevalley}
아핀 스킴 사이의 유한 표시 사상 $f : X \to Y$와
국소 닫힌 부분집합 $T$가 존재하며, 이는 $X$의 부분집합이고
$f(T)$는 $Y$의 국소 닫힌 부분집합들의 유한 합집합이 아니다.
\end{lemma}

\begin{proof}
위의 논의를 보라.
\end{proof}









\section{어떤 닫힌 부분스킴에서도 열리지 않는 국소 닫힌 부분스킴}
\label{examples-section-strange-immersion}

\noindent
이는 『사상』의 예 \ref{morphisms-example-thibaut}를 옮긴 것이다.
열린 몰입 다음에 닫힌 몰입을 합성한 꼴이 아닌 몰입의 예를 들자.
$k$를 체라 하자.
$X = \Spec(k[x_1, x_2, x_3, \ldots])$라 하자.
$U = \bigcup_{n = 1}^{\infty} D(x_n)$라 하자.
그러면 $U \to X$는 열린 몰입이다.
다음 아이디얼을 생각하자.
$$
I_n =
(x_1^n, x_2^n, \ldots, x_{n - 1}^n, x_n - 1, x_{n + 1}, x_{n + 2}, \ldots)
\subset
k[x_1, x_2, x_3, \ldots][1/x_n].
$$
임의의 $m \not = n$에 대해
$I_n k[x_1, x_2, x_3, \ldots][1/x_nx_m] = (1)$임에 주의하라.
따라서 준연접 아이디얼 $\widetilde I_n$은 $D(x_n)$ 위에 있고,
이들은 $D(x_nx_m)$에서 일치한다. 즉 $n \not = m$이면
$\widetilde I_n|_{D(x_nx_m)} = \mathcal{O}_{D(x_n x_m)}$이다.
따라서 이 아이디얼들을 붙이면 준연접 아이디얼층
$\mathcal{I} \subset \mathcal{O}_U$를 얻는다.
$Z \subset U$를 $\mathcal{I}$에 대응하는 닫힌 부분스킴이라 하자.
그러면 $Z \to X$는 몰입이다.

\medskip\noindent
$Z \to X$를 $Z \to \overline{Z} \to X$로 분해할 수 없다고 주장한다.
여기서 $\overline{Z} \to X$는 닫혀 있고
$Z \to \overline{Z}$는 열려 있다고 하자. 그렇다면 $\overline{Z}$는
아이디얼 $I \subset k[x_1, x_2, x_3, \ldots]$로 정의되어야 하고,
이 아이디얼은 $I_n = I k[x_1, x_2, x_3, \ldots][1/x_n]$을
만족해야 한다. 그러나 원소 $f \in k[x_1, x_2, x_3, \ldots]$ 가운데
모든 $I_n$에 들어가는 것은 $0$뿐이다! 따라서 그런 $I$는 없다.

\medskip\noindent
사상 $Z \to X$는 몰입의 스킴론적 상이 나쁘게 행동하는 예도 준다.
위 논증에 의해 몰입 $Z \to X$의 스킴론적 상은 $X$이다. 한편
\begin{enumerate}
\item $Z$는 $X$에서 위상적으로 조밀하지 않다.
\item $Z = Z \cap U \to U$의 스킴론적 상은 $Z$ 자체이다. 이는
$U \cap X = U$와 같지 않으므로, 이 경우 스킴론적 상의 구성은
열린집합으로의 제한과 가환하지 않는다.
\end{enumerate}







\section{적당한 열린집합의 부존재}
\label{examples-section-nonexistence-opens}

\noindent
이 절은 『성질』의 절 \ref{properties-section-finding-affine-opens}의 결과를
보충한다.

\medskip\noindent
$k$를 체라 하고 $A = k[z_1, z_2, z_3, \ldots]/I$라 하자. 여기서 $I$는
모든 쌍곱 $z_iz_j$, $i \not = j$, $i, j \in \mathbf{N}$가 생성하는
아이디얼이다. $S = \Spec(A)$라 놓자. $s \in S$를 극대 아이디얼
$(z_i)$에 대응하는 닫힌점이라 하자. $V \subset S \setminus \{s\}$인
준콤팩트 열린집합 가운데 $S \setminus \{s\}$에서 조밀한 것은 없다고
주장한다. $S \setminus \{s\} = \bigcup D(z_i)$임에 주의하라.
각 $D(z_i)$는 일반점 $\eta_i$를 갖는 열린 기약집합이다. 따라서
모든 $i$에 대해 $\eta_i \in V$이다. 그러나 $S \setminus \{s\}$의
주 아핀 열린집합은 $D(f)$ 꼴이며, 여기서 $f \in (z_1, z_2, \ldots)$이다.
어떤 $n$에 대해 $f \in (z_1, \ldots, z_n)$이므로 $D(f)$는 점
$\eta_i$를 유한 개만 포함한다. 따라서 $V$는 준콤팩트일 수 없다.

\medskip\noindent
$k$를 체라 하고 $B = k[x, z_1, z_2, z_3, \ldots]/J$라 하자. 여기서 $J$는
곱 $xz_i$, $i \in \mathbf{N}$와 모든 쌍곱
$z_iz_j$, $i \not = j$, $i, j \in \mathbf{N}$가 생성하는 아이디얼이다.
$T = \Spec(B)$라 놓자. 주 열린집합 $U = D(x)$를 생각하자.
$V \subset T$인 준콤팩트 열린집합 가운데 $V \cap U = \emptyset$이고
$V \cup U$가 $T$에서 조밀한 것은 없다고 주장한다. $t \in T$를
극대 아이디얼 $(x, z_i)$에 대응하는 닫힌점이라 하자. $T$에서 $U$의 폐포는
$\overline{U} = U \cup \{t\}$이다. 따라서
$V \subset \bigcup_i D(z_i)$는 준콤팩트 열린집합이다. 앞 문단의 논증에
의해 $V$는 $\bigcup D(z_i)$에서 조밀할 수 없다.

\begin{lemma}
\label{examples-lemma-complement-of-affine-does-not-contain-qc-dense-open}
아핀 안의 준콤팩트 열린집합의 부존재:
\begin{enumerate}
\item 아핀 스킴 $S$와 아핀 열린집합 $U \subset S$가 존재하여,
$V \subset S$인 준콤팩트 열린집합 가운데 $U \cap V = \emptyset$이고
$U \cup V$가 $S$에서 조밀한 것은 없다.
\item 아핀 스킴 $S$와 닫힌점 $s \in S$가 존재하여
$S \setminus \{s\}$는 준콤팩트 조밀 열린집합을 포함하지 않는다.
\end{enumerate}
\end{lemma}

\begin{proof}
위의 논의를 보라.
\end{proof}

\noindent
위의 아핀 스킴 $T$ 두 복사본을 아핀 열린집합 $U$를 따라 붙인 것을 $X$라
하자. 그러면 사상 $\pi : X \to T$와 $X = U_1 \cup U_2$가 있어 $\pi$는
$U_i$를 $T$와 동형으로, $U_1 \cap U_2$를 $U$와 동형으로 보낸다.
$X$는 준분리이고(『스킴』의 보조정리
\ref{schemes-lemma-characterize-quasi-separated}) 준콤팩트이다.
$W \subset X$인 분리이고 조밀한 준콤팩트 열린집합은 존재하지 않는다고
주장한다. $x_1 \in U_1$, $x_2 \in U_2$를 위에서 도입한 닫힌점
$t \in T$로 가는 두 닫힌점이라 하자. $\tilde \eta \in U_1 \cap U_2$를
$U$의 유일한 일반점 $\eta$로 가는 일반점이라 하자. 특수화
$\eta \leadsto t$ 위에서 $\tilde\eta \leadsto x_1$과
$\tilde\eta \leadsto x_2$가 성립함에 주의하라.
$\pi|_W : W \to T$가 분리이므로 $x_1$이 그 열린집합에 속하면서 동시에
$x_2 \in W$일 수는 없다(분리성의 값 판정법인 『스킴』의 보조정리
\ref{schemes-lemma-valuative-criterion-separatedness}).
$x_1 \not \in W$라 하자. 그러면 $W \cap U_1$은 ($X$가 준분리이므로)
$U_1$의 준콤팩트 조밀 열린집합이고 $x_1$을 포함하지 않는다.
이제 스킴의 동형사상 $(T, t) \cong (S, s)$가 존재함에 주의하라.
이는 $x$를 $z_1$로, $z_i$를 $z_{i + 1}$로 보내어 얻는다.
따라서 이 절의 첫 문단에 의해 모순이다.

\begin{lemma}
\label{examples-lemma-no-dense-separated-quasi-compact-open-in-qcqs}
분리이고 준콤팩트인 조밀 열린집합을 포함하지 않는 준콤팩트 준분리 스킴
$X$가 존재한다.
\end{lemma}

\begin{proof}
위의 논의를 보라.
\end{proof}






\section{준콤팩트 조밀 열린 부분스킴의 부존재}
\label{examples-section-nonexistence-qc-dense-open-subscheme}

\noindent
$X$를 체 $k$ 위의 준콤팩트 준분리 대수공간이라 하자. 스킴 자취
$X' \subset X$는 조밀한 열린 부분공간이다. 『공간의 성질』의 명제
\ref{spaces-properties-proposition-locally-quasi-separated-open-dense-scheme}를 보라.
사실 이 결과는 $X$가 합리적일 때에도 성립한다. 『양호한 공간』의 명제
\ref{decent-spaces-proposition-reasonable-open-dense-scheme}를 보라.
$X$의 준콤팩트 조밀 열린 부분스킴을 찾을 수 있는지가 자연스러운 질문인데,
일반적으로는 불가능하다.

\medskip\noindent
$k$의 표수가 2가 아니라고 하자.
$B = k[x, z_1, z_2, z_3, \ldots]/J$라 하자. 여기서 $J$는 곱
$xz_i$, $i \in \mathbf{N}$와 모든 쌍곱
$z_iz_j$, $i \not = j$, $i, j \in \mathbf{N}$가 생성하는 아이디얼이다.
$U = \Spec(B)$라 놓자. 모든 좌표가 영인 닫힌점을 $0 \in U$라 쓰자.
다음과 같이 놓는다.
$$
j : R = \Delta \amalg \Gamma \longrightarrow U \times_k U
$$
여기서 $\Delta$는 $k$ 위의 $U$의 대각사상의 상이고
$$
\Gamma = \{((x, 0, 0, 0, \ldots), (-x, 0, 0, 0, \ldots))
\mid x \in \mathbf{A}^1_k, x \not = 0\}.
$$
$s, t : R \to U$가 \'etale임은 분명하므로
$j$는 \'etale 동치관계이다. 몫 $X = U/R$은 대수공간이다
(『공간』의 정리 \ref{spaces-theorem-presentation}).
$j$는 몰입이 아니다. 실제로 $(0, 0) \in \Delta$는 $\Gamma$의 폐포에
속한다. 따라서 $X$는 스킴이 아니다. 한편 $R$이 준콤팩트이므로
$X$는 준분리이다. 점 $0 \in U$의 상을 $0_X$라 쓰자.
$X \setminus \{0_X\}$는 스킴이라고 주장한다. 실제로
$$
X \setminus \{0_X\} =
\Spec\left(k[x^2, x^{-2}]\right) \amalg
\Spec\left(k[z_1, z_2, z_3, \ldots]/(z_iz_j)\right) \setminus \{0\}
$$
이다(세부 사항은 생략한다). 한편 절 \ref{examples-section-nonexistence-opens}에서
오른쪽 스킴이 준콤팩트 조밀 열린집합을 포함하지 않음을 보았다.

\begin{lemma}
\label{examples-lemma-nonexistence-qc-dense-open-subscheme}
준콤팩트 조밀 열린 부분스킴을 포함하지 않는 준콤팩트 준분리 대수공간이
존재한다.
\end{lemma}

\begin{proof}
위의 논의를 보라.
\end{proof}

\noindent
『공간』의 예 \ref{spaces-example-non-representable-descent}의 구성을,
위에서 『공간』의 예 \ref{spaces-example-affine-line-involution}의 구성을
사용한 것과 같은 방식으로 사용하면,
준콤팩트 조밀 열린 부분스킴을 포함하지 않는 준콤팩트하고 준분리이며
국소 분리인 대수공간의 예를 얻는다.


\section{대수공간 위의 아핀 스킴}
\label{examples-section-embedding-affines}

\medskip\noindent
$f : Y \to X$가 스킴의 사상이고 $f$가 국소 유한형이며 $Y$가 아핀이라고
하자. 그러면 $Y$에서 $X$ 위의 아핀 $n$-공간으로 가는 몰입
$Y \to \mathbf{A}^n_X$이 존재한다. 조금 더 일반적인 『사상』의 보조정리
\ref{morphisms-lemma-quasi-affine-finite-type-over-S}를 보라.

\medskip\noindent
이제 $f : Y \to X$가 대수공간의 사상이고 $f$가 국소 유한형이며
$Y$가 아핀 스킴이라고 하자. 이 경우에는 일반적으로 $Y$에서
$X$ 위의 아핀 $n$-공간으로 가는 몰입을 찾을 수 없다.

\medskip\noindent
첫 번째 (고약한) 반례는 $Y = \Spec(k)$,
$X = [\mathbf{A}^1_k/\mathbf{Z}]$이다. 여기서 $k$는 표수 영인 체이고
$\mathbf{Z}$는 $\mathbf{A}^1_k$에 평행이동 $(n, t) \mapsto t + n$으로
작용한다. 실제로 $X$ 위의 임의의 사상 $Y \to \mathbf{A}^n_X$을
$X$의 피복 $\mathbf{A}^1_k$로 당기면
$\mathbf{A}^1_k$들의 무한 서로소 합집합에서 $\mathbf{A}^{n + 1}_k$로 가는
사상을 얻는데, 이는 몰입이 아니다.

\medskip\noindent
두 번째 반례는 $Y = \mathbf{A}^1_k \to X = \mathbf{A}^1_k/R$이고
$R = \{(t, t)\} \amalg \{(t, -t), t \not = 0\}$인 경우이다. 실제로 이 경우
사상 $Y \to \mathbf{A}^n_X$은 $Y$ 위의 정칙함수
$f_1, \ldots, f_n$으로 주어진다. 따라서 $Y$와 피복
$\mathbf{A}^{n + 1}_k \to \mathbf{A}^n_X$의 섬유곱은 스킴
$$
\{(f_1(t), \ldots, f_n(t), t)\} \amalg
\{(f_1(t), \ldots, f_n(t), -t), t \not = 0\}
$$
이다. 이것에서 $\mathbf{A}^{n + 1}_k$로 가는 명백한 사상은 몰입이 아니다.
이는 $X$가 준분리인 반례를 준다는 점에 주의하라.

\begin{lemma}
\label{examples-lemma-cannot-embed-into-affine}
대수공간의 유한형 사상 $Y \to X$ 가운데 $Y$가 아핀이고 $X$가 준분리이며,
$X$ 위의 몰입 $Y \to \mathbf{A}^n_X$이 존재하지 않는 것이 있다.
\end{lemma}

\begin{proof}
위의 논의를 보라.
\end{proof}









\section{준연접 가군의 밂앞}
\label{examples-section-push-quasi-coherent}

\noindent
『스킴』의 보조정리 \ref{schemes-lemma-push-forward-quasi-coherent}에서
$f_*$는 $f$가 준콤팩트 준분리일 때 준연접 가군을 준연접 가군으로 보냄을
증명했다. 두 조건이 모두 필요함을 보이는 예들을 들자.

\medskip\noindent
$Y = \Spec(A)$를 아핀 스킴이라 하고
$X = \coprod_{n \in \mathbf{N}} Y$라 하자. 명백한 사상을
$f : X \to Y$라 할 때 $f_*\mathcal{O}_X$는 준연접이 아니라고 주장한다.
실제로 $a \in A$에 대해
$$
f_*\mathcal{O}_X(D(a)) = \prod\nolimits_{n \in \mathbf{N}} A_a
$$
이다. 따라서 $f_*\mathcal{O}_X$가 준연접이려면 모든 $a \in A$에 대해
$$
\prod\nolimits_{n \in \mathbf{N}} A_a
=
\left(\prod\nolimits_{n \in \mathbf{N}} A\right)_a
$$
이어야 한다. 일반적으로 이는 거짓이다. 예를 들어
$A = \mathbf{Z}$이고 $a = 2$이면 $(1, 1/2, 1/4, 1/8, \ldots)$는 왼쪽에는
속하지만 오른쪽에는 속하지 않는다. $f$는 준콤팩트가 아닌 분리 사상이다.

\medskip\noindent
$k$를 체라 하자. 다음과 같이 놓는다.
$$
A = k[t, z, x_1, x_2, x_3, \ldots]/(tx_1z, t^2x_2^2z, t^3x_3^3z, \ldots)
$$
$Y = \Spec(A)$라 하자. $V \subset Y$를 열린 부분스킴
$V = D(x_1) \cup D(x_2) \cup \ldots$라 하자. $Y$의 두 복사본을 $V$를
따라 붙인 것을 $X$라 하고, $f : X \to Y$를 명백한 사상이라 하자.
그러면 완전열
$$
0 \to f_*\mathcal{O}_X \to
\mathcal{O}_Y \oplus \mathcal{O}_Y \xrightarrow{(1, -1)} j_*\mathcal{O}_V
$$
이 있다. 여기서 $j : V \to Y$는 포함 사상이다. 사상
$$
A \longrightarrow \prod A_{x_n}
$$
이 단사이므로(세부 사항은 생략한다)
$\Gamma(Y, f_*\mathcal{O}_X) = A$이다. 한편 사상
$$
A_t \longrightarrow \prod A_{tx_n}
$$
의 핵은 원소 $z$를 포함하므로 영이 아니다. 따라서
$\Gamma(D(t), f_*\mathcal{O}_X)$는 $A_t$보다 엄밀히 크다. 실제로
$(z, 0)$을 포함한다. 그러므로 $f_*\mathcal{O}_X$는 준연접이 아니다.
$f$는 준콤팩트이지만 준분리가 아닌 사상이다.

\begin{lemma}
\label{examples-lemma-pushforward-quasi-coherent}
『스킴』의 보조정리 \ref{schemes-lemma-push-forward-quasi-coherent}는
준콤팩트성 가정과 준분리성 가정 중 어느 것도 뺄 수 없다는 의미에서
최선이다.
\end{lemma}

\begin{proof}
위의 논의를 보라.
\end{proof}










\section{계수 1의 유한 자유 줄기를 갖는 유한이 아닌 가군}
\label{examples-section-nonfree}

\noindent
$R = \mathbf{Q}[x]$라 하자. $R$의 분수체의 부분가군으로
$M = \sum_{n \in \mathbf{N}} \frac{1}{x - n}R$이라 놓자. 그러면 $M$은
유한 생성이 아니지만, $R$의 모든 소 아이디얼 $\mathfrak p$에 대해
$R_{\mathfrak p}$-가군으로서
$M_{\mathfrak p} \cong R_{\mathfrak p}$이다.

\medskip\noindent
비슷한 예로 $R = \mathbf{Z}$와
$M = \sum_{p \text{ prime}} \frac{1}{p} \mathbf{Z} \subset \mathbf{Q}$를
들 수 있다. 이는 $b$가 영이 아니고 제곱 인수가 없을 때의 분수
$\frac{a}{b}$들의 집합과 같다.










\section{각 줄기에서는 가역이지만 가역이 아닌 아이디얼}
\label{examples-section-locally-invertible-not-invertible}

\noindent
$A$를 정역, $I \subset A$를 영이 아닌 아이디얼이라 하자. $I$가
가역이라는 말은 $I$가 $A$-가군으로서 가역이라는 뜻임을 상기하라.
$I$ 자체는 가역이 아니지만 모든 소 아이디얼 $\mathfrak q \subset A$에 대해
$I_\mathfrak q = (f) \subset A_\mathfrak q$가 영이 아닌 주 아이디얼인 예를
만들겠다. 문헌에서는 모든 소 아이디얼 $\mathfrak q$에 대해
$I_\mathfrak q$가 주 아이디얼이라는 성질을 때때로 ``$I$가 국소 주
아이디얼이다''라고 표현한다. 그러나 우리에게 ``국소''는 항상 ``Zariski
위상에서 국소''(또는 현재 다루는 위상에서 국소)라는 뜻이므로 이 용어를
쓸 수 없다.

\medskip\noindent
$R = \mathbf{Q}[x]$라 하고, 절 \ref{examples-section-nonfree}에서 구성한 가군
$M = \sum \frac{1}{x - n}R$을 취하자. 환\footnote{환 $A$는 국소환들이
뇌터 환인 비뇌터 정역의 예이다.}
$$
A = \text{Sym}^*_R(M)
$$
과 아이디얼 $I = M A = \bigoplus_{d \geq 1} \text{Sym}^d_R(M)$을
생각하자. $M$은 $R$-가군으로서 유한 생성이 아니므로 $I$는 $A$의
아이디얼로서 유한 개의 원소로 생성될 수 없다. 가역 가군은 유한 생성이므로
$I$는 가역이 아니다. 한편 소 아이디얼 $\mathfrak p \subset R$을 잡자.
구성에 의해 $M_\mathfrak p \cong R_\mathfrak p$이다. 따라서
$$
A_\mathfrak p = \text{Sym}^*_{R_\mathfrak p}(M_\mathfrak p) \cong
\text{Sym}^*_{R_\mathfrak p}(R_\mathfrak p) =
R_\mathfrak p[T]
$$
이며 이는 등급 $R_\mathfrak p$-대수의 동형이다. 따라서
$I_\mathfrak p \subset A_\mathfrak p$는 영인자가 아닌 $T$로 생성된다.
그러므로 임의의 소 아이디얼 $\mathfrak q \subset A$에 대해서도
$I_\mathfrak q$는 한 원소로 생성된다.

\begin{lemma}
\label{examples-lemma-locally-principal-not-invertible}
정역 $A$와 영이 아닌 아이디얼 $I \subset A$가 존재하여, 모든 소 아이디얼
$\mathfrak q \subset A$에 대해 $I_\mathfrak q \subset A_\mathfrak q$는
주 아이디얼이지만 $I$는 가역 $A$-가군이 아니다.
\end{lemma}

\begin{proof}
위의 논의를 보라.
\end{proof}










\section{사영 가군이 아닌 유한 평탄 가군}
\label{examples-section-finite-flat-not-projective}

\noindent
이는 『대수학』의 주의 \ref{algebra-remark-warning}를 옮긴 것이다.
유한 $R$-가군이 $R$-평탄하다고 해서 자동으로 사영 가군인 것은 아니다.
반례는 $R = \mathcal{C}^\infty(\mathbf{R})$가 $\mathbf{R}$ 위의 무한 번
미분 가능한 함수들의 환이고 $M = R_{\mathfrak m} = R/I$인 경우이다.
여기서 $\mathfrak m = \{f \in R \mid f(0) = 0\}$이고
$I = \{f \in R \mid \exists \epsilon, \epsilon > 0 :
f(x) = 0\ \forall x, |x| < \epsilon\}$이다.

\medskip\noindent
사상 $\Spec(R/I) \to \Spec(R)$은 열린 사상이 아닌 평탄 닫힌 몰입의
예이기도 하다.

\begin{lemma}
\label{examples-lemma-finite-flat-non-projective}
기묘한 평탄 가군.
\begin{enumerate}
\item 환 $R$과, 사영 가군이 아닌 유한 평탄 $R$-가군 $M$이 존재한다.
\item 평탄하지만 열린 사상은 아닌 닫힌 몰입이 존재한다.
\end{enumerate}
\end{lemma}

\begin{proof}
위의 논의를 보라.
\end{proof}



\section{국소 자유가 아닌 사영 가군}
\label{examples-section-projective-not-locally-free}

\noindent
두 예를 들겠다. 하나는 계수가 $0$과 $1$ 사이이고, 다른 하나는 계수가
$\aleph_0$인 예이다.

\begin{lemma}
\label{examples-lemma-ideal-generated-by-idempotents-projective}
$R$을 환이라 하자. 아이디얼 $I \subset R$이 가산 개의 멱등원으로 생성되면
$I$는 사영 $R$-가군이다.
\end{lemma}

\begin{proof}
$I = (e_1, e_2, e_3, \ldots)$이고 각 $e_n$이 $R$의 멱등원이라고 하자.
귀납적으로 $e_{n + 1}$을 $e_n + (1 - e_n)e_{n + 1}$로 바꾸면
$(e_1) \subset (e_2) \subset (e_3) \subset \ldots$라 가정할 수 있다. 따라서
$I = \bigcup_{n \geq 1} (e_n) = \colim_n e_nR$이다. 이 경우
$$
\Hom_R(I, M) = \Hom_R(\colim_n e_nR, M)
= \lim_n \Hom_R(e_nR, M) = \lim_n e_nM
$$
이다. 전이 사상 $e_{n + 1}M \to e_nM$은 $e_n$을 곱하여 주어지고 전사이다.
따라서 『대수학』의 보조정리 \ref{algebra-lemma-ML-exact-sequence}에 의해
함자 $\Hom_R(I, M)$은 완전하다. 즉 $I$는 사영 $R$-가군이다.
\end{proof}

\noindent
$P \subset Q$를 $R$-가군의 포함이라 하고 $Q$는 유한 $R$-가군,
$P$는 국소 자유라고 하자. 『대수학』의 정의
\ref{algebra-definition-locally-free}를 보라. $Q$가 $N$개의 원소로 생성되는
$R$-가군이라고 하자. 그러면 『대수학』의 보조정리
\ref{algebra-lemma-map-cannot-be-injective}에 의해 $P$는 유한 국소 자유이다
(자유인 부분들의 계수는 $N$ 이하이다). 이 경우 『대수학』의 보조정리
\ref{algebra-lemma-finite-projective}에 의해 $P$는 유한 $R$-가군이다.

\medskip\noindent
이를 위 결과와 결합하면 가산 개의 멱등원으로 생성되는 유한 생성이 아닌
아이디얼은 사영이지만 국소 자유가 아니다. 구체적으로
$R = \prod_{n \in \mathbf{N}} \mathbf{F}_2$로 놓고, $I$를 멱등원
$$
e_n = (1, 1, \ldots, 1, 0, \ldots )
$$
들이 생성하는 아이디얼로 잡을 수 있다. 여기서 $1$의 열은 길이 $n$이다.

\begin{lemma}
\label{examples-lemma-ideal-projective-not-locally-free}
어떤 환 $R$과 아이디얼 $I$가 존재하여, $I$는 사영 $R$-가군이지만
국소 자유 $R$-가군은 아니다.
\end{lemma}

\begin{proof}
위를 보라.
\end{proof}

\begin{remark}
\label{examples-remark-projective-with-dual-finite-not-finite}
위의 $R = \prod_{n \in \mathbf{N}} \mathbf{F}_2$와
$e_n = (1, 1, \ldots, 1, 0, \ldots )$이 생성하는 아이디얼 $I$의 예는 쌍대가
유한 생성이고 심지어 유한 자유이지만 자신은 유한 생성이 아닌 사영 가군
$M = I$의 예도 준다. 실제로 보조정리
\ref{examples-lemma-ideal-generated-by-idempotents-projective}의 증명 전략을 써서
독자는 $\Hom_R(I, R) \cong R$임을 보일 수 있다.
\end{remark}

\begin{lemma}
\label{examples-lemma-chow-group-product}
$K$를 체라 하자. $C_i$ ($i = 1, \ldots, n$)를 $K$ 위의 매끄럽고 사영이며
기하학적으로 기약인 곡선이라 하자. $P_i \in C_i(K)$를 유리점이라 하고
$Q_i \in C_i$를 $[\kappa(Q_i) : K] = 2$인 점이라 하자. 그러면
$U_i = C_i \setminus \{Q_i\}$일 때
$[P_1 \times \ldots \times P_n]$은
$\CH_0(U_1 \times_K \ldots \times_K U_n)$에서 영이 아니다.
\end{lemma}

\begin{proof}
차수 사상
$\deg : \CH_0(C_1 \times_K \ldots \times_K C_n) \to \mathbf{Z}$가 있다.
각 $Q_i$는 $K$ 위에서 차수 $2$이므로 ``경계''
$$
C_1 \times_K \ldots \times_K C_n
\setminus
U_1 \times_K \ldots \times_K U_n
$$
에 지지된 모든 영차원 순환의 차수는 $2$로 나누어진다.
\end{proof}

\noindent
위 보조정리로 사영이지만 국소 자유가 아닌 또 다른 가군을 다음과 같이
구성할 수 있다. $C_n$ ($n = 1, 2, 3, \ldots$)을 $\mathbf{Q}$ 위의 매끄럽고
사영이며 기하학적으로 기약인 곡선으로 잡고, 각각에 점
$P_n, Q_n \in C_n$이 있어 $\kappa(P_n) = \mathbf{Q}$이고
$\kappa(Q_n)$은 $\mathbf{Q}$의 이차 확대라고 하자.
$U_n = C_n \setminus \{Q_n\}$이라 놓으면 이는 아핀 곡선이다.
$\mathcal{L}_n$을 $U_n$ 위에서 $P_n$의 아이디얼층의 역이라 하자.
영차원 순환군 $\CH_0(U_n)$에서
$c_1(\mathcal{L}_n) = [P_n]$임에 주의하라.
$A_n = \Gamma(U_n, \mathcal{O}_{U_n})$라 놓자.
$L_n = \Gamma(U_n, \mathcal{L}_n)$은 계수 $1$인 국소 자유 $A_n$-가군이다.
다음과 같이 놓는다.
$$
B_n = A_1 \otimes_{\mathbf{Q}} A_2 \otimes_{\mathbf{Q}} \ldots
\otimes_{\mathbf{Q}} A_n
$$
$\Spec(B_n) = U_1 \times \ldots \times U_n$이며, 모든 곱은
$\Spec(\mathbf{Q})$ 위에서 취한다. $i \leq n$에 대해
$$
L_{n, i} =
A_1 \otimes_{\mathbf{Q}} \ldots \otimes_{\mathbf{Q}} L_i
\otimes_{\mathbf{Q}} \ldots \otimes_{\mathbf{Q}} A_n
$$
이라 놓자. 이는 계수 $1$인 국소 자유 $B_n$-가군이며
$\text{pr}_i^*\mathcal{L}_n$의 대역 단면이기도 하다. 다음과 같이 놓는다.
$$
B_\infty = \colim_n B_n
\quad\text{and}\quad
L_{\infty, i} = \colim_n L_{n, i}
$$
끝으로 다음과 같이 놓는다.
$$
M = \bigoplus\nolimits_{i \geq 1} L_{\infty, i}.
$$
이는 유한 국소 자유 가군들의 직합이므로 사영이다. $M$은 국소
자유가 아니라고 주장한다. 그렇지 않으면 어떤 영이 아닌 함수
$f \in B_\infty$에 대해 $M_f$가 $(B_\infty)_f$ 위에서 자유라고 하자.
그 기저를 $e_1, e_2, \ldots$라 하자. $n \geq 1$을 택하여
$f \in B_n$이라 하자. $m \geq n + 1$을 택하여
$e_1, \ldots, e_{n + 1}$이 다음 직합에 속하게 하자.
$$
\bigoplus\nolimits_{1 \leq i \leq m} L_{m, i}.
$$
충실히 평탄한 밑변환 뒤에
$e_1, \ldots, e_{n + 1}$이 기저의 일부이므로 Chern 류
$$
c_i(\text{pr}_1^*\mathcal{L}_1 \oplus \ldots \oplus
\text{pr}_m^*\mathcal{L}_m), \quad i = m, m - 1, \ldots, m - n
$$
는
$$
D(f) \subset U_1 \times \ldots \times U_m
$$
의 Chow 군에서 영이다. $f$는 $U_1 \times \ldots \times U_n$ 위 함수의
당김이다. 따라서 체 $K = \mathbf{Q}(C_1 \times \ldots \times C_n)$ 위의 다양체
$$
W = (C_{n + 1} \times \ldots \times C_m)_K
$$
에서는 특히 다음 소멸이 성립한다.
$$
c_{m - n}(\mathcal{O}_W^{\oplus n} \oplus
\text{pr}_1^*\mathcal{L}_{n + 1} \oplus \ldots
\oplus \text{pr}_{m - n}^*\mathcal{L}_m) = 0.
$$
달리 말해 순환
$$
[(P_{n + 1} \times \ldots \times P_m)_K]
$$
는 $W$ 위의 영차원 순환의 Chow 군에서 영이다. 이는 위 보조정리
\ref{examples-lemma-chow-group-product}와 모순이다. 실제로 $n + 1 \leq i \leq m$인
$Q_i$는 $(C_n)_K$ 위의 대응점 $Q_i'$를 유도하고, $K/\mathbf{Q}$가
기하학적으로 기약이므로 $[\kappa(Q_i') : K] = 2$이다.

\begin{lemma}
\label{examples-lemma-projective-not-locally-free}
가산 환 $R$과, 계수 $1$인 국소 자유 가군들을 가산 개 직합한 사영 가군
$M$이 존재하여 $M$은 국소 자유가 아니다.
\end{lemma}

\begin{proof}
위를 보라.
\end{proof}

\section{영이 아닌 평탄 아이디얼을 갖는 영차원 국소환}
\label{examples-section-zero-dimensional-flat-ideal}

\noindent
\cite{Lazard}와 \cite{Autour}에는 영이 아닌 평탄 아이디얼을 갖는 영차원
국소환의 예가 있다. 그 구성은 다음과 같다. $k$를 체라 하자.
$X_i, Y_i$ ($i \geq 1$)를 변수라 하자.
$R = k[X_i, Y_i]/(X_i - Y_i X_{i + 1}, Y_i^2)$로 놓자. 이 환에서
$x_i$와\ $y_i$는 각각 $X_i$와\ $Y_i$의 상이다. 이 환에서
$$
x_i = y_i x_{i + 1} = y_i y_{i +1} x_{i + 2} =
y_i y_{i + 1} y_{i + 2} x_{i + 3} = \ldots
$$
이다. 환 $R$의 소 아이디얼은 $\mathfrak m = (x_i, y_i)$ 하나뿐이다.
아이디얼 $I = (x_i)$가 평탄 $R$-가군임을 주장한다.

\medskip\noindent
$R$에서 $x_i$의 소멸자는 아이디얼
$(x_1, x_2, x_3, \ldots, y_i, y_{i + 1}, y_{i + 2}, \ldots)$이다.
$R$-가군 $M$을 원소 $e_i$ ($i \geq 1$)와 관계
$e_i = y_i e_{i + 1}$로 생성하자. $M$은 $R$의 복사본들의 직접극한
$\colim_i R$이므로 평탄이다. 그 전이 사상은
$$
R \xrightarrow{y_1} R \xrightarrow{y_2} R \xrightarrow{y_3} \ldots
$$
이다. $M$에서 $e_i$의 소멸자도 아이디얼
$(x_1, x_2, x_3, \ldots, y_i, y_{i + 1}, y_{i + 2}, \ldots)$이다.
$M$과\ $I$의 모든 원소는 각각 어떤 $f, h \in R$에 대해
$f e_i$와\ $h x_i$로 쓸 수 있다. 따라서 사상 $M \to I$,
$e_i \to x_i$는 동형사상이고 $I$는 평탄이다.

\begin{lemma}
\label{examples-lemma-zero-dimensional-flat-ideal}
\begin{slogan}
평탄 아이디얼을 갖는 영차원 환.
\end{slogan}
유일한 소 아이디얼을 갖는 국소환 $R$과, 평탄 $R$-가군인 영이 아닌
아이디얼 $I \subset R$이 존재한다.
\end{lemma}

\begin{proof}
위의 논의를 보라.
\end{proof}








\section{전사 함수가 아닌 영차원 환의 범주론적 전사}
\label{examples-section-epimorphism-not-surjective}

\noindent
\cite{Lazard-deux}와 \cite{Autour}에서 다음 예를 찾을 수 있다.
$k$를 체라 하고 환 준동형
$$
k[x_1, x_2, \ldots, z_1, z_2, \ldots]/
(x_i^{4^i}, z_i^{4^i})
\longrightarrow
k[x_1, x_2, \ldots, y_1, y_2, \ldots]/
(x_i^{4^i}, y_i - x_{i + 1}y_{i + 1}^2)
$$
을 생각하자. 이는 $x_i$를 $x_i$로, $z_i$를 $x_iy_i$로 보낸다.
오른쪽에서 $y_i^{4^{i + 1}}$은 영이지만 $y_1$은 영이 아니다
(세부 사항은 생략한다). 이 사상은 전사 함수가 아니다.
$\deg(x_i) = -1$, $\deg(z_i) = 0$, $\deg(y_i) = 1$로 두어 위 사상을
$\mathbf{Z}$-등급 대수의 사상으로 보면 $y_1$이 상에 속하지 않음이
분명하다. 끝으로
$$
y_{i - 1} \otimes 1 = x_i y_i^2 \otimes 1 = y_i \otimes x_i y_i =
x_i y_i \otimes y_i = 1 \otimes x_i y_i^2.
$$
따라서 이 사상은 범주론적 전사이다. 즉 표적을 원천 위에서 자기 자신과
텐서한 것은 표적과 동형이다.

\begin{lemma}
\label{examples-lemma-epi-not-surjective}
차원 $0$인 국소환 사이의 범주론적 전사 가운데 전사 함수가 아닌 것이
존재한다.
\end{lemma}

\begin{proof}
위의 논의를 보라.
\end{proof}



\section{유한형이고 유한 표시는 아니며 한 소 아이디얼에서 평탄인 경우}
\label{examples-section-ft-not-fp-flat-at-prime}

\noindent
$k$를 체라 하고 국소환 $A_0 = k[x, y]_{(x, y)}$를 생각하자.
$\mathfrak p_{0, n} = (y + x^n + x^{2n + 1})$라 쓰면 이는 소 아이디얼이다.
$$
A = A_0[z_1, z_2, z_3, \ldots]/(z_n z_m, z_n(y + x^n + x^{2n + 1}))
$$
라 놓자. $A \to A_0$는 핵의 제곱이 영인 전사이다. 따라서 $A$도
국소환이고 $A$의 소 아이디얼들은 $A_0$의 소 아이디얼들과 일대일로
대응한다. $\mathfrak p_n$을 $A$에서 $\mathfrak p_{0, n}$에 대응하는 소
아이디얼이라 하자. $\mathfrak p_n$은 $z_n$의 $A$에서의 소멸자이다. 다음과 같이 놓는다.
$$
C = A[z]/(xz^2 + z + y)[\frac{1}{2zx + 1}]
$$
$A \to C$는 \'etale 환 준동형이다.
『대수학』의 예 \ref{algebra-example-make-standard-smooth}를 보라.
$\mathfrak q \subset C$를 $x$, $y$, $z$와 모든 $z_n$이 생성하는 극대 아이디얼이라
하자. $A \to C$가 평탄이므로 $z_n$의 $C$에서의 소멸자는
$\mathfrak p_nC$이다. 다음을 계산한다.
\begin{align*}
C/\mathfrak p_n C
& =
A_0[z]/(xz^2 + z + y, y + x^n + x^{2n + 1})[1/(2zx + 1)] \\
& =
k[x]_{(x)}[z]/(xz^2 + z - x^n - x^{2n + 1})[1/(2zx + 1)] \\
& =
k[x]_{(x)}[z]/(z - x^n) \times
k[x]_{(x)}[z]/(xz + x^{n + 1} + 1)[1/(2zx + 1)] \\
& =
k[x]_{(x)} \times k(x)
\end{align*}
이는 $(z - x^n)(xz + x^{n + 1} + 1) = xz^2 + z - x^n - x^{2n + 1}$이기 때문이다.
따라서 $\mathfrak p_nC = \mathfrak r_n \cap \mathfrak q_n$이다. 여기서
$\mathfrak r_n = \mathfrak p_nC + (z - x^n)C$,
$\mathfrak q_n = \mathfrak p_nC + (xz + x^{n + 1} + 1)C$이다.
$\mathfrak q_n + \mathfrak r_n = C$이므로
$\mathfrak p_nC = \mathfrak r_n \mathfrak q_n$도 얻는다.
따라서 $\mathfrak q_n$은 $\xi_n = (z - x^n)z_n$의 소멸자이다.
한편 $\mathfrak r_n \subset \mathfrak q$이고, 다른 한편
$\mathfrak q_n + \mathfrak q = C$이다. 예를 들어 $\mathfrak q_n$이
$C$의 극대 아이디얼이고 $\mathfrak q$와 다르기 때문이다. 마찬가지로
$\mathfrak q_n + \mathfrak q_m = C$가 $n \not = m$이면 성립한다. 이제
$$
B = \Im(C \longrightarrow C_{\mathfrak q})
$$
라 놓자. $\xi_n$들은 $B$에서 영으로 간다. 이는 $xz + x^{n + 1} + 1$이
$\mathfrak q$에 속하지 않기 때문이다. $\mathfrak q' \subset B$를 $\mathfrak q$의 상이라 하자.
구성상 $B$는 유한형 $A$-대수이고
$B_{\mathfrak q'} \cong C_{\mathfrak q}$이다. 특히
$B_{\mathfrak q'}$는 $A$ 위에서 평탄이다.

\medskip\noindent
$g' \in B$, $g' \not \in \mathfrak q'$이고 $B_{g'}$가 $A$ 위의 유한 표시가
되는 원소는 없다고 주장한다. 증명의 개요를 제시한다.
$g \in C$ 중에서 $g' \in B$로 가는 원소를 택하고 $C_g \to B_{g'}$를 생각하자. 『대수학』의 보조정리
\ref{algebra-lemma-finite-presentation-independent}에 의해 $B_g$가 $A$ 위의
유한 표시일 필요충분조건은 $C_g \to B_{g'}$의 핵이 유한 생성인 것이다.
그러나 원소 $g \in C$가 $\mathfrak q$에 속하지 않으므로
$A_0[z]/(xz^2 + z + y)$에서 영이 아닌 원소로 간다. 따라서 $g$는 유한 개의
$\mathfrak q_n$에만 속할 수 있다. 실제로
$(y + x^n + x^{2n + 1}, xz + x^{n + 1} + 1)$들은 2차원 기약 뇌터 공간
$\Spec(k[x, y, z]/(xz^2 + z + y))$의 서로 다른 무한 개의 여차원 1인 점이다.
사상
$$
\bigoplus\nolimits_{g \not \in \mathfrak q_n} C/\mathfrak q_n
\longrightarrow
C_g, \quad
(c_n) \longrightarrow \sum c_n \xi_n
$$
은 단사이고 그 상은 $C_g \to B_{g'}$의 핵이다. 이 명제의 증명은 생략한다.
(힌트: $A = A_0 \oplus I$라 쓰고 이를 $A_0$-가군으로 보자. 여기서
$I$는 $A \to A_0$의 핵이다. 마찬가지로 $C = C_0 \oplus IC$라 쓰고,
$IC = \bigoplus Cz_n \cong \bigoplus (C/\mathfrak r_n \oplus C/\mathfrak q_n)$로
쓴 뒤 $g$를 곱하는 것이 각 직합 성분에 미치는 영향을 조사하라.)
이로써 주장의 증명 개요가 끝난다. 또한 $B_{g'}$가 $A$ 위에서 평탄하지 않음이
위와 같은 모든 $g'$에 대해 증명된다. 평탄이라면 $z_n$의
$B_{g'}$에서의 상의 소멸자는 $\mathfrak p_nB_{g'}$여서 $z - x^n$을 포함하지 않을
것이기 때문이다.

\medskip\noindent
그 결과 \cite[Part I, Remarques (3.4.7) (\romannumeral 5)]{GruRay}에서
제기된 질문에 부정적으로 답할 수 있다. 정확한 명제는 다음과 같다.

\begin{lemma}
\label{examples-lemma-example-raynaud-gruson}
국소환 $A$, 유한형 환 준동형 $A \to B$, 소 아이디얼 $\mathfrak q$가 존재하여
$\mathfrak m_A$ 위에 놓이고 $B_{\mathfrak q}$는 $A$ 위에서 평탄이지만,
$g \in B$, $g \not \in \mathfrak q$인 모든 원소에 대해 $B_g$는 $A$ 위의
유한 표시도 아니고 $A$ 위에서 평탄하지도 않다.
\end{lemma}

\begin{proof}
위의 논의를 보라.
\end{proof}

\section{유한형이고 평탄이지만 유한 표시가 아닌 경우}
\label{examples-section-finite-type-flat-not-finite-presentation}

\noindent
이 절에서는 유한형이고 평탄하지만 유한 표시가 아닌 환 준동형과 사상의
몇 가지 예를 든다.

\medskip\noindent
$R$을 택하되, 아이디얼 $I$가 있어서 $R/I$가 유한 평탄이지만 사영이 아닌
가군이 되게 하자. 구체적인 예는 절
\ref{examples-section-finite-flat-not-projective}을 보라. 이는 $I$가 유한 생성이
아님을 뜻한다. 『대수학』의 보조정리
\ref{algebra-lemma-finitely-generated-pure-ideal}을 보라. 또한
$I = I^2$이다. 『대수학』의 보조정리 \ref{algebra-lemma-pure}를 보라.
우리 예들의 밑환은 $R$, 밑스킴은 $S = \Spec(R)$이다.

\medskip\noindent
환 준동형 $R \to R \oplus R/I\epsilon$을 생각하자. 여기서 관례적으로
$\epsilon^2 = 0$이라 한다. 이는 유한이고 평탄하지만 유한
표시는 아니다. 모든 섬유환은 완전 교차이고 기하학적으로 기약이다.

\medskip\noindent
$A = R[x, y]/(xy, ay; a \in I)$라 하자. $R$-가군으로서
$A = \bigoplus_{i \geq 0} Ry^i \oplus \bigoplus_{j > 0} R/Ix^j$이다.
따라서 $R \to A$는 평탄 유한형이지만 유한 표시가 아닌 환 준동형이다.
각 섬유환은 $\kappa(\mathfrak p)[x, y]/(xy)$ 또는
$\kappa(\mathfrak p)[x]$와 동형이다.

\medskip\noindent
$B = R[X_0, X_1, X_2]/(X_1X_2, aX_2; a \in I)$로 놓아 앞 예를 사영 사상으로
바꿀 수 있다. 이 경우 $X = \text{Proj}(B) \to S$는 고유 평탄이지만 유한
표시가 아닌 사상이고, 각 $s \in S$에 대해 섬유 $X_s$는
$\mathbf{P}^1_s$ 또는 $\mathbf{P}^2_s$에서 $X_1X_2$의 영점으로 정의되는
닫힌 부분스킴과 동형이다(후자는 산술 종수 $0$인 사영
마디 곡선이다).

\medskip\noindent
$M = R \oplus R \oplus R/I$, $B = \text{Sym}_R(M)$를 $M$ 위의 대칭 대수라 하고
$X = \text{Proj}(B)$라 놓자. 그러면 $X \to S$는 고유 평탄이지만 유한 표시가 아닌 사상이고,
$s \in S$에 대한 기하 섬유는 $\mathbf{P}^1_s$ 또는 $\mathbf{P}^2_s$와
동형이다. 특히 이 섬유들은 매끄럽고 기하학적으로 기약이다.

\begin{lemma}
\label{examples-lemma-finite-type-flat-not-finitely-presented}
다음 예들이 존재한다.
\begin{enumerate}
\item 유한 표시가 아니면서 기하학적으로 기약인 완전 교차 섬유환을 갖는
평탄 유한형 환 준동형.
\item 유한 표시가 아니면서 기하학적으로 연결이고 기하학적으로 축소이며 차원 1인
완전 교차 섬유환을 갖는 평탄 유한형 환 준동형.
\item 유한 표시가 아닌 스킴의 고유 평탄 사상 $X \to S$로서, 모든 섬유가
$\mathbf{P}^1_s$ 또는 $X_1X_2$의 $\mathbf{P}^2_s$ 안의 영점 자취와 동형인 것.
\item 유한 표시가 아닌 스킴의 고유 평탄 사상 $X \to S$로서, 모든 섬유가
$\mathbf{P}^1_s$ 또는 $\mathbf{P}^2_s$와 동형인 것.
\end{enumerate}
\end{lemma}

\begin{proof}
위의 논의를 보라.
\end{proof}

\section{유한형 환 준동형의 위상}
\label{examples-section-topology-finite-type}

\noindent
$A \to B$를 국소 정역들의 국소 준동형이라 하자. $A$가 뇌터이고
$A \to B$가 본질적으로 유한형이며
$A/\mathfrak m_A \subset B/\mathfrak m_B$가 유한이면,
소 아이디얼 $\mathfrak q \subset B$가 존재하여 $\mathfrak q \not = \mathfrak m_B$이고
$A \to B/\mathfrak q$는 준유한 환 준동형의 국소화이다.
『사상 더 보기』의 보조정리
\ref{more-morphisms-lemma-quasi-finite-quasi-section-meeting-nearby-open}를 보라.

\medskip\noindent
이 절에서는 $A$가 뇌터가 아닐 때 이 결과가 거짓임을 보이는 예를 든다.
$k$를 체라 하고 다음 두 환을 생각하자.
$$
A = \{a_0 + a_1 x + a_2 x^2 + \ldots \mid a_0 \in k, a_i \in k((y))
\text{ for }i\geq 1\}
$$
및
$$
C = \{a_0 + a_1 x + a_2 x^2 + \ldots \mid a_0 \in k[y], a_i \in k((y))
\text{ for }i\geq 1\}.
$$
포함 $A \to C$는 유한형이다. 이는 $C$가 $y$에 의해 $A$ 위에서
생성되기 때문이다.
$A$는 극대 아이디얼
$\mathfrak m = \{a_1 x + a_2 x^2 + \ldots \in A\}$을 갖는 국소환이고, 소
아이디얼은 $(0)$과 $\mathfrak m$뿐이라고 주장한다. 실제로
원소 $f = a_0 + a_1 x + a_2 x^2 + \ldots$가 $A$에 속할 때
$a_0 \not = 0$이기만 하면 가역이다. $\mathfrak q \subset A$가 영이 아닌 소 아이디얼이고
$f = a_i x^i + \ldots \in \mathfrak q$라 하자. 멱급수의 성질을 쓰면 모든
$g \in k((y))$에 대해 $g^{i + 1} x^{i + 1} \in \mathfrak q$, 즉
$gx \in \mathfrak q$임을 알 수 있다. 따라서 $\mathfrak q = \mathfrak m$이다.

\medskip\noindent
$C$의 스펙트럼에 대해서도 위와 같이 논증하면, 모든 영이 아닌 소
아이디얼은 소 아이디얼
$\mathfrak p = \{a_1 x + a_2 x^2 + \ldots \in C\}$를 포함하며, 이는
$\mathfrak m$ 위에 놓인다. 따라서 $C$에서 $(0)$ 위에 놓인 소 아이디얼은
$(0)$뿐이다. $\mathfrak m_C = yC + \mathfrak p$, $B = C_{\mathfrak m_C}$라 놓자.
그러면 $A \to B$가 원하는 예이다.

\begin{lemma}
\label{examples-lemma-topology-finite-type}
국소 정역들의 국소 준동형 $A \to B$가 존재하여, 본질적으로 유한형이고
$A/\mathfrak m_A \to B/\mathfrak m_B$가 유한이며, 소 아이디얼
$\mathfrak q \not = \mathfrak m_B$를 $B$에서 어떻게 택하든 $A \to B/\mathfrak q$는 어떤
준유한 환 준동형의 국소화도 아니다.
\end{lemma}

\begin{proof}
위의 논의를 보라.
\end{proof}

\section{순수하지만 보편적으로 순수하지 않은 경우}
\label{examples-section-pure-not-universally}

\noindent
$k$를 체라 하자. 다음과 같이 놓는다.
$$
R = k[[x, xy, xy^2, \ldots]] \subset k[[x, y]].
$$
달리 말해 멱급수 $f \in k[[x, y]]$가 $R$에 속할 필요충분조건은
$f(0, y)$가 상수인 것이다. 특히 $R[1/x] = k[[x, y]][1/x]$이고 $R/xR$은
극대 아이디얼의 제곱이 영인 국소환이다. $R[y] \subset k[[x, y]]$를
멱급수 $f \in k[[x, y]]$ 중에서 $f(0, y)$가 $y$의 다항식인 것들의 집합이라 하자.
$R \to R[y]$는 유한형이지만 유한 표시가 아닌 환 준동형이고, $x$를 가역화하면
동형사상을 유도한다. 또한 제곱 영인 핵을 갖는 전사
$R[y]/xR[y] \to k[y]$가 있다. 유한 표시 환 준동형
$R \to S = R[t]/(xt - xy)$를 생각하자. 역시 $R[1/x] \to S[1/x]$는
동형사상이고, 이 경우 $S/xS \cong (R/xR)[t]/(xy)$는 멱영인 핵을 갖고
$k[t]$ 위로 전사한다. 전사 $S \to R[y]$, $t \longmapsto y$가 있으며, 이는
$x$를 가역화하면 동형사상을, $x$로 나눈 뒤에는 멱영인 핵을 갖는 전사를 유도한다.
따라서 $S \to R[y]$의 핵은 국소 멱영이다. 특히 $S \to R[y]$는 스펙트럼
위에서 보편 위상동형을 유도한다.

\medskip\noindent
먼저 $S$가 $S$-가군으로서 $R$ 위에서 상대적으로 순수함을 주장한다.
$x$를 가역화하면 동형사상을 얻으므로 극대 아이디얼
$\mathfrak m \subset R$에서만 확인하면 된다. $R$은 극대 아이디얼에 관해 완비이므로
헨젤 환이다. 따라서 모든 소 아이디얼 $\mathfrak p \subset R$,
$\mathfrak p \not = \mathfrak m$에 대해,
$\mathfrak q$를 $S$에서 $\mathfrak p$ 위에 놓인 유일한 소 아이디얼이라 할 때
$\mathfrak mS + \mathfrak q \not = S$를 만족함만 확인하면 된다.
$\mathfrak p \not = \mathfrak m$이므로 이는 $k[[x, y]][1/x]$의 유일한 소
아이디얼에 대응한다. 따라서 $\mathfrak p = (0)$이거나, $\mathfrak p = (f)$이다. 후자의 경우
$f \in k[[x, y]]$는 $x$와 동반이 아닌 기약원이다(여기서
$k[[x, y]]$가 유일 인수분해 정역임을 썼다. 나중의 참고문헌을 넣을 것).
첫 경우에는 $\mathfrak q = (0)$이어서 결과가 분명하다. 둘째 경우에는
$f$에 가역원을 곱하여 $f \in R[y]$라 할 수 있다(Weierstrass 준비정리;
세부 사항은 생략한다). 그러면 $R[y]/fR[y] \cong k[[x, y]]/(f)$임을 쉽게
알 수 있어 $f$는 $R[y]$의 소 아이디얼을 정의하고
$\mathfrak mR[y] + fR[y] \not = R[y]$이다. $S \to R[y]$가 스펙트럼 위에서
보편 위상동형을 유도하므로 $S$에 대해서도 원하는 결과를 얻는다.

\medskip\noindent
둘째로 $S$는 $R$ 위에서 보편적으로 상대 순수하지 않다고 주장한다.
이를 보이려면 값매김환 $\mathcal{O}$와 국소 환 준동형 $R \to \mathcal{O}$를
찾아 $\Spec(R[y] \otimes_R \mathcal{O}) \to \Spec(\mathcal{O})$가
$\Spec(\mathcal{O})$의 닫힌점을 만나지 않게 하면 충분하다. 동치로,
$\varphi : R \to \mathcal{O}$를 찾아 $\varphi(x) \not = 0$이고
$v(\varphi(x)) > v(\varphi(xy))$이게 하면 된다. 여기서 $v$는 $\mathcal{O}$의
값매김이다(이는 $y$의 값매김이 음수라는 뜻이다). 이를 위해 명백한 환 준동형
$$
R
\longrightarrow
\{a_0 + a_1 x + a_2 x^2 + \ldots \mid a_0 \in k[y^{-1}], a_i \in k((y))\}
$$
을 생각하자. 값매김환 $\mathcal{O}$로서, 오른쪽을 극대 아이디얼
$(y^{-1}, x)$에서 국소화한 환을 지배하는 것을 찾으면 원하는 결과를 얻는다.

\begin{lemma}
\label{examples-lemma-pure-not-universally-pure}
아핀 스킴의 유한 표시 사상 $X \to S$와 유한 표시
$\mathcal{O}_X$-가군 $\mathcal{F}$가 존재하여, $\mathcal{F}$는 $S$에 대해
상대 순수이지만 $S$에 대해 보편적으로 상대 순수하지 않다.
\end{lemma}

\begin{proof}
위의 논의를 보라.
\end{proof}

\section{평탄하지 않은 형식적 매끄러운 환 준동형}
\label{examples-section-formally-smooth-nonflat}

\noindent
$k$를 체라 하자. $k$-대수 $k[\mathbf{Q}]$를 생각하자. 이는
$k$-대수로서 기저가 $x_\alpha, \alpha \in \mathbf{Q}$이고 곱셈이
$x_\alpha x_\beta = x_{\alpha + \beta}$로 정해진다.
(특히 $x_0 = 1$이다.) 다음 $k$-대수 준동형을 생각하자.
$$
k[\mathbf{Q}] \longrightarrow k, \quad
x_\alpha \longmapsto 1.
$$
이는 핵이 $J$인 전사이고, 이 핵은 원소 $x_\alpha - 1$들로 생성된다.
$J/J^2$를 계산하자. $x_\alpha$를 곱하는 사상은 $J/J^2$ 위에서 항등이다.
$z_\alpha$로 $x_\alpha - 1$의 $J^2$에 대한 합동류를 나타내자. 이 류들이
$J/J^2$를 생성한다. 식
$$
(x_\alpha - 1)(x_\beta - 1) = x_{\alpha + \beta} - x_\alpha - x_\beta + 1 =
(x_{\alpha + \beta} - 1) - (x_\alpha - 1) - (x_\beta - 1)
$$
에서 $z_{\alpha + \beta} = z_\alpha + z_\beta$임을 $J/J^2$ 안에서 얻는다.
$J/J^2$의 일반적인 원소는 $\sum \lambda_\alpha z_\alpha$ 꼴이며, 여기서
$\lambda_\alpha \in k$ 중 영이 아닌 것은 유한 개뿐이다. $k$의 표수가
$p > 0$이면
$$
0 = pz_{\alpha/p} = z_{\alpha/p} + \ldots + z_{\alpha/p} = z_\alpha
$$
이므로 $J/J^2 = 0$이다. $k$의 표수가 영이면
$$
J/J^2 = \mathbf{Q} \otimes_{\mathbf{Z}} k \cong k
$$
이고(세부 사항은 생략한다) 영이 아니다.

\medskip\noindent
$k[\mathbf{Q}] \to k$는 $k$의 표수가 양수이면 형식적으로 매끄러운 환
준동형이라고 주장한다. 실선으로 주어진 다음 가환 도식을 생각하자.
$$
\xymatrix{
k \ar[r] \ar@{..>}[rd] & A \\
k[\mathbf{Q}] \ar[u] \ar[r]^\varphi & A' \ar[u]
}
$$
여기서 $A' \to A$는 핵 $I$가 제곱 영인 전사라고 하자.
$k[\mathbf{Q}] \to k$가 형식적으로 매끄럽다는 것을 보이려면 $\varphi$가 $k$를
거쳐 분해됨을 증명해야 한다. $\varphi(x_\alpha - 1)$은 $A$에서 영으로
가므로 $\varphi$는 사상 $\overline{\varphi} : J/J^2 \to I$를 유도하고,
이것의 소멸이 원하는 분해의 장애이다. $J/J^2 = 0$이 표수가 $p > 0$일 때
성립하므로 원하는 결과, 즉 이 경우 $k[\mathbf{Q}] \to k$가 형식적으로
매끄럽다는 결론을 얻는다.
끝으로 이 환 준동형은 평탄하지 않다. 예를 들어 영인자가 아닌
$x_2 - 1$을 영으로 보내기 때문이다.

\begin{lemma}
\label{examples-lemma-formally-smooth-nonflat}
평탄하지 않은 형식적 매끄러운 환 준동형이 존재한다.
\end{lemma}

\begin{proof}
위의 논의를 보라.
\end{proof}

\section{평탄하지 않은 형식적 \'etale 환 준동형}
\label{examples-section-formally-etale-nonflat}

\noindent
이 절에서는 \cite[0, Example 19.10.3(i)]{EGA}의 마지막 문장에 대한
반례를 든다(이는 나중의 정오표 목록에서도 잡히지 않았다).
$A \to A/J$를 생각하자. 여기서 $A$는 국소환이고 $J$는
$J^2 = J$인 영이 아닌 진 아이디얼이다(따라서 $J$는 유한 생성이 아니다).
대수적으로 닫힌 비아르키메데스 체의 값매김환에서 극대 아이디얼을 $J$로
두면 이러한 $(A, J)$의
예를 준다. 이 평탄하지 않은 몫사상은 형식적 \'etale이다.
실제로 다음 가환 도식이 주어졌다고 하자.
$$
\xymatrix{
A/J \ar[r] & R/I \\
A \ar[u] \ar[r]^\varphi & R \ar[u]
}
$$
여기서 $I$는 환 $R$의 아이디얼이고 $I^2 = 0$이다. 그러면 $A \to R$은
$A/J$를 거쳐 유일하게 분해된다. 실제로
$$
\varphi(J) = \varphi(J^2) \subset (\varphi(J)A)^2 \subset I^2 = 0.
$$
따라서 이는 형식적 \'etale 경우, 곧 국소 유한형이면서 형식적 \'etale인
사상에 대한 ``구조 정리''에도 반례를 준다. 이 정리는
\cite[IV, Theorem 18.4.6(i)]{EGA}에 있다
(실제로 사람들이 쓰는 (ii)에 대한 반례는 아니다). 후자 증명의 오류는
증명의 맨 마지막 단계에서 잘못된 \cite[0, Example 19.3.10(i)]{EGA}를
원용하는 데 있으며, 바로 그 때문에 위 반례가 들어온다.

\begin{lemma}
\label{examples-lemma-formally-etale-not-presented}
평탄하지 않은 형식적 \'etale 환 준동형이 존재한다.
\end{lemma}

\begin{proof}
위의 논의를 보라.
\end{proof}

\section{여접복합체가 자명하지 않은 형식적 \'etale 환 준동형}
\label{examples-section-formally-etale-nontrivial-cotangent-complex}

\noindent
$k$를 체라 하고 환
$$
R = k[\{x_n\}_{n \geq 1}, \{y_n\}_{n \geq 1}]/(
x_1y_1, x_{nm}^m - x_n, y_{nm}^m - y_n)
$$
을 생각하자. $A$를 모든 $x_n, y_n$이 생성하는 극대 아이디얼에서
국소화한 것이라 하고, 그 극대 아이디얼을 $J \subset A$라 쓰자.
$B = A/J$라 놓자. 구성상 $J^2 = J$이므로 $A \to B$는 형식적 \'etale이다
(절 \ref{examples-section-formally-etale-nonflat}을 보라). 원소
$x_1 \otimes y_1$은
$$
J \otimes_A J \longrightarrow J.
$$
의 핵에 속하는 영이 아닌 원소라고 주장한다. 실제로 $(A, J)$는 국소화
쌍 $(A_n, J_n)$들의 쌍대극한이다. 이 쌍들은 환
$$
R_n = k[x_n, y_n]/(x_n^n y_n^n)
$$
을 해당 극대 아이디얼에서 국소화하여 얻는다.
$x_1 \otimes y_1$은 원소
$x_n^n \otimes y_n^n \in J_n \otimes_{A_n} J_n$에 대응하고, 직접 계산하면
이는 영이 아니다(계산은 생략한다). $\otimes$가 쌍대극한과 가환하므로
결론이 따른다. \cite[III Section 3.3]{cotangent}에 의해 $J$는 약정칙이
아니다. 따라서 \cite[III Proposition 3.3.3]{cotangent}에 의해
여접복합체 $L_{B/A}$는 영이 아니다. 더 정확히 말할 수 있다.
$H_0(L_{B/A}) = \Omega_{B/A}$이고 $H_1(L_{B/A}) = 0$이다. 이는
$J/J^2 = 0$이기 때문이다.
그러나 Quillen의 기본 스펙트럼 열의 다섯 항 완전열
(『여접』의 주의 \ref{cotangent-remark-elucidate-ss} 또는
\cite[Corollary 8.2.6]{Reinhard})과
$\text{Tor}_2^A(B, B) = \Ker(J \otimes_A J \to J)$가 영이 아니라는 사실에서
$H_2(L_{B/A})$가 영이 아님을 얻는다.

\begin{lemma}
\label{examples-lemma-formally-etale-nontrivial-cotangent-complex}
형식적 \'etale 전사 환 준동형 $A \to B$가 존재하며
$L_{B/A}$는 영이 아니다.
\end{lemma}

\begin{proof}
위의 논의를 보라.
\end{proof}




\section{평탄하고 형식적 비분기라고 해서 형식적 \'etale인 것은 아니다}
\label{examples-section-flat-formally-unramified-not-formally-etale}

\noindent
『사상 더 보기』의 보조정리
\ref{more-morphisms-lemma-unramified-flat-formally-etale}에서는 스킴의 비분기
평탄 사상 $X \to S$가 형식적 \'etale임을 보인다.
이 절의 목적은 `비분기'를 `형식적 비분기'로 바꿀 수 없음을 보여 주는 두
예를 드는 것이다. 첫째 예는 완전체 환의 특수한 성질을 이용하고, 둘째 예는
뇌터 정칙환의 사상에 대해서도 결과가 실패함을 보인다.

\begin{lemma}
\label{examples-lemma-perfect-closure-polynomial-ring}
$A = \mathbb{F}_p[T]$를 $\mathbb{F}_p$ 위의 한 변수 다항식환이라 하고
$A_{perf}$를 $A$의 완전 폐포라 하자. 그러면 $A \rightarrow A_{perf}$는 평탄하고
형식적 비분기이지만 형식적 \'etale은 아니다.
\end{lemma}

\begin{proof}
Frobenius 사상 $F_A : A \to A$ 아래에서 표적 쪽 $A$는
$\{1, T, \dots, T^{p - 1}\}$을 기저로 하는 원천 위의 자유 가군이다.
따라서 $F_A$는 충실히 평탄이고, 충실히 평탄한 사상들의 쌍대극한인
$A \to A_{perf}$도 충실히 평탄이다. $A_{perf}$가 완전체 환이므로 상대
Frobenius $F_{A_{perf}/A}$는 전사이다. 달리 말해
$A_{perf} = A[A_{perf}^p]$이고, 여기서 바로
$\Omega_{A_{perf}/A} = 0$이 따른다. 『사상 더 보기』의 보조정리
\ref{more-morphisms-lemma-formally-unramified-differentials}에 의해
$A \rightarrow A_{perf}$는 형식적 비분기이다.

\medskip\noindent
$A \rightarrow A_{perf}$가 형식적으로 매끄럽지 않음을 보이면 충분하다.
$A$는 매끄러운 $\mathbb{F}_p$-대수이므로 여접복합체
$L_{A/\mathbb{F}_P} \simeq \Omega_{A/\mathbb{F}_p}[0]$은 차수 $0$에
집중되어 있다. 『여접』의 보조정리
\ref{cotangent-lemma-when-projective}를 보라. 더욱이 『여접』의 보조정리
$L_{A_{perf}/\mathbb{F}_p} = 0$은 $D(A_{perf})$에서 성립한다. 이는
『여접』의 보조정리 \ref{cotangent-lemma-perfect-zero}에 의한 것이다.
여접복합체의 다음 구별 삼각형을 생각하자.
$$
L_{A/\mathbb{F}_p} \otimes_A A_{perf} \to
L_{A_{perf}/\mathbb{F}_p} \to
L_{A_{perf}/A} \to
(L_{A/\mathbb{F}_p} \otimes_A A_{perf})[1]
$$
이는 $D(A_{perf})$ 안의 삼각형이다. 『여접』의 절
\ref{cotangent-section-triangle}을 보라. 그러면
$L_{A_{perf}/A} = \Omega_{A/\mathbb{F}_p} \otimes_A A_{perf}[1]$이다.
즉 $L_{A_{perf}/A}$는 계수 $1$인 자유 $A_{perf}$-가군을 차수 $-1$에
놓은 것과 같다. 따라서 $A \rightarrow A_{perf}$는 형식적으로 매끄럽지
않다. 이는 『사상 더 보기』의 보조정리
\ref{more-morphisms-lemma-NL-formally-smooth}와 『여접』의 보조정리
\ref{cotangent-lemma-relation-with-naive-cotangent-complex}에 의한다.
\end{proof}

\noindent
다음 예도 소수 표수의 환을 사용하지만 조금 더 놀랍다. 단점은 앞 예보다
표수 $p$ 현상에 관한 지식을 더 요구한다는 것이다. 소수 표수의 환 $A$에서
$F$-유한이라 함은 $A$ 위의 Frobenius 사상이 유한함을 뜻함을 상기하라.

\begin{lemma}
\label{examples-lemma-completion-etale}
$(A, \mathfrak m, \kappa)$를 표수 $p > 0$의 뇌터 국소환이라 하고
$[\kappa : \kappa^p] < \infty$라 하자. 표준 사상 $A \to A^\wedge$, 즉
$A$의 완비화로 가는 사상은 평탄하고 형식적 비분기이다. 그러나 $A$가 정칙이지만
우수하지 않으면 이 사상은
형식적 \'etale이 아니다.
\end{lemma}

\begin{proof}
완비화의 평탄성은 『대수학』의 보조정리
\ref{algebra-lemma-completion-flat}이다. 형식적 비분기임을 보이려면
$\Omega_{A^\wedge/A} = 0$임을 보이면 충분하다. 『대수학』의 보조정리
\ref{algebra-lemma-characterize-formally-unramified}을 보라.

\medskip\noindent
증명을 개설한다. $x_1, \ldots, x_r \in A$를 택하되, 그 상이
$p$-기저 $\overline{x}_1, \ldots, \overline{x}_r$가 되어 $\kappa$를
생성하게 하자. 즉 $\kappa$는 $\overline{x}_i$에 의해 $\kappa^p$ 위에서
최소 생성된다. $y_1, \ldots, y_s$를 $\mathfrak m$의 최소 생성원 집합으로
택하자. 각 $n$에 대해 원소 $x_1, \ldots, x_r, y_1, \ldots, y_s$는
$A/\mathfrak m^n$을 $(A/\mathfrak m^n)^p$ 위에서 Frobenius를 통해
생성한다. 일부 세부 사항은 생략한다. 따라서
$F : A^\wedge \to A^\wedge$는 유한이다. 이에 따라
$\Omega_{A^\wedge/A}$는 유한 $A^\wedge$-가군이다. 한편 임의의
$a \in A^\wedge$와 $n$에 대해 $a_0 \in A$를 택하여
$a - a_0 \in \mathfrak m^nA^\wedge$가 되게 할 수 있다.
따라서 $\text{d}(a) \in \bigcap \mathfrak m^n \Omega_{A^\wedge/A}$이고,
『대수학』의 보조정리
\ref{algebra-lemma-intersect-powers-ideal-module-zero}에 의해
$\text{d}(a)$는 영이다. 그러므로 $\Omega_{A^\wedge/A} = 0$이다.

\medskip\noindent
$A$가 정칙이라고 하자. Cohen 구조정리를 쓰면
$x_1, \ldots, x_r, y_1, \ldots, y_s$는 $p$-기저로서 $A^\wedge$를
생성한다. 즉
$$
A^\wedge = \bigoplus\nolimits_{I, J} (A^\wedge)^p
x_1^{i_1} \ldots x_r^{i_r} y_1^{j_1} \ldots y_s^{j_s}
$$
이고 $I = (i_1, \ldots, i_r)$, $J = (j_1, \ldots, j_s)$이며
$0 \leq i_a, j_b \leq p - 1$이다. 세부 사항은 생략한다. 특히
$\Omega_{A^\wedge}$는 자유 $A^\wedge$-가군이고, 그 기저는
$\text{d}(x_1), \ldots, \text{d}(x_r), \text{d}(y_1), \ldots, \text{d}(y_s)$이다.

\medskip\noindent
이제 $A \to A^\wedge$가 형식적 \'etale이거나
심지어 형식적으로 매끄럽기만 해도, 『대수학』의 명제
\ref{algebra-proposition-characterize-formally-smooth}에 의해
$\NL_{A^\wedge/A}$의 차수 $-1, 0$ 코호몰로지는 영이다.
환 준동형 $\mathbf{F}_p \to A \to A^\wedge$의 Jacobi-Zariski 열
(『대수학』의 보조정리 \ref{algebra-lemma-exact-sequence-NL})에서
$\Omega_A \otimes_A A^\wedge \cong \Omega_{A^\wedge}$를 얻는다.
따라서 $\Omega_A$는
$\text{d}(x_1), \ldots, \text{d}(x_r), \text{d}(y_1), \ldots, \text{d}(y_s)$를
기저로 하는 자유 가군이다. 분수체를 살피고 $A$가 정규임을 쓰면
$a \in A$가 $p$제곱일 필요충분조건은 $A^\wedge$에서 그 상이 $p$제곱인
것임을 알 수 있다(세부 사항은 생략한다.
『대수학』의 보조정리
\ref{algebra-lemma-derivative-zero-pth-power}를 사용하라).
둘째 결과로 연산자 $\partial/\partial x_a$와
$\partial/\partial y_b$가 $A$ 위에 정의된다.

\medskip\noindent
이들로부터 $A$가 $A^p$ 위에서 유한이고 $p$-기저가
$x_1, \ldots, x_r, y_1, \ldots, y_s$라는 결론이 나옴을 보이겠다.
이는 $A$가 우수하지 않다는 가정과 모순이다. Kunz의 결과
\cite[Corollary 2.6]{Kun76}를 보라. 이제 $a \in A$라 하고
$$
a =  \sum\nolimits_{I, J} (a_{I, J})^p
x_1^{i_1} \ldots x_r^{i_r} y_1^{j_1} \ldots y_s^{j_s}
$$
로 쓰자. 여기서 $a_{I, J} \in A^\wedge$이다. 모든 $a_{I, J} \in A$임을
보이면 증명이 끝난다. 양변에 연산자
$$
(\partial/\partial x_1)^{p - 1} \ldots
(\partial/\partial x_r)^{p - 1}
(\partial/\partial y_1)^{p - 1} \ldots
(\partial/\partial y_s)^{p - 1}
$$
를 적용하면 $a_{I, J}^p \in A$이다. 여기서
$I = (p - 1, \ldots, p - 1)$이고 $J = (p - 1, \ldots, p - 1)$이다.
위 관찰에 의해 $a_{I, J} \in A$이다. $a$를
$a' = a - a_{I, J}^p x^I y^J$로 바꾼 뒤 차수가 $1$ 낮은 미분 연산자를 써서
다른 계수 $a_{I, J}$도 $A$에 속하게 할 수 있다.
이를 계속한다.
\end{proof}

\begin{remark}
\label{examples-remark-reference-existence-regular-nonexcellent-rings}
잉여체가 유한 $p$-기저를 갖는 우수하지 않은 정칙환은
$\mathbb{P}^2_k$의 함수체 안에서도 구성할 수 있다. 여기서 밑체는 표수
$p$인 $k = \overline{k}$이다.
\cite[$\mathsection 4.1$]{DS18}을 보라.
\end{remark}

\noindent
보조정리 \ref{examples-lemma-completion-etale}의 증명은 실제로 조금 더 말해 준다.

\begin{lemma}
\label{examples-lemma-excellent-regular-local-rings}
$(A, \mathfrak m, \kappa)$를 표수 $p > 0$인 정칙 국소환이라 하고
$[\kappa : \kappa^p] < \infty$라 하자. $A$가 우수일 필요충분조건은
$A \to A^\wedge$가 형식적 \'etale인 것이다.
\end{lemma}

\begin{proof}
역방향 함의는 보조정리 \ref{examples-lemma-completion-etale}에서 따른다.
정방향을 위해, 보조정리 \ref{examples-lemma-completion-etale}에 의해
$A \to A^\wedge$가 형식적 비분기, 동치로 $\Omega_{A^\wedge/A}$가
영임을 이미 안다. 따라서 완비화 사상이 $A$가 우수일 때 형식적으로
매끄러움만 보이면 된다. N\'eron-Popescu 비특이화에 의해
$A \to A^\wedge$는 매끄러운 $A$-대수들의 여과 쌍대극한으로
쓸 수 있다(『환 준동형의 평활화』의 정리
\ref{smoothing-theorem-popescu}). 따라서 $\NL_{A^\wedge/A}$의 차수 $-1$
코호몰로지는 영이다. 『대수학』의 명제
\ref{algebra-proposition-characterize-formally-smooth}에 의해
$A \to A^\wedge$는 형식적으로 매끄럽다.
\end{proof}

\section{멱등원들의 집합으로 생성된 아이디얼과 국소화}
\label{examples-section-ideal-locally-idempotents}

\noindent
$R$을 환이라 하고 환
$$
B(R) = R[x_n; n \in \mathbf{Z}]/(x_n(x_n - 1), x_nx_m; n \not = m)
$$
을 생각하자. 모든 소 아이디얼 $\mathfrak q \subset B(R)$은
$$
\mathfrak q = \mathfrak pB(R) + (x_n; n \in \mathbf{Z})
$$
꼴이거나
$$
\mathfrak q =
\mathfrak pB(R) + (x_n - 1) + (x_m; n \not = m, m \in \mathbf{Z}).
$$
꼴임을 쉽게 보일 수 있다. 따라서
$$
\Spec(B(R)) =
\Spec(R) \amalg \coprod\nolimits_{n \in \mathbf{Z}} \Spec(R)
$$
이다. 여기서 위상은 단순한 서로소 합 위상이 아니며 다음 성질을 갖는다.
$n \in \mathbf{Z}$로 첨자 붙인 각 복사본은 열린 부분스킴이며, 구체적으로
표준 열린집합 $D(x_n)$이다. ``가운데'' $\Spec(R)$의 복사본은 다른
$\Spec(R)$의 복사본 중 임의의 무한 개의 합집합의 폐포에 속한다.
이 마지막 $\Spec(R)$ 복사본은 아이디얼 $(x_n, n \in \mathbf{Z})$로
잘리며, 이 아이디얼은 멱등원 $x_n$들로 생성된다. 따라서 $\Spec(R)$이 연결이면 위 분해는
$\Spec(B(R))$의 연결 성분들로의 분해와 정확히 같다.

\medskip\noindent
이제 $A = \mathbf{C}[x, y]/((y - x^2 + 1)(y + x^2 - 1))$라 하자.
$A$의 스펙트럼은 두 기약 성분 $C_1 = \Spec(A_1)$,
$C_2 = \Spec(A_2)$로 이루어진다. 여기서
$A_1 = \mathbf{C}[x, y]/(y - x^2 + 1)$,
$A_2 = \mathbf{C}[x, y]/(y + x^2 - 1)$이다. 이들은 각각
$(x, y) = (t, t^2 - 1)$과 $(x, y) = (t, -t^2 + 1)$로 매개화되고
$P = (-1, 0)$과 $Q = (1, 0)$에서 만난다. $B(A)$의 꼬인 판본을 만들되
$B(A_1)$을 $B(A_2)$에 다음과 같이 붙일 수 있다. $P$ 위에서는
$x_n \in B(A_1) \otimes \kappa(P)$를
$x_n \in B(A_2) \otimes \kappa(P)$와 대응시키지만, $Q$ 위에서는
$x_n \in B(A_1) \otimes \kappa(Q)$를
$x_{n + 1} \in B(A_2) \otimes \kappa(Q)$와 대응시킨다.
$B^{twist}(A)$는 이렇게 얻어진 $A$-대수를 나타낸다. 세부 사항은 생략한다.
구성상 $B^{twist}(A)$는
$A$ 위 Zariski 국소적으로 꼬이지 않은 판본과 동형이다. 실제로 이는 두
주 열린집합 $\Spec(A) \setminus \{P\}$와
$\Spec(A) \setminus \{Q\}$ 위에서 성립한다. 그러나 우리의 붙임은 충분한
``모노드로미''를 만들어 $\Spec(B^{twist}(A))$가 연결이 되게 한다
(세부 사항은 생략한다). 끝으로 가운데 복사본
$\Spec(A) \to \Spec(B^{twist}(A))$가 있다. 이는 그 아이디얼이
$B^{twist}(A)$ 위 Zariski 국소적으로 멱등원들로 생성된 아이디얼에 의해
잘리지만, 전체적으로는 그렇지 않은 닫힌 부분스킴을 준다
($B^{twist}(A)$에는 자명하지 않은 멱등원이 없기 때문이다).

\begin{lemma}
\label{examples-lemma-not-generated-idempotents}
아핀 스킴 $X = \Spec(A)$와 닫힌 부분스킴 $T \subset X$가 존재하여,
$T$는 $X$ 위 Zariski 국소적으로 멱등원들로 생성된 아이디얼에 의해 잘리지만
$T$는 멱등원들로 생성된 하나의 아이디얼에 의해 잘리지는 않는다.
\end{lemma}

\begin{proof}
위를 보라.
\end{proof}

\section{국소환들을 동일시하지만 ind-\'etale이 아닌 환 준동형}
\label{examples-section-not-ind-etale}

\noindent
절 \ref{examples-section-ideal-locally-idempotents}에서 구성한 환 준동형
$R \to B(R)$은 $R$의 유한 곱들의 쌍대극한이다. 따라서 $R \to B(R)$은
ind-Zariski이다. 『Pro-\'etale 코호몰로지』의 정의
\ref{proetale-definition-ind-zariski}를 보라. 이제 절
\ref{examples-section-ideal-locally-idempotents}에서 구성한
환 준동형 $A \to B^{twist}(A)$를 생각하자. 이 환 준동형은
$\Spec(A)$ 위 Zariski 국소적으로 ind-Zariski 환 준동형 $R \to B(R)$과
동형이므로 국소환들을 동일시한다(『Pro-\'etale 코호몰로지』의 보조정리
\ref{proetale-lemma-ind-zariski-implies}를 보라). 절
\ref{examples-section-ideal-locally-idempotents}의 논의는 핵이
멱등원들로 생성되지 않는 단면 $B^{twist}(A) \to A$가 있음을 보인다.
만일 $A \to B^{twist}(A)$가 ind-\'etale, 즉
$B^{twist}(A) = \colim A_i$이고 $A \to A_i$가 \'etale이라면,
$A_i \to A$의 핵은 멱등원 하나로 생성될
것이다(『대수학』의 보조정리 \ref{algebra-lemma-map-between-etale}과
\ref{algebra-lemma-surjective-flat-finitely-presented}). 이는 위 결과와 모순이다.

\begin{lemma}
\label{examples-lemma-not-ind-etale}
국소환들을 동일시하지만 ind-\'etale이 아닌 환 준동형 $A \to B$가 존재한다.
따라서 더더욱 ind-Zariski가 아니다.
\end{lemma}

\begin{proof}
위의 논의를 보라.
\end{proof}


\section{단사 가군에 연관된 비플라스크 준연접 층}
\label{examples-section-nonflasque}

\noindent
이런 종류의 예를 더 보려면 Illusie가 Verdier의 몇 가지 예를 설명하는
\cite[Expos\'e II, Appendix I]{SGA6}을 보라.

\medskip\noindent
아핀 스킴 $X = \Spec(A)$를 생각하자. 여기서
$$
A = k[x, y, z_1, z_2, \ldots]/(x^nz_n)
$$
는 『성질』의 예 \ref{properties-example-does-not-work-in-general}에 나오는
환이다. $I = (x) \subset A$라 하자. $U = D(x)$라는 $X$의 준콤팩트
열린집합을 생각하자. 앞서\ 인용한\ 곳에서 단면 $s \in \mathcal{O}_X(U)$가 있으며
그 어떤 $A$-가군 사상 $I^n \to A$에서도 오지 않음을 보았다. 여기서
$n \geq 0$는 임의이다.

\medskip\noindent
$\alpha : A \to J$를 $A$에서 단사 $A$-가군으로 가는 매장이라 하자.
$Q = J/\alpha(A)$라 하고 몫사상을 $\beta : J \to Q$라 쓰자. 사상
$$
\Gamma(X, \widetilde{J})
\longrightarrow
\Gamma(U, \widetilde{J})
$$
은 전사가 아니라고 주장한다. 실제로 $\alpha(s)$는 그 상에 속하지 않는다.
반대로 원소 $x \in J$가 $\alpha(s)$로 $U$ 위에서 제한된다고 하자.
그러면 $\beta(x) \in Q$는 $0$으로 $U$ 위에서 제한된다. 따라서
$I^n\beta(x) = 0$이며, 여기서 $n$은 어떤 정수이다. 『성질』의 보조정리
\ref{properties-lemma-sections-over-quasi-compact-open-in-affine}를 보라.
그러면 $\varphi : I^n \to A$, $h \mapsto \alpha^{-1}(hx)$라는 사상을 얻는다.
이 $\varphi$가 『성질』의 보조정리
\ref{properties-lemma-sections-over-quasi-compact-open-in-affine}의 사상을
통해 단면 $s$를 준다는 것은 쉽게
보이는데, 이는 모순이다.

\begin{lemma}
\label{examples-lemma-nonflasque}
아핀 스킴 $X = \Spec(A)$와 단사 $A$-가군 $J$가 존재하여
$\widetilde{J}$는 $X$ 위의 플라스크층이 아니다. 심지어 제한
$\Gamma(X, \widetilde{J}) \to \Gamma(U, \widetilde{J})$도 $U$가 표준
열린집합일 때 전사일 필요가 없다.
\end{lemma}

\begin{proof}
위를 보라.
\end{proof}

\noindent
사실 비슷한 구성으로, 연관 준연접 층이 준콤팩트 열린집합 위에서 영이 아닌
코호몰로지를 갖는 단사 가군의 예를 얻을 수 있다. 환
$$
A = k[x, y, w_1, u_1, w_2, u_2, \ldots]/(x^nw_n, y^nu_n, u_n^2, w_n^2)
$$
에서 시작하자. 여기서 $k$는 체이다. 단사 사상 $A \to I$를 택하자.
여기서 $I$는 단사 $A$-가군이다. 원소 $1/xy$는
$A_{xy} \subset I_{xy}$에 속하지만 $I_x \oplus I_y \to I_{xy}$의 상에는
속하지 않는다고 주장한다. 반대로
$$
\frac{1}{xy} = \frac{i}{x^n} + \frac{j}{y^n}
$$
이라고 하자. 여기서 $n \geq 1$이고 $i, j \in I$이다. 분모를 없애면
$$
(xy)^{n + m - 1} = x^my^{n + m}i + x^{n + m}y^mj
$$
을 얻으며, 여기서 $m \geq 0$이다. 여기에 $u_{n + m}w_{n + m}$을 곱하면
$u_{n + m}w_{n + m}(xy)^{n + m - 1} = 0$을 $A$에서 얻는데, 이것이 원하는
모순이다. $U = D(x) \cup D(y) \subset X = \Spec(A)$라 하자. 모든 $A$-가군
$M$에 대해 Mayer-Vietoris에 의한 완전열
$$
0 \to H^0(U, \widetilde{M}) \to M_x \oplus M_y \to M_{xy} \to
H^1(U, \widetilde{M}) \to 0
$$
이 있다. 따라서 $H^1(U, \widetilde{I})$는 영이 아니다.

\begin{lemma}
\label{examples-lemma-nonvanishing}
밑 위상 공간이 뇌터인 아핀 스킴 $X = \Spec(A)$와 단사 $A$-가군 $I$가
존재하여, $\widetilde{I}$는 영이 아닌 $H^1$을 어떤 준콤팩트 열린집합
$U$ 위에서 갖는다. 여기서 이 열린집합은 $X$의 열린 부분이다.
\end{lemma}

\begin{proof}
위를 보라. 위상 공간으로서 $\Spec(A) = \Spec(k[x, y])$임에 주의하라.
\end{proof}




\section{분리되지 않는 평탄 군 스킴}
\label{examples-section-non-separated-group-scheme}

\noindent
체 위의 모든 군 스킴은 분리되어 있다. 『군체』의 보조정리
\ref{groupoids-lemma-group-scheme-over-field-separated}를 보라.
그러나 밑 위의 군 스킴에 대해서는 그렇지 않다.

\medskip\noindent
$k$를 체라 하자. $S = \Spec(k[x]) = \mathbf{A}^1_k$라 하자.
$G$를 원점 $0$을 두 겹으로 만든 아핀 직선
(『스킴』의 예 \ref{schemes-example-affine-space-zero-doubled}를 보라)을
$S$ 위의 스킴으로 본 것으로 하자. 따라서 $G \to S$의 섬유는 한 원소만
갖거나 두 원소를 갖는다(하나는 $U$에, 하나는 $V$에 있다).
그러므로 이 섬유들에 군 구조를 줄 수 있다($U$에 있는 원소를 군 구조의
영원소로 정한다). 더 정확히 말해, $G$에는 $S$로 동형사상되는 두 열린집합
$U, V$가 있고, $U \cap V$는 $S \setminus \{0\}$으로 동형사상된다. 그러면
$$
G \times_S G = U \times_S U \cup V \times_S U \cup
U \times_S V \cup V \times_S V
$$
이며 각 조각은 $S$와 동형이다. 따라서 첫째와 마지막 조각을 $U$로,
가운데 두 조각을 $V$로 보내는 유일한 $S$-사상으로 곱셈
$m : G \times_S G \to G$를 정의할 수 있다. 이는 위에서 준 점별 설명과
일치한다. 이것이 군 스킴 구조를 정의한다는 확인은 생략한다.

\begin{lemma}
\label{examples-lemma-non-separated-group-scheme}
아핀 직선 위에 유한형인 평탄 군 스킴으로서 분리되지 않는 것이 존재한다.
\end{lemma}

\begin{proof}
위 논의를 보라.
\end{proof}

\begin{lemma}
\label{examples-lemma-non-quasi-separated-group-scheme}
무한 차원 아핀 공간 위에 유한형인 평탄 군 스킴으로서 준분리되지 않는 것이
존재한다.
\end{lemma}

\begin{proof}
위와 같은 구성을 무한 차원 아핀 공간
$S = \mathbf{A}^\infty_k = \Spec k[x_1, x_2, \ldots]$을 밑으로 하고,
극대 아이디얼 $(x_1, x_2, \ldots)$에 대응하는 원점 $0 \in S$를 $G$에서
두 겹으로 만드는 닫힌점으로 삼아 수행할 수 있다. 그 결과인 군 스킴
$G \rightarrow S$는 『스킴』의 예
\ref{schemes-example-not-quasi-separated}에서 설명한 대로 준분리되지 않는다.
\end{proof}



\section{항등성분은 평탄하지만 평탄하지 않은 군 스킴}
\label{examples-section-non-flat-group-scheme}

\noindent
$X \to S$를 스킴의 단사사상이라 하자. $G = S \amalg X$라 하자.
$S$-사상 $m : G \times_S G \to G$를
$$
G \times_S G = X \times_S X \amalg X \amalg X \amalg S
\longrightarrow G = X \amalg S
$$
에서 성분 $X \times_S X$와 $S$는 $S$로 보내고, 두 성분 $X$는 항등사상으로
$X$에 보내는 사상으로 정하자. 이는 군 법칙을 정의한다. 이를 보이려면
$G \times_S G \times_S G \to G$인 사상으로서
$m \circ (m \times \text{id}_G) = m \circ (\text{id}_G \times m)$임을 보여야 한다.
$G \times_S G \times_S G$를 위와 같이 성분들로 분해하면, 각 $8$개 조각에
제한한 등식을 확인하면 됨을 알 수 있다. 각 조각은 $S$, $X$,
$X \times_S X$, 또는 $X \times_S X \times_S X$ 가운데 하나와 동형이다.
또한 두 사상은 각각 이 조각들을 $S$, $X$, $S$, $X$로 보낸다.
그런데 $X \to S$가 단사사상이므로 $X \times_S X \cong X$이고
$X \times_S X \times_S X \cong X$이며, 각 경우에 $S$-사상 $S \to S$ 또는
$X \to X$가 정확히 하나씩 있다. 따라서
$m \circ (m \times \text{id}_G) = m \circ (\text{id}_G \times m)$임을 알 수 있다.
이제 $X \to S$를 평탄하지 않은 임의의 스킴 단사사상(예를 들어 닫힌
몰입)으로 택하면, 항등성분은 $S$라서 평탄하지만 그 자체는 평탄하지 않은
밑 $S$ 위의 군 스킴을 얻는다.

\begin{lemma}
\label{examples-lemma-non-flat-group-scheme}
어떤 밑 $S$ 위에 군 스킴 $G$가 존재하여, 그 항등성분은 $S$ 위에서
평탄하지만 그 자체는 $S$ 위에서 평탄하지 않다.
\end{lemma}

\begin{proof}
위 논의를 보라.
\end{proof}




\section{체 위의 분리되지 않는 군 대수공간}
\label{examples-section-non-separated-group-space}

\noindent
체 위의 모든 군 스킴은 분리되어 있다. 『군체』의 보조정리
\ref{groupoids-lemma-group-scheme-over-field-separated}를 보라.
그러나 체 위의 군 대수공간에 대해서는 그렇지 않다
(긍정적인 결과는 이 절 끝을 보라).

\medskip\noindent
$k$를 표수가 영인 체라 하자. 『공간』의 예
\ref{spaces-example-affine-line-translation}에 나오는 대수공간
$G = \mathbf{A}^1_k/\mathbf{Z}$를 생각하자. 구성상 $G$는 $k$ 위 스킴의
범주에서 준층
$$
T \longmapsto \Gamma(T, \mathcal{O}_T) / \mathbf{Z}
$$
에 연관된 fppf 층이다. 준층 위의 자명한 덧셈 법칙은 덧셈
$m : G \times G \to G$를 유도하여 $G$를 $\Spec(k)$ 위의 군 대수공간으로
만든다. $G$는 분리되어 있지 않다(심지어 준분리되거나 국소적으로
분리되어 있지도 않다). 한편 $G \to \Spec(k)$는 유한형이다!

\begin{lemma}
\label{examples-lemma-non-separated-group-space}
체 위에 유한형인 군 대수공간으로서 분리되어 있지 않은 것이 존재한다
(심지어 준분리되거나 국소적으로 분리되어 있지도 않다).
\end{lemma}

\begin{proof}
위 논의를 보라.
\end{proof}

\noindent
긍정적인 결과는 다음과 같다. 군 대수공간 $G$가 준분리되거나,
국소적으로 분리되거나, 더 일반적으로 양호한 대수공간이면 $G$는 실제로
분리되어 있다. 『공간에서의 군체에 관하여 더 자세히』의 보조정리
\ref{spaces-more-groupoids-lemma-group-scheme-over-field-separated}를 보라.
또한 체 위의 유한형 분리 군 대수공간은 사실 스킴이다. 『공간에서의
군체에 관하여 더 자세히』의 보조정리
\ref{spaces-more-groupoids-lemma-group-space-scheme-locally-finite-type-over-k}를
보라. 증명의 발상은 스킴 자취가 열린 조밀집합이라는 것이다. 『공간의
성질』의 명제
\ref{spaces-properties-proposition-locally-quasi-separated-open-dense-scheme}
또는 『양호한 공간』의 정리
\ref{decent-spaces-theorem-decent-open-dense-scheme}를 보라.
이 열린집합을 평행이동하면 $G$의 모든 점이 스킴인 열린 근방을 갖는다는
것을 알 수 있다.




\section{\'etale 사상의 한 섬유 안에서 점들 사이의 특수화}
\label{examples-section-specializations-fibre-etale}

\noindent
$f : X \to Y$가 스킴의 \'etale 사상, 또는 더 일반적으로 국소 준유한
사상이면 섬유의 점들 사이에는 특수화가 없다. 『사상』의 보조정리
\ref{morphisms-lemma-locally-quasi-finite-fibres}를 보라.
그러나 대수공간의 사상에 대해서는 일반적으로 성립하지 않는다.

\medskip\noindent
예를 주기 위해 $k$를 체라 하자. 다음과 같이 놓자.
$$
P = k[u, u^{-1}, y, \{x_n\}_{n \in \mathbf{Z}}].
$$
$\tau(u) = u$, $\tau(y) = uy$, $\tau(x_n) = x_{n + 1}$이라는 규칙으로 주어진
자기동형사상 $\tau$가 생성하는 $k$-대수 사상들에 의한 $\mathbf{Z}$의
$P$ 위 작용을 생각하자. $d \geq 1$에 대해
$I_d = ((1 - u^d)y, x_n - x_{n + d}, n \in \mathbf{Z})$라 놓자.
그러면 $V(I_d) \subset \Spec(P)$는 $\tau^d$의 고정점 자취이다.
$S \subset P$를 $y$와 모든 $1 - u^d$, $d \in \mathbf{N}$이 생성하는
곱셈적 부분집합이라 하자. 그러면 $\mathbf{Z}$가
$U = \Spec(S^{-1}P)$ 위에 자유롭게 작용함을 알 수 있다.
$X = U/\mathbf{Z}$를 몫 대수공간이라 하자. 『공간』의 정의
\ref{spaces-definition-quotient}를 보라.

\medskip\noindent
$S^{-1}P$의 소 아이디얼
$\mathfrak p_n = (x_n, x_{n + 1}, \ldots)$을 생각하자.
$\tau(\mathfrak p_n) = \mathfrak p_{n + 1}$임에 주의하라. 따라서 이들은 각각
점 $\xi_n \in U$를 정의하고, 이 점들의 $X$에서의 상은 모두 동일한
$X$의 점 $x$이다. 더욱이 다음 특수화들이 있다.
$$
\ldots \leadsto \xi_n \leadsto \xi_{n - 1} \leadsto \ldots
$$
따라서 $U \to X$가 약속한 유형의 예이다.

\begin{lemma}
\label{examples-lemma-specializations-fibre-etale}
대수공간의 \'etale 사상 $f : X \to Y$와 $|X|$ 안의 점들의 자명하지 않은
특수화 $x \leadsto x'$가 존재하여 $|Y|$에서 $f(x) = f(x')$이다.
\end{lemma}

\begin{proof}
위 논의를 보라.
\end{proof}





\section{fppf 토서가 아닌 토서}
\label{examples-section-torsor-not-fppf}

\noindent
『군체』의 주석 \ref{groupoids-remark-fun-with-torsors}에서 임의의
$G$-토서가 fppf 위상에 대한 $G$-토서인지 물었다. 이 절에서는 항상
그렇지는 않음을 보인다.

\medskip\noindent
$k$를 체라 하자. 이하 모든 스킴과 스택은 $k$ 위에 있다.
$G \to \Spec(k)$를 다음 군 스킴이라 하자.
$$
G = (\mu_{2, k})^\infty =
\mu_{2, k} \times_k \mu_{2, k} \times_k \mu_{2, k} \times_k \ldots =
\lim_n (\mu_{2, k})^n
$$
여기서 $\mu_{2, k}$는 $\Spec(k)$ 위의 이차 단위근 군 스킴이다.
『군체』의 예 \ref{groupoids-example-roots-of-unity}를 보라.
아핀 스킴들의 역극한이므로 $G$는 아핀 군 스킴이다. 실제로 이는 환
$k[t_1, t_2, t_3, \ldots]/(t_i^2 - 1)$의 스펙트럼이다. 곱셈 사상
$m : G \times_k G \to G$는 대수 수준에서 $t_i \mapsto t_i \otimes t_i$로
주어진다.

\medskip\noindent
$k$ 위의 임의의 $G$-토서는 어떤 $a_i \in k^*$에 대해
$$
P = \Spec(k[x_1, x_2, x_3, \ldots]/(x_i^2 - a_i))
$$
꼴이며, $G$-작용 $G \times_k P \to P$는 대수 수준에서
$x_i \to t_i \otimes x_i$로 주어진다고 주장한다. 증명은 생략한다.
사실 이 예에 필요한 것은 $P$가 $G$-토서라는 사실뿐이며, 이는
$k' = k(\sqrt{a_1}, \sqrt{a_2}, \ldots)$ 위에서 $P$가 $G$와
$G$-동치적으로 동형이 되므로 분명하다. $P$가 자명할 필요충분조건은
$k' = k$이다. 실제로 $P$에 $k$-유리점이 있으면 모든 $a_i$가 제곱이다.

\medskip\noindent
$P$가 fppf 토서일 필요충분조건은 체 확대
$k' = k(\sqrt{a_1}, \sqrt{a_2}, \ldots)/k$가 유한인 것이라고 주장한다.
$k'$가 $k$ 위에서 유한이면 $\{\Spec(k') \to \Spec(k)\}$는 $P$를
자명하게 만드는 fppf 덮개이므로 $P$는 실제로 fppf 토서이다. 반대로
$P$가 fppf $G$-토서라고 하자. 이는 각 $P_{S_i}$가 자명해지는 fppf 덮개
$\{S_i \to \Spec(k)\}$가 존재한다는 뜻이다. $S_i$가 공집합이 아닌
$i$ 하나를 택하고 닫힌점 $s \in S_i$를 택하자. 『다양체』의 보조정리
\ref{varieties-lemma-locally-finite-type-Jacobson}에 의해 체 확대
$\kappa(s)/k$는 유한이고, 구성상 $P_{\kappa(s)}$에는
$\kappa(s)$-유리점이 있다. 따라서 $k \subset k' \subset \kappa(s)$이고
$k'$는 $k$ 위에서 유한이다.

\medskip\noindent
구체적인 예를 얻으려면 $k = \mathbf{Q}$, $a_i = i$로 택하면 된다
(원한다면 $a_i$를 $i$번째 소수로 택해도 된다).

\begin{lemma}
\label{examples-lemma-torsors-principal-spaces-not-equal}
$S$를 스킴이라 하고 $G$를 $S$ 위의 군 스킴이라 하자.
주동차 $G$-공간을 분류하는 스택 $G\textit{-Principal}$
(『스택의 예』의 소절
\ref{examples-stacks-subsection-principal-homogeneous-spaces}를 보라)과
fppf $G$-토서를 분류하는 스택 $G\textit{-Torsors}$
(『스택의 예』의 소절 \ref{examples-stacks-subsection-fppf-torsors}를 보라)은
일반적으로 동치가 아니다.
\end{lemma}

\begin{proof}
위 논의는 함자 $G\textit{-Torsors} \to G\textit{-Principal}$이
일반적으로 본질적으로 전사적이지 않음을 보여 준다.
\end{proof}








\section{준콤팩트 평탄 덮개를 갖지만 대수적이지 않은 스택}
\label{examples-section-not-algebraic-stack}

\noindent
이 절에서는 Brian Conrad가 제시한 예를 간단히 설명한다. 이 예는 온라인의
\href{https://mathoverflow.net/questions/15082/fpqc-covers-of-stacks/15269#15269}{이곳}에서
찾을 수 있다. 여기서 드는 예는 약간 다르다.

\medskip\noindent
$k$를 대수적으로 닫힌 체라 하자. 이하 모든 스킴과 스택은 $k$ 위에 있다.
$G \to \Spec(k)$를 아핀 군 스킴이라 하자. 『스택의 예』의 보조정리
\ref{examples-stacks-lemma-classifying-stacks}에서
$\mathcal{X} = [\Spec(k)/G]$를 $(\Sch/\Spec(k))_{fppf}$ 위의 군체 스택으로
보는 서로 동치인 여러 방법을 제시했다. 특히 $\mathcal{X}$는 fppf
$G$-토서를 분류한다. 더 정확히 말해, $1$-사상 $T \to \mathcal{X}$는
$T$ 위의 fppf $G_T$-토서 $P$에 대응하고, $2$-화살표는 토서의
동형사상에 대응한다. 따라서 대각 $1$-사상
$$
\Delta :
\mathcal{X}
\longrightarrow
\mathcal{X} \times_{\Spec(k)} \mathcal{X}
$$
은 표현가능하고 아핀이다. 실제로 스킴 $T/k$ 위의 임의의 fppf
$G_T$-토서 두 개 $P_1, P_2$에 대해 스킴 $\mathit{Isom}(P_1, P_2)$는
$T$ 위에서 아핀이다. $\Spec(k)$ 위의 자명한 $G$-토서는 $1$-사상
$$
f : \Spec(k) \longrightarrow \mathcal{X}.
$$
을 정의한다. 이것은 전사인 $1$-사상이라고 주장한다. 그 이유는 정의상
임의의 $1$-사상 $T \to \mathcal{X}$에 대해 $P_{T_i}$가 자명한
$G_{T_i}$-토서와 동형이 되는 fppf 덮개 $\{T_i \to T\}$가 존재하기
때문이다. 따라서 합성 $T_i \to T \to \mathcal{X}$는 $f$를 통해 분해된다.
그러므로 사영 $T \times_\mathcal{X} \Spec(k) \to T$가 전사임은 분명하다
(이것이 $f$가 전사라는 성질의 정의이다. 『대수적 스택』의 정의
\ref{algebraic-definition-relative-representable-property}를 보라).
비슷하게 $f$가 준콤팩트하고 평탄함을 보일 수 있다(세부사항은 생략한다).
또 다음 동형도 기록해 둔다.
$$
\Spec(k) \times_\mathcal{X} \Spec(k) \cong G
$$
이는 $k$ 위 스킴의 동형이다.

\medskip\noindent
$U$가 스킴인 전사 매끄러운 사상 $p : U \to \mathcal{X}$가 존재한다고 하자.
섬유곱
$$
\xymatrix{
W \ar[d] \ar[r] & U \ar[d] \\
\Spec(k) \ar[r] & \mathcal{X}
}
$$
을 생각하자. 그러면 $W$는 $k$ 위의 공집합이 아닌 매끄러운 스킴이므로
$k$-점을 갖는다. 이는 $f$를 $U$를 통해 분해할 수 있다는 뜻이다.
따라서
$$
G \cong
\Spec(k) \times_\mathcal{X} \Spec(k) \cong
(\Spec(k) \times_k \Spec(k))
\times_{(U \times_k U)}
(U \times_\mathcal{X} U)
$$
를 얻는다. 사영 $U \times_\mathcal{X} U \to U$들이 매끄럽다고
가정했으므로 『사상』의 보조정리
\ref{morphisms-lemma-permanence-finite-type}에 의해
$U \times_\mathcal{X} U \to U \times_k U$는 국소 유한형이다.
따라서 이 경우 $G$는 $k$ 위에서 국소 유한형이다. 종합하면 다음
보조정리를 증명했다(이는 상당히 일반화할 수 있다).

\begin{lemma}
\label{examples-lemma-BG-algebraic}
$k$를 체라 하고 $G$를 $k$ 위의 아핀 군 스킴이라 하자. 스택
$[\Spec(k)/G]$가 스킴에 의한 매끄러운 덮개를 가지면 $G$는 $k$ 위에서
유한형이다.
\end{lemma}

\begin{proof}
위 논의를 보라.
\end{proof}

\noindent
이 절의 제목과 같은 구체적인 예를 얻으려면 절
\ref{examples-section-torsor-not-fppf}의 군 스킴
$G = (\mu_{2, k})^\infty$를 택하면 된다. 이는 $k$ 위에서 국소 유한형이
아니다. 위 논의에 의해 $\mathcal{X} = [\Spec(k)/G]$는 『대수적 스택』의
정의 \ref{algebraic-definition-algebraic-stack}의 성질 (1), (2)를 갖지만
성질 (3)은 갖지 않는다. 따라서 $\mathcal{X}$는 대수적 스택이 아니다.
한편 스킴 $U$와 전사이고 평탄하고 준콤팩트인 사상
$U \to \mathcal{X}$는 존재한다. 바로 위에서 살펴본
$f : \Spec(k) \to \mathcal{X}$이다.






\section{대상에 대해서는 극한 보존이지만 극한 보존은 아닌 경우}
\label{examples-section-limit-preserving}

\noindent
$S$를 공집합이 아닌 스킴이라 하자. $\mathcal{G}$를
$(\Sch/S)_{fppf}$ 위의 단사 아벨 층이라 하자. 그러면 $S$ 위의 군체
스택
$$
\mathcal{G}\textit{-Torsors} \longrightarrow (\Sch/S)_{fppf}
$$
을 얻는다. 『스택의 예』의 보조정리
\ref{examples-stacks-lemma-torsors-sheaf-stack-in-groupoids}를 보라.
모든 $\mathcal{G}$-토서가 자명하므로 이 스택은 $(\Sch/S)_{fppf}$ 위에서
대상에 대해 극한을 보존한다(『표현가능성의 판정기준』의 절
\ref{criteria-section-limit-preserving}를 보라). 한편 $\mathcal{G}$는
층으로서 극한을 보존할 필요가 없으므로 $\mathcal{G}\textit{-Torsors}$는
일반적으로 극한을 보존하지 않는다(『Artin의 공리』의 정의
\ref{artin-definition-limit-preserving}를 보라). 예를 들어 영이 아닌
임의의 단사층 $\mathcal{I}$를 택하고
$\mathcal{G} = \prod_{n \in \mathbf{Z}} \mathcal{I}$라 놓으면 예를 얻는다.

\begin{lemma}
\label{examples-lemma-limit-preserving-on-objects-not-limit-preserving}
$S$를 공집합이 아닌 스킴이라 하자. 군체 스택
$p : \mathcal{X} \to (\Sch/S)_{fppf}$가 존재하여, $p$는 대상에 대해서는
극한을 보존하지만 $\mathcal{X}$ 자체는 극한을 보존하지 않는다.
\end{lemma}

\begin{proof}
위 논의를 보라.
\end{proof}



\section{대수적이지 않은 분류 스택}
\label{examples-section-non-algebraic}

\noindent
$S = \Spec(\mathbf{F}_p)$라 하고 $\mu_p$를 $S$ 위의 $p$차 단위근 군
스킴이라 하자. 『공간에서의 군체』의 절
\ref{spaces-groupoids-section-stacks}에서 몫 스택 $[S/\mu_p]$를
도입했고, 『스택의 예』의 절
\ref{examples-stacks-section-group-quotient-stacks}에서
$[S/\mu_p]$가 fppf $\mu_p$-토서의 분류 스택임을 보였다. 즉 $S$ 위의
스킴 $T$가 주어지면 범주 $\Mor_S(T, [S/\mu_p])$는 $T$ 위의 fppf
$\mu_p$-토서 범주와 표준적으로 동치이다. 끝으로 『표현가능성의
판정기준』의 정리
\ref{criteria-theorem-flat-groupoid-gives-algebraic-stack}에서
$[S/\mu_p]$가 대수적 스택임을 보았다.

\medskip\noindent
이제 다음 질문을 할 수 있다. ``\'etale $\mu_p$-토서를 분류하는 군체
범주 $\mathcal{S}$는 어떤가?'' (다시 말해 $\mathcal{S}$는 $\Sch/S$ 위의
범주이고, 스킴 $T$ 위의 섬유 범주는 $T$ 위의 \'etale $\mu_p$-토서
범주이다.)

\medskip\noindent
첫 번째 문제는 대상의 강하가 성립하지 않으므로 이것이 fppf 위상에 대한
스택이 아니라는 점이다. 예를 들어 $T = \Spec(\mathbf{F}_p(T))$ 위의
$\mu_p$-토서 $\Spec(\mathbf{F}_p(t)[x]/(x^p - t))$는 fppf 국소적으로
자명하지만 \'etale 국소적으로는 자명하지 않다.

\medskip\noindent
이 첫 문제를 고치려면 \'etale 위상에서 작업하면 되고, 이 경우에는
대상의 강하가 성립한다. 실제로 $\mathcal{S}$는 $(\Sch/S)_\etale$ 위의
군체 스택이다. 또한 대각
$\Delta : \mathcal{S} \to \mathcal{S} \times \mathcal{S}$는 (스킴으로)
표현가능하다. 이는 두 $\mu_p$-토서가 주어지면(\'etale 국소적으로
자명하든 아니든) 그들 사이의 동형사상 층이 스킴으로 표현가능하기
때문이다.

\medskip\noindent
따라서 마침내 스킴 $U$와 매끄럽고 전사인 $1$-사상
$U \to \mathcal{S}$가 존재하는지 물을 수 있다. 이것이 불가능함을 두
방법으로 보일 것이다. 하나는 직접적인 논증(독자에게는 건너뛰기를
권한다)이고, 다른 하나는 일반적인 결과를 쓰는 논증이다.

\medskip\noindent
직접 논증(개요): $1$-사상
$\mathcal{S} \to \Spec(\mathbf{F}_p)$는 형식적 매끄러움에 대한 무한소
올림 판정법을 만족함에 주의하라. 실제로 스킴의 일차 무한소 두꺼워짐
$T \to T'$가 주어지면 $\mu_p(T') \to \mu_p(T)$의 핵은 $T'$ 안에서
$T$를 정의하는 아이디얼 층의 단면들과 동형이므로
$H^1_\etale(T, \mu_p) = H^1_\etale(T', \mu_p)$이다. 또한 $\mathcal{S}$는
극한을 보존하는 스택이다. 따라서 $U \to \mathcal{S}$가 매끄러우면
$U \to \Spec(\mathbf{F}_p)$는 극한을 보존하고 형식적 매끄러움에 대한
무한소 올림 판정법을 만족한다. 이는 $U$가 $\mathbf{F}_p$ 위에서
매끄럽다는 것을 뜻한다. 특히 $U$는 축약 스킴이므로
$H^1_\etale(U, \mu_p) = 0$이다. 따라서 $U \to \mathcal{S}$는
$U \to \Spec(\mathbf{F}_p) \to \mathcal{S}$로 분해되고 첫 화살표는
매끄럽다. 매끄러움의 강하에 의해 $U \to \mathcal{S}$가 매끄러우면
$\Spec(\mathbf{F}_p) \to \mathcal{S}$도 매끄러워야 한다. 그러나
$\Spec(\mathbf{F}_p) \times_\mathcal{S} \Spec(\mathbf{F}_p)$는
$\mu_p$이고, 이는 $\Spec(\mathbf{F}_p)$ 위에서 매끄럽지 않으므로
그렇지 않다.


\medskip\noindent
구조적 논증: 『표현가능성의 판정기준』의 절
\ref{criteria-section-stacks-etale}에서 대수적 스택을, \'etale 위상에
대한 군체 스택 가운데 대각이 대수공간으로 표현가능하고 매끄러운 덮개를
갖는 것들로 생각할 수 있음을 보았다. 따라서 매끄러운 전사
$U \to \mathcal{S}$가 존재한다면 $\mathcal{S}$는 대수적 스택이고, 특히
fppf 위상에서 강하를 만족한다. 그러나 위에서 $\mathcal{S}$가 fppf
위상에서 강하를 만족하지 않음을 보았다.

\medskip\noindent
대략 말해 위 논증들은, 밑 $B$ 위의 군 스킴 $G$에 대해 \'etale
국소적으로 자명한 토서를 분류하는 \'etale 위상의 분류 스택이 대수적일
필요충분조건이 $G$가 $B$ 위에서 매끄러운 것임을 보여 준다. fppf
위상에서 작업하는 이점 하나는 $G \to B$가 평탄하고 국소 유한표현이라고
가정하는 것으로 충분하다는 점이다. 실제로 (fppf 위상에 대한) 몫 스택
$[B/G]$가 대수적일 필요충분조건은 $G \to B$가 평탄하고 국소
유한표현인 것이다. 『표현가능성의 판정기준』의 보조정리
\ref{criteria-lemma-BG-algebraic}를 보라.






\section{준콤팩트 평탄 덮개를 갖지만 대수적이지 않은 층}
\label{examples-section-not-algebraic}

\noindent
함자 $F = (\mathbf{P}^1)^\infty$를 생각하자. 즉 스킴 $T$에 대한 값
$F(T)$는 각 $f_i : T \to \mathbf{P}^1$가 스킴의 사상인
$f = (f_1, f_2, f_3, \ldots)$들의 집합이다. $\mathbf{P}^1$은 fpqc 덮개에
대한 층 성질을 만족한다. 『강하』의 보조정리
\ref{descent-lemma-fpqc-universal-effective-epimorphisms}를 보라.
층들의 곱은 층이므로 $F$도 fpqc 위상에 대한 층 성질을 만족한다.
$F$의 대각은 표현가능하다. 사상 $f : T \to F$와 $g : S \to F$가 주어지면
$T \times_F S$는 스킴 $T \times S$ 안의 닫힌 부분스킴
$T \times_{f_i, \mathbf{P}^1, g_i} S$들의 스킴론적 교집합이다.
전사 매끄러운 아핀 사상 $\text{SL}_2 \to \mathbf{P}^1$을 갖는 군 스킴
$\text{SL}_2$를 생각하자. 다음으로 $U = (\text{SL}_2)^\infty$와 그
표준적인 (곱) 사상 $U \to F$를 생각하자. $U$는 아핀 스킴임에 주의하라.
사상 $U \to F$가 평탄하고 전사이며 보편적으로 열린 사상이라고 주장한다.
실제로 $f : T \to F$를 사상이라 하자. 그러면 $Z = T \times_F U$는 $T$ 위의
스킴 $Z_i = T \times_{f_i, \mathbf{P}^1} \text{SL}_2$들의 무한 섬유곱이다.
각 사상 $Z_i \to T$는 전사이고 매끄럽고 아핀이므로
$$
Z = Z_1 \times_T Z_2 \times_T Z_3 \times_T \ldots
$$
는 $Z$ 위에서 평탄하고 아핀인 스킴이다. 간단한 극한 논증으로
$Z \to T$가 열린 사상이기도 함을 알 수 있다.

\medskip\noindent
한편 $F$는 대수공간이 아니라고 주장한다. 실제로 $F$가 대수공간이라면
(위에서 설명한 $F$ 위의 섬유곱에 의해) 준콤팩트 분리 대수공간일 것이다.
따라서 준연접 층의 코호몰로지는 어떤 차수 경계 위에서 소멸할 것이다
(『공간의 코호몰로지』의 명제
\ref{spaces-cohomology-proposition-vanishing} 및 그 앞의 주석들을 보라).
그러나 사영
$$
(\mathbf{P}^1)^\infty \longrightarrow (\mathbf{P}^1)^n
$$
아래에서 $\mathcal{O}(-2, -2, \ldots, -2)$를 당겨오면(이 사영에는 단면이
있다) 차수 $n$에서 코호몰로지가 영이 아닌 준연접 층을 얻을 수 있음은
분명하다. 종합하면 mathoverflow에서 Anton Geraschenko가 물은 질문의
답을 얻는다.

\begin{lemma}
\label{examples-lemma-not-algebraic}
함자 $F : \Sch^{opp} \to \textit{Sets}$가 존재하여 fpqc 위상에 대한 층
조건을 만족하고, 표현가능한 대각 $\Delta : F \to F \times F$를 가지며,
전사이고 평탄하고 보편적으로 열리고 준콤팩트인 사상 $U \to F$가
존재한다. 여기서 $U$는 스킴이지만 $F$는 대수공간이 아니다.
\end{lemma}

\begin{proof}
위 논의를 보라.
\end{proof}






\section{\'etale 위상과 자리스키 및 유한 \'etale 덮개의 비교}
\label{examples-section-etale-zariski-finite-etale}

\noindent
이 절에서는 C.~Deninger의 질문에 답하여 B.~Bhatt가 찾아낸 예를 제시한다.
이 예는 $\mathbf{C}$ 위의 대수다양체에서 \'etale 위상이 열린 몰입과 유한
\'etale 사상들로 생성된 위상보다 진정으로 더 섬세함을 보여 준다.
이런 예들은 전문가들에게 분명히 알려져 있으며, 어쩌면 Artin과
Grothendieck까지 거슬러 올라갈 것이다\footnote{전해지는 이야기에 따르면,
Grothendieck은 처음에는 자리스키 열린 덮개와 유한 \'etale 덮개를 써서
\'etale 위상을 정의하려 했으나, Artin이 (아마도 여기 기록한 것과 같은
예를 바탕으로) 달리 생각하도록 설득했다고 한다.}.
출판된 예를 알고 있다면
\href{mailto:stacks.project@gmail.com}{stacks.project@gmail.com}으로
전자우편을 보내 달라.

\medskip\noindent
$X$를 $\mathbf{C}$ 위의 매끄럽고 연결된 아핀 곡선이라 하자.
매끄럽고 연결된 곡선 $Y$에서 오는 차수 $3$의 유한 사상
$g: Y \to X$를 다음 성질을 갖도록 택하자. 어떤 점 $x \in X$가 있어서
그 섬유 $g^{-1}(x) = \{u, y\}$가 비분기점 $u$(차수 1)와 분기점
$y$(차수 2)로 이루어진다. $X$를 $x$의 자리스키 열린 근방으로 바꾸어
$g$가 $y$를 제외한 모든 곳에서 비분기라고 가정할 수 있다.
$U = Y \setminus \{y\}$라 하자. 그러면 $f=g|_U: U \to X$는 전사인
\'etale 사상이고 $f^{-1}(x) = \{u\}$이다.

\begin{lemma}
\label{examples-lemma-etale-not-dominated-zariski-finite-etale}
\'etale 덮개 $U \to X$는 열린 몰입과 유한 \'etale 사상들의 어떠한 유한
합성으로도 세분될 수 없다.
\end{lemma}

\begin{proof}
모순을 얻기 위해 다음 자료가 존재한다고 가정하자. 사상의 탑
$$
W_m \to W_{m-1} \to \dots \to W_1 \to W_0 = X
$$
에서 각 사상 $W_i \to W_{i - 1}$은 열린 몰입이거나 유한 \'etale
사상이고, $W_m \to X$를 $f : U \to X$를 따라 올린 사상 $W_m \to U$가
있으며, 점 $w \in W_m$는 $u \in U$(따라서 $x \in X$)로 간다고 하자.
각 $W_k$를 $w$의 상 $w_k \in W_k$를 포함하는 연결성분으로 바꾸어,
모든 $W_k$가 매끄럽고 연결되어 있다고 가정할 수 있다.

\medskip\noindent
각 $0 \le i \le m$에 대해 섬유곱 $Z_i = Y \times_X W_i$를 생각하자.
$W_i \to X$는 \'etale이고 $Y$는 매끄러우므로 곡선 $Z_i$는 매끄럽고
연결된 곡선들의 서로소 합이며, 밑변환에 의해 사상 $Z_i \to W_i$는
차수 $3$의 유한 사상이다. $w_i$ 위의 $Z_i \to W_i$의 섬유는 분명히
$Y$에서 각각 $y, u$로 가는 두 점 $z_i, v_i$로 이루어진다.
$Z_i \to W_i$는 $v_i$에서 비분기이고 $z_i$에서 차수 $2$로 분기한다.
$n_i = \#\pi_0(Z_i)$라 하자. $Z_i \to W_i$의 섬유에 점이 $2$개이므로
$n_i \in \{1, 2\}$이다. 또한 전이사상 $Z_i \to Z_{i - 1}$들이 우세하므로
$n_i$는 $i$와 함께 증가할 수만 있다. 분명히 $n_0 = 1$이다. 더욱이
$n_m \geq 2$여야 한다. $X$-사상 $W_m \to U \subset Y$가 있으므로
밑변환 $Z_m \to W_m$에는 단면이 있다. 따라서 $n_{i - 1} = 1$이고
$n_i = 2$인 최소 지표 $i \ge 1$이 있다. $n_j$들은 $Z_i$가 기약인지
여부로 결정되므로 사상 $W_i \to W_{i - 1}$은 열린 몰입일 수 없고,
따라서 유한 \'etale이어야 한다. $n_i = 2$이므로 매끄러운 곡선 $Z_i$는
두 연결성분을 갖는다. $W_i$로 가는 사상이 유한이므로 이 성분 가운데
하나는 $W_i$로 동형사상되어야 한다. 따라서
$Z_i = W_i \sqcup Z'_i$라 쓸 수 있고 $Z'_i \to W_i$는 차수 $2$이다.
유도된 사상 $W_i \subset Z_i \to Z_{i - 1}$은 연결된 곡선들 사이의 우세
사상이고, 둘 다 $W_{i - 1}$ 위에서 유한이다. 이런 사상은 항상
전사이다. 그러므로 $z_{i - 1}$로 가는 점 $t \in W_i$를 택할 수 있다.
이산 값매김환들의 확대
$$
\mathcal{O}_{W_{i - 1}, w_{i - 1}} \to
\mathcal{O}_{Z_i, z_i} \to
\mathcal{O}_{W_i, z}
$$
를 얻는다. 이제 첫 확대는 분기하지만 합성은 $W_i \to W_{i - 1}$이
\'etale이었으므로 비분기이다. 이는 불가능하므로 모순을 얻는다.
\end{proof}







\section{층과 특수화}
\label{examples-section-sheaves}

\noindent
이하에서는 『위상』의 정의 \ref{topologies-definition-big-etale-site}에서
구성한 큰 \'etale 사이트 $\Sch_\etale$를 고정한다. 또한 스킴은 이
사이트의 대상이라고 한다. 스킴 $X$의 점 $x, x'$에 대해
$x \in \overline{\{x'\}}$이면 $x$를 $x'$의 {\it 특수화}라고 부르며
$x' \leadsto x$라고 쓴다는 것을 상기하라. 특히 $x = x'$이면 이 조건이
성립한다.

\medskip\noindent
다음 규칙으로 정의한 함자 $F : \Sch_\etale \to \textit{Ab}$를 생각하자.
$$
F(X) = \prod\nolimits_{x \in X}
\prod\nolimits_{x' \in X, x' \leadsto x, x' \not = x}
\mathbf{Z}/2\mathbf{Z}
$$
스킴 $X$가 주어졌을 때 그 점들의 밑집합을 $|X|$라 쓴다.
원소 $a \in F(X)$는 집합의 사상
$|X| \times |X| \to \mathbf{Z}/2\mathbf{Z}$, $(x, x') \mapsto a(x, x')$로
보는데, $x = x'$이거나 $x$가 $x'$의 특수화가 아니면 그 값은 영이다.
스킴의 사상 $f : X \to Y$가 주어졌을 때
$$
F(f) : F(Y) \longrightarrow F(X)
$$
를 $b \in F(Y)$에 대해 다음 규칙으로 정의한다.
$$
F(f)(b)(x, x') =
\left\{
\begin{matrix}
0 & \text{if }x\text{ is not a specialization of }x' \\
b(f(x), f(x')) & \text{else.}
\end{matrix}
\right.
$$
이것은 실제로 $F(X)$의 원소를 정의함에 주의하라. 사상 $f : X \to Y$와
$g : Y \to Z$가 합성 가능하면 $F(f) \circ F(g) = F(g \circ f)$라고
주장한다. 실제로 $c \in F(Z)$라 하고 $x' \leadsto x$를 $X$의 점들의
특수화라 하면
$$
F(g \circ f)(x, x') = c(g(f(x)), g(f(x'))) = F(g)(F(f)(c))(x, x')
$$
이다. 왜냐하면 $f(x') \leadsto f(x)$이기 때문이다.
(이는 $f(x) = f(x')$일 때도 성립한다.)

\medskip\noindent
$G$를 \'etale 위상에서 $F$를 층화한 것이라 하자.

\medskip\noindent
$X$가 스킴이고 $x' \leadsto x$가 $x' \not = x$인 특수화이면
$G(X) \not = 0$이라고 주장한다. 실제로 $a$를 함수
$|X| \times |X| \to \mathbf{Z}/2\mathbf{Z}$로 보았을 때 $(x, x')$에서
영이 아닌 원소 $a \in F(X)$를 택하자. $\{f_i : U_i \to X\}$를 $X$의
\'etale 덮개라 하자. 어떤 $i$와 $f_i(u_i) = x$인 점 $u_i \in U_i$를
택할 수 있다. 일반화는 평탄 사상을 따라 올라가므로(『사상』의 보조정리
\ref{morphisms-lemma-generalizations-lift-flat}를 보라)
$f_i(u'_i) = x'$인 특수화 $u'_i \leadsto u_i$를 찾을 수 있다. 위 구성에
따라 $F(f_i)(a) \not = 0$이다. 따라서 $a$는 $G(X)$의 영이 아닌 원소를
정의한다.

\medskip\noindent
$X = \Spec(k)$이고 $k$가 체이면(더 일반적으로 모든 소 아이디얼이 극대인
환이어도 된다) $F(X) = 0$임에 주의하라. 또 \'etale 사상 $U \to X$마다
\'etale 사상의 섬유 안에서 서로 다른 점 사이에는 특수화가 없으므로
$F(U) = 0$이다. 따라서 $G(X) = 0$이다.

\medskip\noindent
$X \subset X'$가 두꺼워짐이라고 하자. 『사상에 관하여 더 자세히』의 정의
\ref{more-morphisms-definition-thickening}를 보라. 밑변환 함자
$U' \mapsto U = U' \times_{X'} X$에 의해 $X'$ 위의 \'etale 스킴 범주와
$X$ 위의 \'etale 스킴 범주는 동치이다. 『\'Etale 코호몰로지』의 정리
\ref{etale-cohomology-theorem-topological-invariance}를 보라. 이 상황에서
항상 $F(U) = F(U')$이므로 $G(X) = G(X')$이기도 하다.

\medskip\noindent
변형으로, 스킴 $X$에 다음 사상들의 모임을 대응시키는 준층 $F_n$을
생각할 수 있다.
$a : |X|^{n + 1} \to \mathbf{Z}/2\mathbf{Z}$에서 $a(x_0, \ldots, x_n)$은
$x_n \leadsto \ldots \leadsto x_0$가 특수화들의 열이고
$x_n \not = x_{n - 1} \not = \ldots \not = x_0$일 때에만 영이 아닐 수 있다.
$G_n$을 $F_n$에 연관된 층이라 하자. 위와 정확히 같은 방법으로
$\dim(X) \geq n$이면 $G_n$이 영이 아니고, $\dim(X) < n$이면 영임을
보일 수 있다.

\begin{lemma}
\label{examples-lemma-sheaf-zero-on-low-dimension}
$\Sch_\etale$ 위에 다음 성질을 갖는 아벨 군의 층 $G$가 존재한다.
\begin{enumerate}
\item $\dim(X) < n$이면 $G(X) = 0$이다.
\item $\dim(X) \geq n$이면 $G(X)$는 영이 아니다.
\item $X \subset X'$가 두꺼워짐이면 $G(X) = G(X')$이다.
\end{enumerate}
\end{lemma}

\begin{proof}
위 논의를 보라.
\end{proof}

\begin{remark}
\label{examples-remark-specialization}
몇 가지 주석을 덧붙인다.
\begin{enumerate}
\item 준층 $F$와 $F_n$은 분리 준층이다.
\item 실제로 $F$, $F_n$은 층이 아니다.
\item $G$, $G_n$은 사실 fppf 위상에 대한 층임을 보일 수 있다.
\end{enumerate}
필요할 때 이 결과들을 증명할 것이다.
\end{remark}

\section{층과 구성가능 함수}
\label{examples-section-constructible-functions}

\noindent
이하에서는 『위상』의 정의 \ref{topologies-definition-big-etale-site}에서
구성한 큰 \'etale 사이트 $\Sch_\etale$를 고정한다. 또한 스킴은 이
사이트의 대상이라고 한다. 이 절에서 스킴 $X$의 {\it 구성가능 분할}이란
국소 유한 서로소 합 분해 $X = \coprod_{i \in I} X_i$로서 각
$X_i \subset X$가 $X$의 국소 구성가능 부분집합인 것을 말한다. 국소
유한이라는 것은 임의의 준콤팩트 열린집합 $U \subset X$에 대해
$i \in I$ 가운데 $X_i \cap U$가 공집합이 아닌 것이 유한 개뿐이라는 뜻이다.
$f : X \to Y$가 스킴의 사상이고 $Y = \coprod Y_j$가 구성가능 분할이면
$X = \coprod f^{-1}(Y_j)$는 $X$의 구성가능 분할임에 주의하라. 집합 $S$와
스킴 $X$가 주어졌을 때 {\it 구성가능 함수} $f : |X| \to S$란
$X = \coprod_{s \in S} f^{-1}(s)$가 $X$의 구성가능 분할인 사상을 말한다.
$G$가 (추상군)이고 $a, b : |X| \to G$가 구성가능 함수이면
$ab : |X| \to G, x \mapsto a(x)b(x)$도 구성가능 함수이다. 그 이유는
임의의 두 구성가능 분할을 동시에 세분하는 세 번째 분할이 있기 때문이다.

\medskip\noindent
$A$를 임의의 아벨 군이라 하자. 각 스킴 $X$에 대해 다음과 같이 정의한다.
$$
F(X) =
\frac{\{a : |X| \to A \mid a \text{는 구성가능 함수이다}\}}{\{\text{국소 상수 함수 }|X| \to A\}}
$$
$F(X)$의 원소 $a$는 단순히 국소 상수 함수를 더하는 것까지 잘 정의된
함수로 생각한다. 스킴의 사상 $f : X \to Y$와 원소 $b \in F(Y)$가 주어지면
$F(f)(b) = b \circ f$로 정의한다. 따라서 $F$는 $\Sch_\etale$ 위의
준층이다.

\medskip\noindent
$\{f_i : U_i \to X\}$가 fppf 덮개이고 $a \in F(X)$가
$F(f_i)(a) = 0$을 모든 지표에 대해 $F(U_i)$에서 만족하면, $a \circ f_i$는 각 $i$에 대해
국소 상수 함수임에 주의하라. 사상 $f_i$들이 열린 사상이므로 이는 다시
$a$가 국소 상수 함수라는 뜻이다. 따라서 $F(X)$에서 $a = 0$이다.
그러므로 $F$는 분리 준층이다(fppf 위상에서, 따라서 더욱이 \'etale
위상에서도 그렇다).

\medskip\noindent
$G$를 \'etale 위상에서 $F$를 층화한 것이라 하자. $F$는 분리되어 있고,
예를 들어 $X$가 이산 값매김환의 스펙트럼일 때 $F(X) \not = 0$이므로
$G$는 영이 아니다.

\medskip\noindent
$X = \Spec(k)$이고 $k$가 체라 하자. 그러면 $X$의 임의의 \'etale 덮개는
$k'/k$가 유한 분리가능 체 확대인 덮개
$\{\Spec(k') \to \Spec(k)\}$로 지배된다. $F(\Spec(k')) = 0$이므로
$G(X) = 0$이다.

\medskip\noindent
$X \subset X'$가 두꺼워짐이라고 하자. 『사상에 관하여 더 자세히』의 정의
\ref{more-morphisms-definition-thickening}를 보라. 밑변환 함자
$U' \mapsto U = U' \times_{X'} X$에 의해 $X'$ 위의 \'etale 스킴 범주와
$X$ 위의 \'etale 스킴 범주는 동치이다. 『\'Etale 코호몰로지』의 정리
\ref{etale-cohomology-theorem-topological-invariance}를 보라. 이 상황에서
$F(U) = F(U')$이므로 $G(X) = G(X')$이기도 하다.

\medskip\noindent
층 $G$는 극한을 보존한다. 『공간의 극한』의 정의
\ref{spaces-limits-definition-locally-finite-presentation}를 보라.
실제로 환 $R$이 환들의 유향 쌍대극한 $R = \colim_i R_i$로 쓰였다고 하자.
$X = \Spec(R)$, $X_i = \Spec(R_i)$라 놓으면 $X = \lim_i X_i$이다. 이때
$G(X) = \colim_i G(X_i)$이다. 이를 증명하려면 먼저 $\Spec(R)$의 구성가능
분할이 어떤 $\Spec(R_i)$의 구성가능 분할에서 온다는 것을 증명한다.
이로부터 $F$에 대한 결과를 얻는다. 층화에 대한 결과를 얻으려면 임의의
\'etale 환 사상 $R \to R'$가 어떤 $i$에 대한 \'etale 환 사상
$R_i \to R_i'$에서 온다는 것을 사용하라. 세부사항은 생략한다.

\begin{lemma}
\label{examples-lemma-weird-sheaf}
$\Sch_\etale$ 위에 다음 성질을 갖는 아벨 군의 층 $G$가 존재한다.
\begin{enumerate}
\item $G(\Spec(k)) = 0$이다. 여기서 $k$는 임의의 체이다.
\item $G$는 극한을 보존한다.
\item $X \subset X'$가 두꺼워짐이면 $G(X) = G(X')$이다.
\item $G$는 영이 아니다.
\end{enumerate}
\end{lemma}

\begin{proof}
위 논의를 보라.
\end{proof}



\section{lisse-\'etale 사이트는 함자적이지 않다}
\label{examples-section-lisse-etale-not-functorial}

\noindent
$X$의 {\it lisse-\'etale} 사이트 $X_{lisse,\etale}$는 $X$ 위에서
매끄러운 스킴들의 범주에 (통상적인) \'etale 덮개를 준 것이다.
『스택의 코호몰로지』의 절
\ref{stacks-cohomology-section-lisse-etale}를 보라. $f : X \to Y$를
스킴의 사상이라 하자. 연속 함자
$$
u : Y_{lisse,\etale} \longrightarrow X_{lisse,\etale},\quad
V/Y \longmapsto V \times_Y X
$$
가 있다. 따라서 수반 함자쌍
$$
u^s :
\Sh(X_{lisse,\etale})
\longrightarrow
\Sh(Y_{lisse,\etale}),
\quad
u_s :
\Sh(Y_{lisse,\etale})
\longrightarrow
\Sh(X_{lisse,\etale}),
$$
을 얻는다. 『사이트』의 절 \ref{sites-section-continuous-functors}를 보라.
일반적으로 $u$가 사이트의 사상을 정의하지 {\bf 않는다}고 주장한다.
『사이트』의 정의 \ref{sites-definition-morphism-sites}를 보라. 다시 말해
$u_s$는 일반적으로 왼쪽 완전하지 않다고 주장한다. 표현가능 준층은
lisse-\'etale 사이트 위의 층임에 주의하라. 따라서 『사이트』의
보조정리 \ref{sites-lemma-pullback-representable-sheaf}에 의해
$u_sh_V = h_{V \times_Y X}$이다. 이제 $Y$ 위에서 매끄러운 스킴
$V_1, V_2$ 사이의 두 사상
$$
\xymatrix{
V_1 \ar[rd] \ar@<1ex>[rr]^a \ar@<-1ex>[rr]_b & & V_2 \ar[ld] \\
& Y
}
$$
을 생각하자. $u_s$가 왼쪽 완전하다면
$$
u_s \text{Equalizer}(h_a, h_b : h_{V_1} \to h_{V_2})
=
\text{Equalizer}(h_{a \times 1}, h_{b \times 1} :
h_{V_1 \times_Y X} \to h_{V_2 \times_Y X})
$$
를 얻을 것이다. $Y$ 위에서 매끄러운 스킴에서
$(a = b) \subset V_1$로 가는 사상이 존재하지 않도록, 즉 왼쪽 항이
공집합 층이 되도록 하되 $X$로 밑변환한 뒤에는 등화자가 공집합이 아니고
$X$ 위에서 매끄럽도록 사상 $a, b : V_1 \to V_2$를 택할 것이다.
간단한 예로 $X = \Spec(\mathbf{F}_p)$, $Y = \Spec(\mathbf{Z})$,
$V_1 = V_2 = \mathbf{A}^1_\mathbf{Z}$로 택하고 $a(x) = x$,
$b(x) = x + p$로 정하면 된다. $a$와 $b$의 등화자는
$\mathbf{A}^1_\mathbf{Z}$의 $(p)$ 위의 섬유임에 주의하라.

\begin{lemma}
\label{examples-lemma-lisse-etale-not-functorial}
lisse-\'etale 사이트는 스킴의 사상에 대해서조차 함자적이지 않다.
\end{lemma}

\begin{proof}
위 논의를 보라.
\end{proof}





\section{뇌터 스킴의 범주 위의 층}
\label{examples-section-sheaves-locally-Noetherian}

\noindent
$S$를 국소 뇌터 스킴이라 하자. 『Artin의 공리』의 절
\ref{artin-section-noetherian}에서처럼 포함 함자
$$
u : (\textit{Noetherian}/S)_{fppf} \longrightarrow (\Sch/S)_{fppf}
$$
를 생각하자. 이는 $S$ 위의 국소 뇌터 스킴들의 fppf 사이트를 $S$의
큰 fppf 사이트에 포함시킨다. 인용한 절에서 설명했듯 이 함자는
연속이다. 따라서 수반 함자쌍
$$
u^s :
\Sh((\Sch/S)_{fppf})
\longrightarrow
\Sh((\textit{Noetherian}/S)_{fppf})
$$
과
$$
u_s :
\Sh((\textit{Noetherian}/S)_{fppf})
\longrightarrow
\Sh((\Sch/S)_{fppf})
$$
을 얻는다. 『사이트』의 절 \ref{sites-section-continuous-functors}를
보라. 그러나 일반적으로 $u$가 사이트의 사상을 정의하지 않는다고
주장한다. 『사이트』의 정의 \ref{sites-definition-morphism-sites}를
보라. 다시 말해 함자 $u_s$는 일반적으로 왼쪽 완전하지 않다고 주장한다.

\medskip\noindent
$p$를 소수라 하고 $S = \Spec(\mathbf{F}_p)$라 놓자.
$(\textit{Noetherian}/S)_{fppf}$ 위의 층들의 단사 사상
$$
a : \mathcal{F} \longrightarrow \mathcal{G}
$$
을 다음과 같이 정의한다. $U$가 $S$ 위의 국소 뇌터 스킴이면
$$
\mathcal{G}(U) = \Gamma(U, \mathcal{O}_U)^* =
\Mor_S(U, \mathbf{G}_{m, S})
$$
로 정의하고
$$
\mathcal{F}(U) = \{f \in \mathcal{G}(U) \mid \text{fppf 국소적으로 }f
\text{는 임의 차수의 }p\text{-거듭제곱근을 갖는다}\}
$$
로 둔다. 뇌터 $\mathbf{F}_p$-대수 $A$의 멱영근기 $I \subset A$는
멱영이고, $A$ 안에서 $1$의 $p$-거듭제곱근은 $1 + a$, $a \in I$ 꼴의
원소이다. 따라서 $1$의 자명하지 않은 어떠한 $p$-거듭제곱근도 임의
차수의 $p$-거듭제곱근을 갖지 않는다. 그러므로 국소 뇌터 스킴 $U$마다
$\mathcal{F}(U)$는 $p$-꼬임이 없는 아벨 군이다. 몇 가지 세부사항은
생략한다. 따라서 $p : \mathcal{F} \to \mathcal{F}$는
$(\textit{Noetherian}/S)_{fppf}$ 위의 아벨 층들의 단사 사상이다.

\medskip\noindent
모순을 얻기 위해 $u_s$가 완전하다고 가정하자. 그러면
$p : u_s\mathcal{F} \to u_s\mathcal{F}$도 단사이고, $S$ 위의 임의의
$V$에 대해 $(u_s\mathcal{F})(V)$는 $p$-꼬임이 없는 아벨 군이다.
표현가능 준층은 fppf 사이트 위의 층이므로 『사이트』의 보조정리
\ref{sites-lemma-pullback-representable-sheaf}에 의해
$u_s\mathcal{G}$는 $\mathbf{G}_{m, S}$로 표현된다.
$u_s\mathcal{F} \to u_s\mathcal{G}$가 단사라는 사실을 사용하면,
$S$ 위의 모든 스킴 $V$에 대해 $p$-꼬임이 없는 부분군
$$
(u_s\mathcal{F})(V) \subset \Gamma(V, \mathcal{O}_V)^*
$$
를 얻으며 이는 다음 성질을 갖는다. $S$ 위 스킴의 사상 $V \to U$에서
$U$가 국소 뇌터이면 부분군
$$
\mathcal{F}(U) \subset \Gamma(U, \mathcal{O}_U)^*
$$
는 제한 사상
$\Gamma(U, \mathcal{O}_U)^* \to \Gamma(V, \mathcal{O}_V)^*$에 의해
부분군 $(u_s\mathcal{F})(V)$ 안으로 간다.

\medskip\noindent
이제 실제 모순은 다음과 같이 얻는다.
$k = \bigcup_{n \geq 0} \mathbf{F}_p(t^{1/{p^n}})$라 하고
$$
B = k \otimes_{\mathbf{F}_p(t)} k
$$
및 $V = \Spec(B)$라 놓자. 두 쌍대곱 사상 $k \to B$에 대응하는 두 사영
$V \to \Spec(k)$가 있고 $\Spec(k)$는 뇌터이므로 부분군
$$
(u_s\mathcal{F})(V) \subset B^*
$$
는 $k^* \otimes 1$과 $1 \otimes k^*$를 포함한다. 이는 모순이다. 실제로
$$
(t^{1/p} \otimes 1) \cdot (1 \otimes t^{-1/p}) =
t^{1/p} \otimes t^{-1/p}
$$
는 $B$의 자명하지 않은 $p$-꼬임 가역원이다.

\begin{lemma}
\label{examples-lemma-not-a-morphism-of-sites-noetherian-to-all}
$S = \Spec(\mathbf{F}_p)$일 때 포함 함자
$(\textit{Noetherian}/S)_{fppf} \to (\Sch/S)_{fppf}$는 사이트의 사상을
정의하지 않는다.
\end{lemma}

\begin{proof}
위 논의를 보라.
\end{proof}



\section{준연접 가군의 유도 밂앞}
\label{examples-section-derived-push-quasi-coherent}

\noindent
$k$를 표수 $p > 0$인 체라 하자. $S = \Spec(k[x])$라 하자.
$G = \mathbf{Z}/p\mathbf{Z}$를 추상군 또는 $S$ 위의 상수 군 스킴으로 보자.
$G$가 $S$에 자명하게 작용할 때 대수적 스택
$\mathcal{X} = [S/G]$를 생각하자. 『스택의 예』의 주석
\ref{examples-stacks-remark-X-mod-G-group}과 『표현가능성의 판정기준』의
보조정리 \ref{criteria-lemma-BG-algebraic}를 보라. 구조 사상
$$
f : \mathcal{X} \longrightarrow S
$$
을 생각하자. 이 사상은 준콤팩트하고 준분리이다. 따라서 함자
$$
Rf_{\QCoh, *} :
D^{+}_\QCoh(\mathcal{O}_\mathcal{X})
\longrightarrow
D^{+}_\QCoh(\mathcal{O}_S),
$$
를 얻는다. 『스택의 유도범주』의 명제
\ref{stacks-perfect-proposition-derived-direct-image-quasi-coherent}를 보라.
$Rf_{\QCoh, *}\mathcal{O}_\mathcal{X}$를 계산하자.
$D_\QCoh(\mathcal{O}_S)$는 $k[x]$-가군의 유도범주와 동치이므로
(『스킴의 유도범주』의 보조정리
\ref{perfect-lemma-affine-compare-bounded}를 보라), 이는
$R\Gamma(\mathcal{X}, \mathcal{O}_\mathcal{X})$를 계산하는 것과 같다.
이를 위해 덮개 $S \to \mathcal{X}$와 스펙트럼 열
$$
H^q(S \times_\mathcal{X} \ldots \times_\mathcal{X} S, O)
\Rightarrow H^{p + q}(\mathcal{X}, \mathcal{O}_\mathcal{X})
$$
을 쓸 수 있다. 『스택의 코호몰로지』의 명제
\ref{stacks-cohomology-proposition-smooth-covering-compute-cohomology}를 보라.
다음 등식에 주의하라.
$$
S \times_\mathcal{X} \ldots \times_\mathcal{X} S = S \times G^p
$$
이는 아핀이다. 따라서 복합체
$$
k[x] \to \text{Map}(G, k[x]) \to \text{Map}(G^2, k[x]) \to \ldots
$$
가 $R\Gamma(\mathcal{X}, \mathcal{O}_\mathcal{X})$를 계산한다. 여기서
$\varphi \in \text{Map}(G^{p - 1}, k[x])$의 미분은
$(g_1, \ldots, g_p)$를 다음 원소로 보내는 사상이다.
$$
\varphi(g_2, \ldots, g_p) +
\sum\nolimits_{i = 1}^{p - 1}
(-1)^i\varphi(g_1, \ldots, g_i + g_{i + 1}, \ldots, g_p)
+ (-1)^p\varphi(g_1, \ldots, g_{p - 1}).
$$
이는 $k[x]$에 자명하게 작용하는 $G$의 군 코호몰로지를 계산하는
복합체일 뿐이다(나중에 참조를 삽입하라). 순환군 $G$가 $k[x]$에
자명하게 작용할 때의 코호몰로지는 각 코호몰로지 차수 $\geq 0$에서 정확히 $k[x]$ 한 개이다
(나중에 참조를 삽입하라). 따라서
$$
Rf_*\mathcal{O}_\mathcal{X} = \bigoplus\nolimits_{n \geq 0} \mathcal{O}_S[-n]
$$
이다. 이제 복합체
$$
E = \bigoplus\nolimits_{m \geq 0} \mathcal{O}_\mathcal{X}[m]
$$
을 생각하자. 이는 $D_\QCoh(\mathcal{O}_\mathcal{X})$의 대상이다.
일반적인 결과를 위해 잠시 논의를 중단한다.

\begin{lemma}
\label{examples-lemma-is-limit}
$\mathcal{X}$를 대수적 스택이라 하자. $K$를
$D(\mathcal{O}_\mathcal{X})$의 대상으로서 그 코호몰로지 층들이 국소
준연접이고(『스택 위의 층』의 정의
\ref{stacks-sheaves-definition-locally-quasi-coherent}) 평탄 밑변환 성질을
만족한다고 하자(『스택의 코호몰로지』의 정의
\ref{stacks-cohomology-definition-flat-base-change}). 그러면
$D(\mathcal{O}_\mathcal{X})$ 안에 구별삼각형
$$
K \to
\prod\nolimits_{n \geq 0} \tau_{\geq -n} K \to
\prod\nolimits_{n \geq 0} \tau_{\geq -n} K \to K[1]
$$
이 존재한다. 다시 말해 $K$는 그 표준 절단들의 유도극한이다.
\end{lemma}

\begin{proof}
$\mathcal{X}$의 ``큰 fppf 사이트'' $\mathcal{X}_{fppf}$에서 작업한다는
것을 상기하라(『스택 위의 층』과 『스택 위의 코호몰로지』 장에서
$\mathcal{O}_\mathcal{X}$-가군의 층에 관해 채택한 관례에 따른다).
$\mathcal{B}$를 아핀 스킴 $U$ 위에 놓이는 $\mathcal{X}_{fppf}$의 대상
$x$들의 집합이라 하자. 『스택 위의 층』의 보조정리들
\ref{stacks-sheaves-lemma-compare-fppf-etale},
\ref{stacks-sheaves-lemma-cohomology-restriction}, 『강하』의 보조정리
\ref{descent-lemma-quasi-coherent-and-flat-base-change}, 그리고 『스킴의
코호몰로지』의 보조정리
\ref{coherent-lemma-quasi-coherent-affine-cohomology-zero}를 종합하면,
$\mathcal{F}$가 국소 준연접이고 $x \in \mathcal{B}$이면
$H^p(x, \mathcal{F}) = 0$임을 알 수 있다. 이제 $d = 0$으로 둔
『사이트 위의 코호몰로지』의 보조정리
\ref{sites-cohomology-lemma-is-limit-dimension}에서 주장이 따른다.
\end{proof}

\begin{lemma}
\label{examples-lemma-sum-is-product}
$\mathcal{X}$를 대수적 스택이라 하자. $\mathcal{F}_n$이 $\mathcal{X}$ 위의
국소 준연접 층들로 이루어진 모임이고 평탄 밑변환 성질을 만족하면,
$D(\mathcal{O}_\mathcal{X})$에서
$\oplus_n \mathcal{F}_n[n] \to \prod_n \mathcal{F}_n[n]$은 동형사상이다.
\end{lemma}

\begin{proof}
보조정리 \ref{examples-lemma-is-limit}에 의해 직합이 곱과 동형이므로 성립한다.
\end{proof}

\noindent
논의를 계속하자. 준연접 가군은 국소 준연접이고 평탄 밑변환 성질을
만족하므로(『스택 위의 층』의 보조정리
\ref{stacks-sheaves-lemma-quasi-coherent}) 다음을 얻는다.
$$
E = \prod\nolimits_{m \geq 0} \mathcal{O}_\mathcal{X}[m]
$$
코호몰로지는 극한과 가환하므로
$$
Rf_*E = \prod\nolimits_{m \geq 0}
\left(\bigoplus\nolimits_{n \geq 0} \mathcal{O}_S[m - n]\right)
$$
이다. 이 복합체는 $D_\QCoh(\mathcal{O}_S)$의 대상이 아님에 주의하라.
실제로 차수 $0$의 코호몰로지 층은 $\mathcal{O}_S$의 복사본들의 무한
곱이며, 이는 국소 준연접 $\mathcal{O}_S$-가군조차 아니다.

\begin{lemma}
\label{examples-lemma-push-not-OK}
대수적 스택의 준콤팩트 준분리 사상
$f : \mathcal{X} \to \mathcal{Y}$는 함자
$Rf_* : D_\QCoh(\mathcal{O}_\mathcal{X}) \to
D_\QCoh(\mathcal{O}_\mathcal{Y})$를 항상 유도하는 것은 아니다.
\end{lemma}

\begin{proof}
위 논의를 보라.
\end{proof}

\section{큰 아벨 범주}
\label{examples-section-big}

\noindent
이 절의 목적은 ``큰'' 아벨 범주 $\mathcal{A}$와 대상 $M, N$을 주어
확장들의 동형류 모임 $\Ext_\mathcal{A}(M, N)$이 집합이 아닌 예를 보이는
것이다. 이 예는 Freyd가 제시했다. \cite[page 131, Exercise A]{Freyd}를
보라.

\medskip\noindent
$\mathcal{A}$를 다음과 같이 정의한다. $\mathcal{A}$의 대상은 삼중항
$(M, \alpha, f)$로 이루어진다. 여기서 $M$은 아벨 군, $\alpha$는 서수,
$f : \alpha \to \text{End}(M)$은 사상이다. 사상
$(M, \alpha, f) \to (M', \alpha', f')$은 아벨 군의 준동형사상
$\varphi : M \to M'$로서 {\it 모든} 서수 $\beta$에 대해
$$
\varphi \circ f(\beta) = f'(\beta) \circ \varphi
$$
를 만족하는 것이다. 여기서 $\beta$가 $\alpha$에 속하지 않으면
$f(\beta) = 0$으로 정하고, 마찬가지로 $\beta$가 $\alpha'$의 원소가
아니면 $f'(\beta)$를 영으로 정한다. 이렇게 정의한 범주가 아벨
범주라는 확인은 생략한다.

\medskip\noindent
대상 $Z = (\mathbf{Z}, \emptyset, f)$, 즉 모든 작용소가 영인 대상을
생각하자. 핵심 관찰은 $\mathcal{A}$에서 계산한 군
$\Ext^1_\mathcal{A}(Z, Z)$가 집합이 아닌 진클래스라는 것이다.
실제로 각 서수 $\alpha$에 대해 $Z$에 의한 $Z$의 확장
$(M, \alpha + 1, f)$을 찾을 수 있는데, 그 밑군은
$M = \mathbf{Z} \oplus \mathbf{Z}$이고 $f$의 값은 다음 값을 제외하면 항상
영이다.
$$
f(\alpha) =
\left(
\begin{matrix}
0 & 1 \\
0 & 0
\end{matrix}
\right).
$$
이는 확장들의 동형류로 이루어진 진클래스를 분명히 만든다. 특히
$\mathcal{A}$의 유도범주에서 사상들의 모임은 진클래스이다.
『유도범주』의 보조정리 \ref{derived-lemma-ext-1}을 보라. 이는 삼각범주의
Verdier 몫을 정의할 때 어느 정도 주의해야 함을 뜻한다.

\begin{lemma}
\label{examples-lemma-big-abelian-category}
$\Ext$-군들이 진클래스인 ``큰'' 아벨 범주 $\mathcal{A}$가 존재한다.
\end{lemma}

\begin{proof}
위 논의를 보라.
\end{proof}



\section{약하게 연관된 점과 스킴론적 조밀성}
\label{examples-section-weak-ass-dense}

\noindent
$k$를 체라 하자. $i$가 $\mathbf{N}$의 원소들을 달릴 때
$R = k[z, x_i, y_i]/(z^2, zx_iy_i)$라 하자. 여기서
$R = R_0 \oplus M_0$이고, $R_0 = k[x_i, y_i]$는 부분환이며 $M_0$는 제곱이
영인 아이디얼이고 $R_0$-가군으로서 $M_0 \cong R_0/(x_iy_i)$임에
주의하라. 소 아이디얼 $\mathfrak p = (z, x_i)$는 $R$을 $R$-가군으로
보았을 때 약하게 연관되어 있다(『대수학』의 정의
\ref{algebra-definition-weakly-associated}). 실제로 $R_\mathfrak p$의
원소 $z$는 영이 아니지만 $\mathfrak pR_\mathfrak p$에 의해 소멸된다.
한편 열린 부분스킴
$$
U = \bigcup D(x_i) \subset \Spec(R) = S
$$
을 생각하자. $U \subset S$가 스킴론적으로 조밀하다고 주장한다
(『사상』의 정의 \ref{morphisms-definition-scheme-theoretically-dense}).
이를 증명하려면 $j : U \to S$가 포함사상일 때
$\mathcal{O}_S \to j_*\mathcal{O}_U$가 단사임을 보이면 충분하다.
『사상』의 보조정리
\ref{morphisms-lemma-characterize-scheme-theoretically-dense}를 보라.
이를 대수학으로 옮기면 모든 $g \in R$에 대해 사상
$$
R_g \longrightarrow \prod R_{x_ig}
$$
가 단사임을 보여야 한다. $g = g_0 + m_0$로 쓰자. 여기서
$g_0 \in R_0$, $m_0 \in M_0$이다. 그러면 $R_g = R_{g_0}$이다
(세부사항은 생략한다). 따라서 $g \in R_0$라고 가정할 수 있다.
또 $g$가 영이 아니라고 가정할 수 있다. 이제
$R_g = (R_0)_g \oplus (M_0)_g$이다. $R_0$는 정역이므로
$(R_0)_g \to \prod (R_0)_{x_ig}$는 단사이다. $g \in (x_iy_i)$이면
$(M_0)_g = 0$이므로 보일 것이 없다. $g \not \in (x_iy_i)$이면
$(x_iy_i)$가 $R_0$의 근기 아이디얼이므로
$M_0 \to \prod (M_0)_{x_ig}$가 단사임을 보이면 된다.
$R_0 \to M_0 \to (M_0)_{x_n}$의 핵은 $(x_iy_i, y_n)$이다.
$(x_iy_i, y_n)$은 근기 아이디얼이므로
$g \not \in (x_iy_i, y_n)$이면 $R_0 \to M_0 \to (M_0)_{x_ng}$의 핵은
$(x_iy_i, y_n)$이다. 모든 $n \gg 0$에 대해
$g \not \in (x_iy_i, y_n)$이므로 그 핵은 원하는 대로
$\bigcap_{n \gg 0} (x_iy_i, y_n) = (x_iy_i)$에 포함된다.

\medskip\noindent
두 번째 예는 Ofer Gabber가 제시했다. $k$를 체라 하고 $R$과 $R'$를
각각 함수 $\mathbf{N} \to k$의 환과 결국 상수인 함수 $\mathbf{N} \to k$의
환이라 하자. 그러면 $\Spec(R)$과 $\Spec(R')$은 각각
Stone-{\v C}ech 콤팩트화\footnote{모든 $f \in R$는 $u$가 가역원이고
$e$가 멱등원인 $ue$ 꼴이다. 그러면 『대수학』의 보조정리
\ref{algebra-lemma-ring-with-only-minimal-primes}는 $\Spec(R)$가
하우스도르프임을 보여 준다. 한편 이산 위상을 준 $\mathbf{N}$은 조밀한
열린 부분집합으로 볼 수 있다. 집합 사상 $\mathbf{N} \to X$가 하우스도르프
준콤팩트 위상 공간 $X$로 갈 때 환 사상 $\mathcal{C}^0(X; k) \to R$을
얻는다. 여기서 $\mathcal{C}^0(X; k)$는 국소 상수 사상 $X \to k$의
$k$-대수이다. 이로부터 $\Spec(R) \to \Spec(\mathcal{C}^0(X; k)) = X$를
얻어 보편 성질이 증명된다.} $\beta\mathbf{N}$과 한 점
콤팩트화\footnote{여기서는 $R'$에 추가되는 극대 아이디얼이 실제로
하나뿐이라고 논증한다.} $\mathbf{N}^* = \mathbf{N} \cup \{\infty\}$이다.
환 $R$과 $R'$의 모든 소 아이디얼이 극소이므로 모든 점은 약하게
연관되어 있다.

\begin{lemma}
\label{examples-lemma-example-schematically-dense-missing-weakly-associated-point}
축약 스킴 $X$와 스킴론적으로 조밀한 열린집합 $U \subset X$가 존재하여,
어떤 약하게 연관된 점 $x \in X$가 $U$에 속하지 않는다.
\end{lemma}

\begin{proof}
첫 번째 예에서는 구성상 $\mathfrak p \not \in U$이다. Gabber의 예에서
스킴 $\Spec(R)$ 또는 $\Spec(R')$는 축약이다.
\end{proof}

\section{대각합의 비가법성의 예}
\label{examples-section-non-additive}

\noindent
$k$를 체라 하고 $R = k[\epsilon]$을 $k$ 위의 쌍대수 환이라 하자.
다시 말해 $R = k[x]/(x^2)$이고 $\epsilon$은 $R$에서 $x$의 합동류이다.
복합체들의 짧은 완전열
$$
\xymatrix{
0 \ar[d] \ar[r] &
R \ar[d]^\epsilon \ar[r]_1 &
R \ar[d] \\
R \ar[r]^1 & R \ar[r] & 0
}
$$
을 생각하자. 여기서 열들이 복합체이고 첫째 행은 차수 $0$, 둘째 행은
차수 $1$에 놓인다. 첫째 복합체(즉 왼쪽 열)를 $A^\bullet$, 둘째를
$B^\bullet$, 셋째를 $C^\bullet$이라 쓰자. 그림
\begin{equation}
\label{examples-equation-commutes-up-to-homotopy}
\vcenter{
\xymatrix{
A^\bullet \ar[d]_{1 + \epsilon} \ar[r] &
B^\bullet \ar[r] \ar[d]_1 &
C^\bullet \ar[d]_1 \\
A^\bullet \ar[r] & B^\bullet \ar[r] & C^\bullet
}
}
\end{equation}
은 $K(R)$에서 가환한다고, 즉 호모토피까지 가환하는 복합체들의
그림이라고 주장한다. 실제로 오른쪽 사각형은 가환하고, 왼쪽 사각형의
어긋남은 호모토피 $1 : A^1 \to B^0$이다. 한편
$$
\text{Tr}_{A^\bullet}(1 + \epsilon) + \text{Tr}_{C^\bullet}(1)
\not = \text{Tr}_{B^\bullet}(1).
$$

\begin{lemma}
\label{examples-lemma-nonadditivity-of-trace}
환 $R$, 호모토피 범주 $K(R)$ 안의 구별삼각형
$(K, L, M, \alpha, \beta, \gamma)$, 그리고 이 구별삼각형의 자기준동형
$(a, b, c)$가 존재하여, $K$, $L$, $M$은 완전 복합체이고
$\text{Tr}_K(a) + \text{Tr}_M(c) \not = \text{Tr}_L(b)$이다.
\end{lemma}

\begin{proof}
위의 예를 생각하자. 사상 $\gamma : C^\bullet \to A^\bullet[1]$은 차수
$0$에서 $\epsilon$을 곱하는 사상으로 주어진다. 『유도범주』의 정의
\ref{derived-definition-distinguished-triangle}를 보라. 따라서
$\epsilon(1 + \epsilon) = \epsilon$이므로
$$
\xymatrix{
C^\bullet \ar[d] \ar[r]_\gamma & A^\bullet[1] \ar[d] \\
C^\bullet \ar[r]^\gamma & A^\bullet[1]
}
$$
도 $K(R)$에서 가환한다. 그러므로 실제로 구별삼각형들의 사상을 얻었다.
\end{proof}




\section{사영 스킴 구조층의 유한하지 않은 코호몰로지}
\label{examples-section-nonfinite-h0}

\noindent
$R$을 임의의 비뇌터 환이라 하고 $I$를 유한 생성이 아닌 $R$의
아이디얼이라 하자. $S = R[x, y]/xyI$를 $x$와 $y$의 차수가 $1$인 등급환으로
보자. 그러면 $X = \text{Proj}(S)$는 $R$ 위의 사영 스킴이다.
$H^0(X, \mathcal{O}_X)$는 유한 생성 $R$-가군이 아니라고 주장한다.

\medskip\noindent
환 $S_x$, $S_y$, $S_{xy}$를 이중차수에 따라 분해하면
$$
S_x = \bigoplus_{\substack{(a, b) \in \mathbf{Z}^2 \\ b \geq 0}} R'_{a, b},
\quad\quad
S_y = \bigoplus_{\substack{(a, b) \in \mathbf{Z}^2 \\ a \geq 0}} R''_{a, b},
\quad\quad
S_{xy} = \bigoplus_{\substack{(a, b) \in \mathbf{Z}^2}} R/I,
$$
로 쓸 수 있다. 여기서 $(a, b)$-성분은 차수 $a + b$에 놓이고
$$
R'_{a, b} =
\left\{
\begin{matrix}
R & b = 0, \\
R/I & b \ne 0,
\end{matrix}
\right.
\quad\quad
R''_{a, b} =
\left\{
\begin{matrix}
R & a = 0, \\
R/I & a \ne 0.
\end{matrix}
\right.
$$
이다. $R$-가군 $H^0(X, \mathcal{O}_X)$는
$$
(S_x)_0 \times (S_y)_0 \to (S_{xy})_0
$$
에서 $(F, G) \mapsto F - G$로 정의된 사상의 핵으로 계산할 수 있다.
여기서 $A_0$라는 표기는 $\mathbf{Z}$-등급환 $A$의 차수 $0$ 성분을
뜻한다. 위에서 $S_x, S_y, S_{xy}$를 구체적으로 설명했으므로 이 핵은
$F - G \in I$인 쌍 $(F, G)$들로 이루어진 부분가군 $M \subset R^2$와
동형이다. $(F, G) \mapsto F - G$로 주어진 $R$-가군의 사상 $M \to I$가
전사이므로 $M$은 유한 생성 $R$-가군이 아니다.

\medskip\noindent
전사 $R[x, y] \to S$는 닫힌 몰입 $i : X \to \mathbf{P}^1_R$에 대응한다.
그러면 $\mathcal{F} = i_*\mathcal{O}_X$는 유한형 준연접
$\mathcal{O}_{\mathbf{P}^1_R}$-가군이고, 그 $H^0$은 위에서 계산한 유한
생성이 아닌 가군 $M$과 같다.

\begin{lemma}
\label{examples-lemma-nonfinite-h0}
유한하지 않은 $H^0$.
\begin{enumerate}
\item 환 $R$과 $R$ 위의 사영 스킴 $X$가 존재하여
$H^0(X, \mathcal{O}_X)$는 유한 $R$-가군이 아니다.
\item 환 $R$과 $\mathbf{P}^1_R$ 위의 유한형 준연접 가군 $\mathcal{F}$가
존재하여 $H^0(\mathbf{P}^1_R, \mathcal{F})$는 유한 $R$-가군이 아니다.
\end{enumerate}
\end{lemma}

\begin{proof}
위 논의를 보라.
\end{proof}








\section{사영성은 밑 위에서 국소적인 성질이 아니다}
\label{examples-section-non-descending-property-projective}

\noindent
강하 장에서 사상의 많은 성질이 밑 위에서, 심지어 fpqc 위상에서도
국소적임을 보았다. 『강하』의 절들
\ref{descent-section-descending-properties-morphisms},
\ref{descent-section-descending-properties-morphisms-fpqc}, 그리고
\ref{descent-section-descending-properties-morphisms-fppf}를 보라.
사상의 사영성에 대해서는 그렇지 않다.

\begin{lemma}
\label{examples-lemma-non-descending-property-projective}
성질들
\begin{enumerate}
\item[] $\mathcal{P}(f) =$ ``$f$는 사영이다'', 그리고
\item[] $\mathcal{P}(f) =$ ``$f$는 준사영이다''
\end{enumerate}
는 밑 위에서 자리스키 국소적이지 않다. 따라서 더욱이 밑 위에서 fpqc
국소적이지 않다.
\end{lemma}

\begin{proof}
Hironaka \cite[Example B.3.4.1]{H}를 따라, 자리스키 국소적으로는
사영이지만 사영은 아닌 매끄러운 복소 3차원 다양체 사이의 고유 사상
$f:V_Y\to Y$를 정의한다. $f$는 고유이지만 사영이 아니므로 준사영도
아니다.

\medskip\noindent
$Y$를 복소수 ${\mathbf C}$ 위의 사영 3차원 공간이라 하자. $C$와 $D$를
$Y$ 안의 매끄러운 이차곡선으로 택하되, 닫힌 부분스킴 $C\cap D$가
축약이고 두 복소점 $P$와 $Q$로 이루어지게 하자. (예를 들어
$C=\{ [x,y,z,w]: xy=z^2, w=0\}$,
$D=\{ [x,y,z,w]: xy=w^2, z=0\}$,
$P=[1,0,0,0]$, $Q=[0,1,0,0]$로 택하라.)
$Y-Q$ 위에서는 먼저 곡선 $C$를 부풀리고, 이어 곡선 $D$의 엄밀
변환을 부풀린다(『인자』의 정의
\ref{divisors-definition-strict-transform}). $Y-P$ 위에서는 먼저 곡선
$D$를 부풀리고 이어 곡선 $C$의 엄밀 변환을 부풀린다. $Y-P-Q$ 위에서
이렇게 구성한 두 다양체는 표준적으로 동형이므로 $Y-P-Q$ 위에서 이들을
붙일 수 있다. 그 결과는 ${\mathbf C}$ 위의 매끄럽고 고유한 3차원
다양체 $V_Y$이다. 사상 $f:V_Y\to Y$는 고유이고 자리스키 국소적으로
사영이다($Y-P$ 위와 $Y-Q$ 위에서 부풀리기이기 때문이다). 『인자』의
보조정리 \ref{divisors-lemma-blowing-up-projective}를 보라.
$V_Y$가 ${\mathbf C}$ 위에서 사영이 아님을 보일 것이다. 그러면 $f$가
사영이 아님이 따른다.

\medskip\noindent
이를 위해 $L$을 $C-P-Q$의 한 복소점의 $V_Y$ 안의 역상이라 하고,
$M$을 $D-P-Q$의 한 복소점의 역상이라 하자. 그러면 $L$과 $M$은 사영직선
${\mathbf P}^1_{{\mathbf C}}$와 동형이다.
다음으로 $E$를 $V_Y$ 안에서 $C\cup D\subset Y$의 $V_Y$ 안의 역상이라 하자. 따라서
$E\rightarrow C\cup D$는 고유 사상이고, $(C\cup D)-\{P,Q\}$ 위의
섬유들은 ${\mathbf P}^1$과 동형이다. $E$에서 $P$의 역상은 두 직선
$L_0$와 $M_0$의 합집합이고, $E$ 위에서 순환의 유리동치
$L\sim L_0+M_0$와 $M\sim M_0$가 있다($C$와 $D$가 ${\mathbf P}^1$과
동형임을 쓴다). 두 곡선을 부푼 순서에서 생기는 비대칭성에 주의하라.
$Q$ 근방에서는 반대 현상이 일어난다. 따라서 $Q$의 역상은 두 직선
$L_0'$와 $M_0'$의 합집합이고 $E$ 위에서 유리동치
$L\sim L_0'$와 $M\sim L_0'+M_0'$가 있다. 이 동치들을 결합하면
$E$ 위에서, 따라서 $V_Y$ 위에서 $L_0+M_0'\sim 0$을 얻는다.
$V_Y$가 ${\mathbf C}$ 위에서 사영이라면 풍부한 선다발 $H$를 가져서
$V_Y$의 모든 곡선 위에서 차수가 $> 0$일 것이다. 특히 $H$는
$L_0+M_0'$ 위에서 양의 차수를 가져야 하는데, 이는 체 위의 고유 스킴에서
선다발의 차수가 유리동치를 법으로 한 1-순환 위에서 잘 정의된다는 사실과
모순이다(『Chow 호몰로지』의 보조정리
\ref{chow-lemma-proper-pushforward-rational-equivalence}
및 보조정리 \ref{chow-lemma-factors}).
따라서 $V_Y$는 ${\mathbf C}$ 위에서 사영이 아니다.
\end{proof}

\noindent
다른 용어로 말하면 Hironaka의 3차원 다양체 $V_Y$는 축약 부분스킴
$C\cup D$를 따라 $Y$를 부푼 $Y'$의 작은 해소이다. 여기서 $Y'$에는
두 개의 마디 특이점이 있다. $C$를 따라 $Y$를 부풀린 뒤 $D$의 엄밀
변환을 따라 부풀려 $Z$를 정의하면 $Z$는 매끄러운 사영 3차원
다양체이고, 사영이 아닌 3차원 다양체 $V_Y$는 $Y-P$ 위에서의
``플롭''에 의해 $Z$와 달라진다.

\section{사영 스킴의 효과적이지 않은 강하 데이터}
\label{examples-section-non-effective-descent-projective}

\noindent
강하 장에서 fpqc 사상에 대한 스킴의 강하 데이터가 여러 사상 부류에
대해 효과적임을 보았다. 특히 아핀 사상, 더 일반적으로 준아핀 사상은
fpqc 덮개에 대한 강하를 만족한다(『강하』의 보조정리
\ref{descent-lemma-quasi-affine}). 사영 사상에 대해서는 그렇지 않다.

\begin{lemma}
\label{examples-lemma-non-effective-descent-projective}
스킴들의 에탈 덮개 $X\to S$와 $X\to S$에 대한 강하 데이터
$(V/X,\varphi)$가 존재하여 $V\to X$는 사영이지만 이 강하 데이터는
스킴의 범주에서 효과적이지 않다.
\end{lemma}

\begin{proof}
스킴이 아닌 매끄러운 분리 복소 대수공간의 Hironaka 예
\cite[Example B.3.4.2]{H}를 모방한다.

\medskip\noindent
군 $G = \mathbf{Z}/2 = \{1, g\}$가 복소수 위의 사영 3차원 공간
$\mathbf{P}^3$에 다음과
같이 작용한다고 하자.
$$
g[x,y,z,w] = [y,x,w,z].
$$
이 작용은 ${\mathbf P}^3$ 안의 서로소인 두 직선
$L_1=\{ [x,x,z,z]\}$와 $L_2=\{ [x,-x,z,-z]\}$ 바깥에서 자유롭다.
$Y={\mathbf P}^3-(L_1\cup L_2)$라 하자. ${\mathbf C}$ 위에 매끄러운
준사영 스킴 $S=Y/G$가 있어서 $Y\to S$는 $G$-토서이다
(『군체』의 정의 \ref{groupoids-definition-principal-homogeneous-space}).
구체적으로 $S$는 ${\mathbf P}^3$의 열린 부분집합 $Y$를 사상
\begin{align*}
{\mathbf P}^3 & \to \text{Proj } {\mathbf C}[x,y,z,w]^G\\
   & = \text{Proj } {\mathbf C}[u_0,u_1,v_0,v_1,v_2]/(v_0v_1=v_2^2),
\end{align*}
으로 보낸 상으로 정의할 수 있다. 여기서
$u_0=x+y$, $u_1=z+w$, $v_0=(x-y)^2$, $v_1=(z-w)^2$,
$v_2=(x-y)(z-w)$이고, 환에는 $u_0,u_1$의 차수가 1,
$v_0,v_1,v_2$의 차수가 2인 등급을 준다.

\medskip\noindent
$C=\{ [x,y,z,w]: xy=z^2, w=0\}$,
$D=\{ [x,y,z,w]: xy=w^2, z=0\}$라 하자. 이들은 ${\mathbf P}^3$ 안의
매끄러운 이차곡선이고 $G$-불변 열린 부분집합 $Y$에 포함되며
$g(C)=D$이다. 또한 $C\cap D$는 두 점 $P:=[1,0,0,0]$와
$Q:=[0,1,0,0]$으로 이루어지고, $G$의 작용은 이 두 점을 서로 바꾼다.

\medskip\noindent
$V_Y\to Y$를 다음과 같은 스킴이라 하자. $Y-P$ 위에서는 $D$를
부풀리고 이어 $C$의 엄밀 변환을 부풀려 정의하고, $Y-Q$ 위에서는
$C$를 부풀리고 이어 $D$의 엄밀 변환을 부풀려 정의한다. (이는 보조정리
\ref{examples-lemma-non-descending-property-projective}의 증명과 같은 구성이지만,
여기서 $Y$는 ${\mathbf P}^3$ 전체가 아니라 ${\mathbf P}^3$의 열린 부분집합을 뜻한다.)
그러면 $Y$ 위의 $G$의 작용은 $V_Y$ 위의 $G$의 작용으로 올라가며,
이는 $Y-P$와 $Y-Q$의 역상들을 서로 바꾼다. $V_Y$ 위의 이 $G$-작용은
$G$-토서 $Y\to S$에 대한 $V_Y$ 위의 강하 데이터
$(V_Y/Y,\varphi_Y)$를 준다. 보조정리
\ref{examples-lemma-non-descending-property-projective}의 증명에서 보였듯 사상
$V_Y\to Y$는 고유이지만 사영은 아니다.

\medskip\noindent
$X$를 열린 부분집합 $Y-P$와 $Y-Q$의 서로소 합이라 하자. 그러면 전사인
에탈 사상 $X\to Y\to S$가 있다. $V$를 $V_Y\to Y$의 $X$로의
당김뒤라 하자. $V_Y\to Y$는 열린 부분집합 $Y-P$, $Y-Q$ 각각의
위에서 부풀리기이므로 $V\to X$는 사영이다. 또한 강하 데이터
$(V_Y/Y,\varphi_Y)$는 에탈 덮개 $X\to S$에 대한 강하 데이터
$(V/X,\varphi)$로 당겨진다.

\medskip\noindent
이 강하 데이터가 스킴의 범주에서 효과적이라고 가정하자. 즉 그 강하
데이터와 함께 $V\to X$로 당겨지는 스킴 $U\to S$가 있다고 하자.
그러면 $U$는 $G$-작용에 의한 $V_Y$의 몫일 것이다.
$$
\xymatrix{
V \ar[r]\ar[d]&  X\ar[d] \\
V_Y \ar[r]\ar[d]& Y\ar[d] \\
U \ar[r]& S
}
$$

\medskip\noindent
$E$를 $C\cup D\subset Y$의 $V_Y$ 안의 역상이라 하자. 따라서
$E\rightarrow C\cup D$는 고유 사상이고, $(C\cup D)-\{P,Q\}$ 위의
섬유들은 ${\mathbf P}^1$과 동형이다. $E$에서 $P$의 역상은 두 직선
$L_0$와 $M_0$의 합집합이다. 따라서 $E$에서 $Q=g(P)$의 역상은 두 직선
$L_0'=g(M_0)$와 $M_0'=g(L_0)$의 합집합이다. 보조정리
\ref{examples-lemma-non-descending-property-projective}의 증명에서 보였듯
$E$ 위에 유리동치 $L_0+M_0'=L_0+g(L_0)\sim 0$이 있다.

\medskip\noindent
닫힌 부분스킴의 강하에 의해 ${\mathbf P}^1$과 동형인 곡선
$L_1\subset U$가 있고 그 $V_Y$ 안의 역상은 $L_0\cup g(L_0)$이다.
(닫힌 몰입이 아핀 사상임에 주의하면서 『강하』의 보조정리
\ref{descent-lemma-affine}를 사용하라.) $R$을 $L_1$의 복소점이라 하자.
$U$가 스킴이라고 가정했으므로 국소환 $O_{U,R}$에서 $R$에서는
소멸하지만 곡선 $L_1$ 전체에서는 소멸하지 않는 함수 $f$를 택할 수 있다.
$D_{\text{loc}}$를 $\Spec O_{U,R}$ 안의 닫힌 부분집합 $\{f = 0\}$의
기약성분이라 하자. 그러면 $D_{\text{loc}}$의 여차원은 1이다.
$U$ 안에서 $D_{\text{loc}}$의 폐포는 $R$을 포함하지만 곡선 $L_1$
전체를 포함하지 않는 $U$의 기약 인자 $D_U$이다. $V_Y$ 안에서 $D_U$의
역상은 $L_0\cup g(L_0)$와 만나지만 곡선 $L_0$, $g(L_0)$ 어느 것도
포함하지 않는 유효 인자 $D$이다.

\medskip\noindent
복소 3차원 다양체 $V_Y$가 매끄러우므로 $O(D)$는 $V_Y$ 위의 선다발이다.
여기서는 정칙 국소환이 유일인수분해정역, 다시 말해 UFD라는 사실을 쓴다.
『대수학에 관하여 더 자세히』의 보조정리
\ref{more-algebra-lemma-regular-local-UFD}를 보라. $O(D)$를 고유 곡면
$E\subset V_Y$에 제한하면, $D$에 관해 아는 사실에 의해 1-순환
$L_0+g(L_0)$ 위에서 양의 차수를 갖는 선다발이 된다. 그러나
$E$ 위에서 $L_0+g(L_0)\sim 0$이므로, 이는 체 위의 고유 스킴에서 선다발의
차수가 유리동치를 법으로 한 1-순환 위에서 잘 정의된다는 사실과 모순이다
(『Chow 호몰로지』의 보조정리
\ref{chow-lemma-proper-pushforward-rational-equivalence}
및 보조정리 \ref{chow-lemma-factors}). 따라서 강하 데이터
$(V/X,\varphi)$는 실제로 효과적이지 않다. 즉 $U$는 스킴으로 존재하지
않는다.
\end{proof}

\noindent
이 예에서 강하 데이터는 대수공간의 범주에서는 {\it 효과적이다}.
더 정확히 말해 $U$는 ${\mathbf C}$ 위의 매끄러운 분리 3차원
대수공간으로 존재한다. 예를 들어 『대수공간』의 보조정리
\ref{spaces-lemma-quotient}에 따른다. Hironaka의 3차원 대수공간 $U$는
매끄러운 준사영 3차원 다양체 $S$를 기약 마디곡선 $(C\cup D)/G$를 따라
부푼 $S'$의 작은 해소이다. 3차원 다양체 $S'$에는 마디 특이점이 하나
있다. $S'$의 다른 작은 해소($U$와 하나의 ``플롭''으로 달라지는 것)도 역시 스킴이
아닌 대수공간이다.




\section{전체 공간이 스킴이 아닌 곡선족}
\label{examples-section-family-of-curves}

\noindent
『몫』의 절 \ref{quot-section-curves}에서 스킴 $S$ 위의 곡선족을 상대
차원 $\leq 1$인, 대수공간 $X$에서 $S$로 가는 고유 평탄 유한 표시
사상으로 정의한다. $S$가 완비 뇌터 국소환의 스펙트럼이면 $X$는
스킴이다. 『공간의 사상에 관한 추가 내용』의 보조정리
\ref{spaces-more-morphisms-lemma-projective-over-complete}를 보라.
이 절에서는 이것이 일반적으로는 참이 아님을 보인다.

\medskip\noindent
$k$를 체라 하자. 다음 성질을 갖는 고유 평탄 사상
$$
Y \longrightarrow \mathbf{A}^1_k
$$
과 점 $y \in Y(k)$에서 시작하자. 이 점은
$0 \in \mathbf{A}^1_k(k)$ 위에 놓이며 다음 성질들을 갖는다.
\begin{enumerate}
\item 섬유 $Y_0$는 $k$ 위의 매끄럽고 기하적으로 기약인 곡선이다.
\item 임의의 진부분 닫힌 부분스킴 $T \subset Y$가
$\mathbf{A}^1_k$를 지배하면 교집합 $T \cap Y_0$는 $y$와 다른 점을 적어도
하나 포함한다.
\end{enumerate}
그러한 곡면이 주어지면 다음과 같이 예를 구성한다.
$$
\xymatrix{
Y \ar[rd] & Z \ar[d] \ar[l] \ar[r] & X \ar[ld] \\
& \mathbf{A}^1_k
}
$$
여기서 $Z \to Y$는 $Y$를 $y$에서 부풀린 것이다. $E \subset Z$를
예외 인자라 하고 $C \subset Z$를 $Y_0$의 엄밀 변환이라 하자.
스킴론적으로 $Z_0 = E \cup C$이다(이를 보려면 $Y$가 $y$에서 매끄럽고,
더욱이 $Y \to \mathbf{A}^1_k$가 $y$에서 매끄럽다는 것을 이용한다).
Artin의 결과들(\cite{ArtinII}; $C$의 법선 다발이 음임을 보이려면
『반안정 축소』의 보조정리 \ref{models-lemma-properties-form}을 이용한다)에
의해 곡선 $C$를 $Z$에서 수축하여 도표와 같은 대수공간 $X$를 얻을 수
있다. $x \in X(k)$를 $C$의 상이라 하자.

\medskip\noindent
$X$가 스킴이 아니라고 주장한다. 실제로 스킴이라면 아핀 열린 근방
$U \subset X$가 있을 것이며, 이는 $x$의 근방이다.
$T = X \setminus U$라 두자. 그러면
$T$는 $\mathbf{A}^1_k$를 지배한다($X \to \mathbf{A}^1_k$의 섬유는
고유이고 차원이 $1$인 반면, $U \to \mathbf{A}^1_k$의 섬유는 아핀이므로
서로 다르다). 닫힌 부분스킴 $T' \subset Z$를 $T$와 동형으로 사상되는
것으로 취하자($x \not \in T$이기 때문이다). 그러면 $T'$의 $X$에서의 상은
위의 조건 (2)와 모순된다($T' \cap Z_0$가 예외 인자 $E$, 즉 부풀리기
$Z \to Y$의 예외 인자에 포함되기 때문이다).

\medskip\noindent
논의를 끝내려면 $Y$를 구성해야 한다. $k$의 표수가 $3$이 아니라고
가정하자. $\mathbf{A}^1_k = \Spec(k[t])$라 쓰고 다음을 취하자.
$$
Y \quad : \quad T_0^3 + T_1^3 + T_2^3 - tT_0T_1T_2 = 0
$$
이는 $\mathbf{P}^2_{k[t]}$ 안에 있다. $t = 0$에서의 섬유는 매끄러운
종수 $1$ 사영 곡선이다.
아핀 조각 $V_+(T_0)$에서는 아핀 방정식
$$
1 + x^3 + y^3 - txy = 0
$$
을 얻으며, 이는 $k$ 위의 매끄러운 곡면을 정의한다. 대칭성에 의해 다른
아핀 조각들에서도 마찬가지이므로 $Y$가 매끄러운 곡면임을 알 수 있다.
마지막으로 이 아핀 방정식에서 분수체가 $k(x, y)$임도 알 수 있으므로
$Y$는 유리 곡면이다. 이제 유리 곡면의 Picard 군은 유한 생성이다
(향후 참고문헌을 여기에 삽입).
따라서 성질 (2)를 갖는 $y \in Y_0(k)$를 택하려면 다음을 만족하는
$y$를 택하는 것으로 충분하다.
\begin{equation}
\label{examples-equation-three}
\mathcal{O}_{Y_0}(ny) \not \in \Im(\Pic(Y) \to \Pic(Y_0))
\text{ 모든 }n > 0
\end{equation}
실제로 차원이 $1$인 기약 성분들의 합을 생각하자. 이들은 (2)에 어긋나는
어떤 $T$의 기약 성분들이다. 그 합은 $Y_0$와의 교집합이 인자 $ny$인
유효 Cartier 인자를 주는데, 여기서
어떤 $n \geq 1$이다. 따라서 $\mathcal{O}_{Y_0}(ny)$가 제한 사상의 상에 속한다고
결론 내릴 수 있다. $Y_0$의 종수가 $\geq 1$이므로 사상
$$
Y_0(k) \to \Pic(Y_0),\quad y \mapsto \mathcal{O}_{Y_0}(y)
$$
은 단사임에 유의하라. 이제 $k$가 비가산 대수적으로 닫힌 체이면
$\Pic(Y)$의 가산성과 방금 한 관찰을 이용하여 어떤 $y \in Y_0(k)$를
찾을 수 있다. 이 점은 (\ref{examples-equation-three})을, 따라서 (2)를 만족한다.

\begin{lemma}
\label{examples-lemma-family-of-curves-not-scheme}
어떤 체 $k$와 곡선족 $X \to \mathbf{A}^1_k$가 존재하여 $X$는
스킴이 아니다.
\end{lemma}

\begin{proof}
위의 논의를 보라.
\end{proof}



\section{유도 밑변환}
\label{examples-section-derived-base-change}

\noindent
$R \to R'$을 환 사상이라 하자. 『대수학에 관한 추가 내용』의 절
\ref{more-algebra-section-derived-base-change}에서 유도 밑변환 함자
$- \otimes_R^\mathbf{L} R' : D(R) \to D(R')$을 구성한다.
다음으로 $R \to A$를 두 번째 환 사상이라 하자. 다음 그림을 보라.
$$
\xymatrix{
	A \ar[r] & A \otimes_R R' \ar@{=}[r] & A' \\
R \ar[u] \ar[r] & R' \ar[u] \ar[ur]
}
$$
$A$-가군 $M$이 주어지면 텐서곱 $M \otimes_R R'$은
$A \otimes_R R'$-가군, 즉 $A'$-가군이다. 환 사상 $A \to A'$에
대해서는 유도 함자
$$
- \otimes_A^\mathbf{L} A' : D(A) \longrightarrow D(A')
$$
가 있지만, 이 함자는 일반적으로 $- \otimes_R^\mathbf{L} R'$과
일치하지 않는다. 더 정확히 말해 $K \in D(A)$에 대한 표준 사상
$$
K \otimes_R^{\mathbf{L}} R' \longrightarrow
K \otimes_A^{\mathbf{L}} A'
$$
은 $D(R')$ 안에 있으며, 『대수학에 관한 추가 내용』의 식
(\ref{more-algebra-equation-comparison-map})에서 구성된다. 이 사상은
일반적으로 동형이 아니다. 따라서 ``유도 밑변환 함자''
$T : D(A) \to D(A')$, 즉 $T(K)$가
$K \otimes_R^\mathbf{L} R'$으로 $D(R')$에서 사상되는 함자가 존재하는지 물을 수
있다. 이 절에서는 일반적으로 그러한 함자가 존재하지 않음을 보인다.

\medskip\noindent
$k$를 체라 하자. $R = k[x, y]$라 두고,
$R' = R/(xy)$ 및 $A = R/(x^2)$라 두자. 대상
$A \otimes_R^\mathbf{L} R'$은 $D(R')$에 속하고
$$
x^2 : R' \longrightarrow R'
$$
로 표현되고,
$H^0(A \otimes_R^\mathbf{L} R') = A \otimes_R R'$이다.
대상 $E$가 $D(A \otimes_R R')$에 속하면서
$A \otimes_R^\mathbf{L} R'$으로 $D(R')$에서 사상되는 일은 없다고
주장한다. 실제로 그러한 $E$가 있다면 가군 $H^0(E)$는 자유여야 하므로
$E$는 $H^0(E)[0] \oplus H^{-1}(E)[1]$로 분해될 것이다. 그러나
$A \otimes_R^\mathbf{L} R'$이 $D(R')$에서 그 코호몰로지 군들의 합과
동형이 아님은 쉽게 알 수 있다.

\begin{lemma}
\label{examples-lemma-no-derived-base-change}
$R \to R'$과 $R \to A$를 환 사상이라 하자. 일반적으로 삼각 범주의
함자 $T : D(A) \to D(A \otimes_R R')$로서, $A$-가군 $M$에 대상
$T(M)$을 대응시키고 이 대상은 $D(A \otimes_R R')$에 속하며
$M \otimes_R^\mathbf{L} R'$으로 사상
$D(A \otimes_R R') \to D(R')$ 아래에서 사상되게 하는 것은 존재하지 않는다.
\end{lemma}

\begin{proof}
위의 논의를 보라.
\end{proof}




\section{흥미로운 콤팩트 대상}
\label{examples-section-interesting-compact}


\noindent
$R$을 환이라 하자. $(A, \text{d})$를 미분 등급 $R$-대수라 하자.
$A = R$이면 $D(A, \text{d}) = D(R)$의 모든 콤팩트 대상은 유한 사영
가군들의 유한 복합체로 표현됨을 알고 있다. 달리 말해 콤팩트 대상은
완전하다. 『대수학에 관한 추가 내용』의 명제
\ref{more-algebra-proposition-perfect-is-compact}를 보라.
미분 등급 가군의 언어에서 이에 대응하는 질문은 다음과 같다.
``$D(A, \text{d})$의 모든 콤팩트 대상은 유한하고 등급 사영인 미분 등급
$A$-가군 $P$로 표현되는가?''

\medskip\noindent
일반적인 미분 등급 대수에 대해서는 이것이 참이 아니다. 실제로
$k$를 표수 $2$인 체라 하자(그러면 부호를 걱정하지 않아도 된다).
$A = k[x, y]/(y^2)$라 하고
\begin{enumerate}
\item $x$의 차수는 $0$,
\item $y$의 차수는 $-1$,
\item $\text{d}(x) = 0$, 그리고
\item $\text{d}(y) = x^2 + x$
\end{enumerate}
라 하자. 그러면 $x : A \to A$는 $K(A, \text{d})$에서 사영자이다.
따라서
$$
A = \Ker(x) \oplus \Im(1 - x)
$$
가 $K(A, \text{d})$에서 성립한다. 『미분 등급 대수』의 보조정리
\ref{dga-lemma-homotopy-direct-sums}와 『유도 범주』의 보조정리
\ref{derived-lemma-projectors-have-images-triangulated}를 보라.
$A$가 $D(A, \text{d})$의 콤팩트 대상임은 분명하다. 그러면
$\Ker(x)$도 $D(A, \text{d})$의 콤팩트 대상이다. 이는 『유도 범주』의
보조정리 \ref{derived-lemma-compact-objects-subcategory}에서 따른다.

\medskip\noindent
다음으로 $M$이 미분 등급 (오른쪽) $A$-가군으로서 $\Ker(x)$를 표현하고,
$M$이 등급 $A$-가군으로서 유한 사영이라고 하자.
$k[x, y]/(y^2)$ 위의 모든 유한 등급 사영 가군은 등급 자유이므로,
$M$은 등급 $k[x, y]/(y^2)$-가군으로서 유한 자유이다(즉 미분을 잊었을
때 그러하다). $N = M/M(x^2 + x)$라 두자. 완전열
$$
0 \to M \xrightarrow{x^2 + x} M \to N \to 0
$$
을 생각하자. $x^2 + x$는 차수 $0$이고 $A$의 중심에 있으며
$\text{d}(x^2 + x) = 0$이므로, 이것은 미분 등급 $A$-가군들의 짧은
완전열이다. 더욱이 $\text{d}(y) = x^2 + x$이므로 $N$ 위의 미분은
선형이다. 사상들
$$
H^{-1}(N) \to H^0(M)
\quad\text{및}\quad
H^0(M) \to H^0(N)
$$
은 동형이다. 실제로 $H^*(M) = H^0(M) = k$인데, 이는
$M \cong \Ker(x)$가 $D(A, \text{d})$에서 성립하기 때문이다. 경계 사상을 계산하면 등급
가군으로서 $H^*(N) = k[x, y]/(x, y^2)$임을 알 수 있다. 자세한 내용은
생략한다. $N$은 자유 $k[x, y]/(y^2, x^2 + x)$-가군이므로 미분과 양립하는
분해
$$
\ldots \to N[2] \xrightarrow{y} N[1] \xrightarrow{y} N \to N/Ny \to 0
$$
를 갖는다. $N$이 유계이고 $H^0(N) = k[x,y]/(x, y^2)$이므로
『호몰로지』의 보조정리 \ref{homology-lemma-first-quadrant-ss}에서
$H^0(N/Ny) = k[x]/(x)$임을 얻는다. 그러나 $N/Ny$는
$k[x]/(x^2 + x) = k \times k$ 위의 유한 자유 가군들로 이루어진 유한
복합체이므로 그 코호몰로지는 짝수 차원이어야 한다. 이는 모순이다.

\begin{lemma}
\label{examples-lemma-no-good-representatif-compact-object}
미분 등급 대수 $(A, \text{d})$와 콤팩트 대상 $E$가 존재한다. 이 대상은
$D(A, \text{d})$에 속하지만, $E$는 유한하고 등급 사영인 미분 등급 $A$-가군으로 표현될 수
없다.
\end{lemma}

\begin{proof}
위의 논의를 보라.
\end{proof}






\section{두 미분 등급 범주}
\label{examples-section-nongraded-differential-graded}

\noindent
이 절에서는 『미분 등급 대수』의 상황 \ref{dga-situation-ABC}에서와 같은
공리 (A), (B), (C)를 만족하지만 대상에 $\mathbf{Z}$-등급이 함께 주어지지
않는 두 미분 등급 범주를 구성한다.

\medskip\noindent
{\bf 예 I.} $X$를 위상 공간이라 하자. $\underline{\mathbf{Z}}$로
값이 $\mathbf{Z}$인 상수층을 나타내자. $A$를
$\underline{\mathbf{Z}}$-토서라 하자. 이 상황에서 아벨 군의 층
$\mathcal{F}$가 {\it $A$-등급}이라는 것은, 국소 단면 $a \in A(U)$가
주어질 때마다 사영자 $p_a : \mathcal{F}|_U \to \mathcal{F}|_U$가 있어서
국소 동형 $\underline{\mathbf{Z}}|_U \to A|_U$가 있을 때마다
$\mathcal{F}|_U = \bigoplus_{n \in \mathbf{Z}} p_n(\mathcal{F})$가
성립한다는 뜻이다. 달리 말해 $X$ 위에서 국소적으로 아벨 층
$\mathcal{F}$는 $\mathbf{Z}$-등급을 갖지만, 겹치는 부분에서 서로 다른
등급의 선택들은 토서 $A$의 전이 함수가 주는 차수의 이동만큼 서로
다르다. 쌍 $(\mathcal{F}, \text{d})$가
{\it 아벨 층들의 $A$-등급 복합체}라는 것은 $\mathcal{F}$가
$A$-등급 아벨 층이고 $\text{d} : \mathcal{F} \to \mathcal{F}$가
미분, 즉 $\text{d}^2 = 0$이며
$p_{a + 1} \circ \text{d} = \text{d} \circ p_a$를 만족한다는 뜻이다.
이는 모든 국소 단면 $a$에 대해 성립하며, 그 단면은 $A$의 단면이다.
달리 말해 $\text{d}(p_a(\mathcal{F}))$는
$p_{a + 1}(\mathcal{F})$에 포함된다.

\medskip\noindent
다음 조건으로 범주 $\mathcal{A}$를 생각하자.
\begin{enumerate}
\item 대상은 아벨 층들의 $A$-등급 복합체이고,
\item 대상 $(\mathcal{F}, \text{d})$, $(\mathcal{G}, \text{d})$에 대해
$$
\Hom_\mathcal{A}((\mathcal{F}, \text{d}), (\mathcal{G}, \text{d})) =
\bigoplus \Hom^n(\mathcal{F}, \mathcal{G})
$$
로 둔다. 여기서 $\Hom^n(\mathcal{F}, \mathcal{G})$는 아벨 층의 사상
$f$들의 군이다. 이 사상들은
$f(p_a(\mathcal{F})) \subset p_{a + n}(\mathcal{G})$를 만족하며, 이는 모든
국소 단면 $a$에 대해 성립하고 그 단면은 $A$의 단면이다. 미분으로는
$\text{d}(f) = \text{d} \circ f - (-1)^n f \circ \text{d}$를 취한다.
『미분 등급 대수』의 예 \ref{dga-example-category-complexes}를 보라.
\end{enumerate}
이것이 실제로 (A), (B), (C)를 만족하는 미분 등급 범주라는 확인은
생략한다. 모든 성질은 $X$ 위에서 국소적으로 확인할 수 있는데, 그곳에서는
아벨 층의 복합체들로 이루어진 미분 등급 범주를 그대로 되찾는다. 따라서
삼각 범주 $K(\mathcal{A})$를 얻는다.

\medskip\noindent
$X$의 꼬인 유도 범주. 대상 $(\mathcal{F}, \text{d})$가 주어지고 이것이
$\mathcal{A}$에 속하면 잘 정의된 $A$-등급 코호몰로지 층
$H(\mathcal{F}, \text{d})$가 있음에 유의하라. 따라서
$K(\mathcal{A})$에서 준동형사상이 무엇을 뜻하는지 분명하다.
준동형사상들을 가역화하여 {\it 유도 범주 $D(\mathcal{A})$, 즉
$A$-등급 층의 복합체들의 유도 범주}를 얻을 수 있다. $A$가 자명한 토서이면
$D(\mathcal{A})$는 $D(X)$와 같지만, 영이 아닌 토서에 대해서는
$X$의 일종의 {\it 꼬인} 유도 범주를 얻는다.

\medskip\noindent
{\bf 예 II.} $C$를 완전체 $k$ 위의 매끄러운 곡선이라 하자. 그 표수는
$2$이다.
그러면 $\Omega_{C/k}$에는 표준 제곱근이 주어진다. 즉
$\Omega_{C/k} = \mathcal{L}^{\otimes 2}$로 쓸 수 있고, 모든 국소 함수
$f$, 즉 $C$ 위의 국소 함수에 대해 단면 $\text{d}(f)$는
$s^{\otimes 2}$와 같으며, 여기서 $s$는 $\mathcal{L}$의 어떤 국소
단면이다. 그 ``이유''는
$$
\text{d}(a_0 + a_1t + \ldots +a_dt^d) =
(\sum\nolimits_{i\text{ 홀수}} a_i^{1/2} t^{(i - 1)/2})^2\text{d}t
$$
이기 때문이다(향후 참고문헌을 여기에 삽입). 특히 이것은 표준 접속
$$
\nabla_{can} :
\Omega_{C/k}
\longrightarrow
\Omega_{C/k} \otimes_{\mathcal{O}_C} \Omega_{C/k}
$$
을 결정하며, 그 $2$-곡률은 영이다(즉 제곱들의 미분이 영이 되게 하는
유일한 접속이다).
접속을 갖는 벡터 다발들의 범주는 텐서 범주이므로 표준 접속
$\nabla_{can}$을 가역층 $\Omega_{C/k}^{\otimes n}$ 위에도 얻으며,
이는 모든 $n \in \mathbf{Z}$에 대해 성립한다.

\medskip\noindent
다음과 같은 범주 $\mathcal{A}$를 취하자.
\begin{enumerate}
\item 대상은 쌍 $(\mathcal{F}, \nabla)$이다. 이 쌍은 유한 국소 자유층
$\mathcal{F}$와 그 위의 접속
$$
\nabla :
\mathcal{F}
\longrightarrow
\mathcal{F} \otimes_{\mathcal{O}_C} \Omega_{C/k}
$$
으로 이루어지며, 이 접속의 $2$-곡률은 영이고,
\item $(\mathcal{F}, \nabla_\mathcal{F})$와
$(\mathcal{G}, \nabla_\mathcal{G})$ 사이의 사상들은
$$
\Hom_\mathcal{A}((\mathcal{F}, \nabla_\mathcal{F}),
(\mathcal{G}, \nabla_\mathcal{G})) =
\bigoplus \Hom_{\mathcal{O}_C}(\mathcal{F},
\mathcal{G} \otimes_{\mathcal{O}_C} \Omega_{C/k}^{\otimes n})
$$
로 주어진다. 원소
$f : \mathcal{F} \to \mathcal{G} \otimes \Omega_{C/k}^{\otimes n}$가
주어지고 그 차수가 $n$이면, 적절한 동일시 아래
$$
\text{d}(f) =
\nabla_{\mathcal{G} \otimes \Omega_{C/k}^{\otimes n}} \circ f +
f \circ \nabla_\mathcal{F}
$$
로 둔다.
\end{enumerate}
이것이 성질 (A), (B), (C)를 갖는 미분 등급 범주를 이룬다는 확인은
생략한다. 따라서 삼각 호모토피 범주 $K(\mathcal{A})$를 얻는다.

\medskip\noindent
$C = \mathbf{P}^1_k$이면 $K(\mathcal{A})$는 영 범주이다. 그러나 $C$가
종수 $> 1$인 매끄러운 고유 곡선이면 $K(\mathcal{A})$는 영이 아니다.
실제로 $\mathcal{N}$이 차수 $0 \leq d < g - 1$이고 영이 아닌 단면
$\sigma$를 갖는 가역층이라고 하자. 그리고
$(\mathcal{F}, \nabla_\mathcal{F}) = (\mathcal{O}_C, \text{d})$ 및
$(\mathcal{G}, \nabla_\mathcal{G}) = (\mathcal{N}^{\otimes 2}, \nabla_{can})$
로 두자. 그러면
$$
\Hom_\mathcal{A}^n((\mathcal{F}, \nabla_\mathcal{F}),
(\mathcal{G}, \nabla_\mathcal{G})) =
\left\{
\begin{matrix}
0 & \text{이면} & n < 0 \\
\Gamma(C, \mathcal{N}^{\otimes 2}) & \text{이면} & n = 0 \\
\Gamma(C, \mathcal{N}^{\otimes 2} \otimes \Omega_{C/k}) & \text{이면} & n = 1
\end{matrix}
\right.
$$
을 알 수 있다. 첫 번째가 $0$인 것은
$\mathcal{N}^{\otimes 2} \otimes \Omega_{C/k}^{\otimes -1}$의 차수가
음이기 때문이다. 이는 조건 $d < g - 1$에서 따른다. 이제 단면
$\sigma^{\otimes 2}$의 미분은 영이므로
준동형 군
$$
\Hom_{K(\mathcal{A})}((\mathcal{F}, \nabla_\mathcal{F}),
(\mathcal{G}, \nabla_\mathcal{G}))
$$
은 영이 아니다.






\section{고유 대수공간의 스택은 대수적이지 않다}
\label{examples-section-proper-spaces-not-algebraic}

\noindent
『몫』의 절 \ref{quot-section-stack-of-spaces}에서 군체들의 스택
$$
p'_{fp, flat, proper} :
\Spacesstack'_{fp, flat, proper}
\longrightarrow
\Sch_{fppf}
$$
을 도입하고 연구했다. 이 스택의 스킴 $S$ 위의 단면 범주는 $S$ 위의
평탄하고 고유하며 유한 표시인 대수공간들의 범주이다. 우리는 이것이
Artin의 공리 중 많은 것을 만족함을 증명했다. 이 절에서는 형식적
유효성이 일반적으로 성립하지 않음을 보여 이 스택이 대수적이지 않은
이유를 보인다.

\medskip\noindent
표준적인 예는 차원 $g$인 아벨 다양체의 보편 변형 공간에는 $g^2$개의
형식 매개변수가 있는 반면, 유효한 형식 변형은 차원
$\leq g(g + 1)/2$인 완비 국소환 위에서 정의될 수 있다는 사실을
이용한다. 우리의 예는 K3 곡면의 적절한 비유효 변형을 명시하여 구성할
것이다. 논증의 개요만 제시하고 모든 세부 사항을 주지는 않겠다.

\medskip\noindent
$k = \mathbf{C}$를 복소수체라 하자. $X \subset \mathbf{P}^3_k$를
차수가 $4$이고 $k$ 위에서 매끄러운 곡면이라 하자. 그러면
$\omega_X \cong \Omega^2_{X/k} \cong \mathcal{O}_X$이다.
마지막으로
$\dim_k H^0(X, T_{X/k}) = 0$,
$\dim_k H^1(X, T_{X/k}) = 20$, 그리고
$\dim_k H^2(X, T_{X/k}) = 0$이다.
$L_{X/k} = \Omega_{X/k}$이다. 이는 $X$가 $k$ 위에서 매끄럽기 때문이다.
$\Ext^i_{\mathcal{O}_X}(\Omega_{X/k}, \mathcal{O}_X) =
H^i(X, T_{X/k})$이며, 『여접공간』의 보조정리
\ref{cotangent-lemma-find-obstruction-ringed-topoi}가 있으므로
$X$의 보편 변형이
$$
k[[x_1, \ldots, x_{20}]]
$$
위에 존재함을 알 수 있다. 이 보편 변형이 유효하다고 가정하자
(『Artin의 공리』의 절 \ref{artin-section-formal-objects}에서의 의미로).
그러면 평탄 고유 사상
$$
f : Y \longrightarrow \Spec(k[[x_1, \ldots, x_{20}]])
$$
을 얻으며, 여기서 $Y$는 보편 변형을 복원하는 대수공간이다.
다음 이유로 이는 불가능하다. $Y$가 분리이므로 아핀 열린 부분스킴
$V \subset Y$를 찾을 수 있다. 특수 섬유 $X$는 $Y$의 섬유이고 매끄러우므로
$f$가 매끄러움을 알 수 있다. 따라서 $Y$는 정칙인 것 위에서
매끄러우므로 정칙이다. 여집합 $D$를 생각하자. 이는 $V$의 여집합이며
$Y$ 안에서 취한다. 그러면 이 여집합은 유효 Cartier 인자이다.
그러면 $\mathcal{O}_Y(D)$는 $\Pic(Y)$의 자명하지 않은
원소이다(이를 증명하려면 체 위의 고유 매끄러운 대수공간에서 공집합이
아닌 아핀 열린집합의 여집합은 언제나 Picard 군에서 자명하지 않음을
보이고, 이를 $f$의 일반 섬유에 적용한다). 마지막으로 모순을 얻기 위해
$\Pic(Y) = 0$임을 보인다. 실제로 사상 $\Pic(Y) \to \Pic(X)$는
단사이다. 이는 $H^1(X, \mathcal{O}_X) = 0$이어서
$\mathcal{O}_X$의 $Y \times \Spec(k[[x_i]]/\mathfrak m^n)$으로의 모든
변형이 자명하고, Grothendieck의 존재 정리가 있기 때문이다(이 정리는
두꺼워짐들 위에서 같은 층을 만드는 연접 가군들이 동형이라고 말한다).
$X$가 충분히 일반적이면 $\Pic(X) = \mathbf{Z}$이고
$\mathcal{O}_X(1)$에 의해 생성된다. 따라서 $\mathcal{O}_X(n)$이 일차
근방으로 변형되지 않음을 보이면 충분하며, 여기서 $n > 0$이다\footnote{이
논증은 사상 $c_1 : \Pic(X) \to H^1(X, \Omega_{X/k})$가 단사인 한
성립한다. 이는 표수 영인 임의의 체 $k$와 임의의 매끄러운 초곡면 $X$에
대해 참이며, 이 초곡면의 차수는 $4$이고 $\mathbf{P}^3_k$ 안에 있다.}.
컵곱
$$
H^1(X, \Omega_{X/k}) \times H^1(X, T_{X/k})
\longrightarrow H^2(X, \mathcal{O}_X)
$$
을 생각하자. 이는 연접 쌍대성에 의해 비퇴화 쌍대화이다.
계산하면 Hodge 코호몰로지에서 Chern 특성류
$c_1(\mathcal{O}_X(n)) \in H^1(X, \Omega_{X/k})$가 영이 아님을 알 수
있다. 따라서 $c_1(\mathcal{O}_X(n))$과의 컵곱이 영이 아닌 일차 변형이
존재한다. 그리고 마지막으로 이 컵곱이 올림에 대한 장애 특성류임을
보인다.

\begin{lemma}
\label{examples-lemma-proper-spaces-not-algebraic}
군체들의 스택
$$
p'_{fp, flat, proper} :
\Spacesstack'_{fp, flat, proper}
\longrightarrow
\Sch_{fppf}
$$
에서 스킴 $S$ 위의 단면 범주가 $S$ 위의 평탄하고 고유하며 유한 표시인
대수공간들의 범주인 것
(『몫』의 절 \ref{quot-section-stack-of-spaces}를 보라)은 대수적 스택이
아니다.
\end{lemma}

\begin{proof}
이것이 대수적 스택이었다면 모든 형식적 대상은 유효했을 것이다.
『Artin의 공리』의 보조정리 \ref{artin-lemma-effective}를 보라.
위의 논의는 $\Spec(\mathbf{C})$로 밑변환한 뒤에는 그렇지 않음을 보인다.
따라서 결론을 얻는다.
\end{proof}






\section{대수적이지 않은 Hom-스택의 예}
\label{examples-section-non-algebraic-hom-stack}

\noindent
$\mathcal{Y}, \mathcal{Z}$를 스킴 $S$ 위의 대수적 스택이라 하자.
{\it Hom-스택} $\underline{\Mor}_S(\mathcal{Y}, \mathcal{Z})$는 $S$ 위의
군체들의 스택으로, 스킴 $T$ 위의 단면 범주는
$$
\Mor_T(\mathcal{Y} \times_S T, \mathcal{Z} \times_S T)
$$
로 주어지며 그 대상은 $1$-사상이고 그 사상은 $2$-사상이다.
이것이 실제로 $(\Sch/S)_{fppf}$ 위의 군체들의 스택이라는 증명은
생략한다(향후 참고문헌을 여기에 삽입). 물론 일반적으로 Hom-스택은
대수적이지 않을 것이다. 이 절에서는 그것이 대수적이지 않으면서
$\mathcal{Y}$는 $S$ 위의 고유 평탄 스킴으로 표현가능하고
$\mathcal{Z}$는 $S$ 위에서 매끄럽고 고유인 예를 제시한다.

\medskip\noindent
$k$를 유한체의 대수적 폐포가 아닌 대수적으로 닫힌 체라 하자.
$S = \Spec(k[[t]])$ 및 $S_n = \Spec(k[t]/(t^n)) \subset S$라 하자.
다음을 만족하는 사상 $f : X \to S$를 취하자.
\begin{enumerate}
\item $f$는 사영이고 평탄하며 그 섬유들은 기하적으로 연결인 곡선이다.
\item 섬유 $X_0 = X \times_S S_0$는 마디 곡선이고, 그 매끄러운 기약
성분들의 쌍대 그래프에는 유리 곡선들로 이루어진 고리가 있다.
\item $X$는 정칙 스킴이다.
\end{enumerate}
그러한 곡면 $X$를 만들기 위해, 예를 들어 다음을 취할 수 있다.
$$
X\quad :\quad T_0T_1T_2 - t(T_0^3 + T_1^3 + T_2^3) = 0
$$
이는 $\mathbf{P}^2_{k[[t]]}$ 안에 있다. $A_0$를 $k$ 위의 영이 아닌 아벨 다양체, 예를 들어
타원곡선이라 하자. $A = A_0 \times_{\Spec(k)} S$를
$S$ 위의 상수 아벨 스킴\footnote{밑 $S$ 위의 아벨 스킴이란 $S$ 위의
평탄 고유 유한 표시 군 스킴으로서 모든 섬유가 아벨 다양체인 것을
말한다.}이라 하자. 이것은 $A_0$에 연관된다. 스택
$\mathcal{X} = \underline{\Mor}_S(X, [S/A]))$가 대수적이지 않음을 보일
것이다.

\medskip\noindent
$[S/A]$는 한편으로 $A$가 $S$에 자명하게 작용하여 얻는 몫 스택이고, 다른
한편으로 fppf $A$-토서를 분류하는 스택과 같다. 『스택의 예』의 명제
\ref{examples-stacks-proposition-equal-quotient-stacks}를 보라.
$[S/A] = [\Spec(k)/A_0] \times_{\Spec(k)} S$임에 유의하라. 이를 이용하면
스킴 $T$ 위의 섬유 범주를 다음과 같이 기술할 수 있다.
\begin{align*}
\mathcal{X}_T
& =
\underline{\Mor}_S(X, [S/A])_T \\
& =
\Mor_T(X \times_S T, [S/A] \times_S T) \\
& =
\Mor_S(X \times_S T, [S/A]) \\
& =
\Mor_{\Spec(k)}(X \times_S T, [\Spec(k)/A_0])
\end{align*}
이는 임의의 $S$-스킴 $T$에 대해 성립한다. 달리 말해 군체
$\mathcal{X}_T$는 fppf $A_0$-토서들의 군체이며, 이 토서들은
$X \times_S T$ 위에 있다.
$\mathcal{X}$가 대수적 스택이 아닌 이유를 논하기 전에 몇 가지
보조정리가 필요하다.

\begin{lemma}
\label{examples-lemma-torsors-over-two-dimensional-regular}
$W$를 함수체가 $K$인 이차원 정칙 정역 뇌터 스킴이라 하자.
$G \to W$를 아벨 스킴이라 하자. 그러면 사상
$H^1_{fppf}(W, G) \to H^1_{fppf}(\Spec(K), G)$는 단사이다.
\end{lemma}

\begin{proof}[증명 개요]
$P \to W$를 일반점에서 자명한 fppf $G$-토서라 하자. 그러면 사상
$\Spec(K) \to P$가 있고 이는 $W$ 위에 있다. 그 스킴론적 상 $Z \subset P$를 취할 수
있다. $P \to W$는 고유이므로($W$ 위의 고유 군 대수공간의 토서이기
때문이다), $Z \to W$는 고유 쌍유리 사상이다. 『체 위의 공간』의
보조정리 \ref{spaces-over-fields-lemma-finite-in-codim-1}에 의해 사상
$Z \to W$는 $W$의 유한 개 닫힌 점들을 제외하면 유한이다.
(곡면에 대해 이차 변환들을 부풀려 사상의 부정점을 해소하는 것에 관한
향후 참고문헌을 여기에 삽입)에 의해 $Z \to W$의 기하적 섬유들의 기약
성분들은 유리 곡선이다. 『공간의 준군에 관한 추가 내용』의 보조정리
\ref{spaces-more-groupoids-lemma-no-nonconstant-morphism-from-P1-to-group}에
의해 유리 곡선에서 그러한 군 스킴이나 토서로 가는 비상수 사상은 없다.
따라서 $Z \to W$는 유한이고, 이에 따라 $Z$는 스킴이며
$Z \to W$는 동형이다. 이는 『사상』의 보조정리
\ref{morphisms-lemma-finite-birational-over-normal}에서 따른다.
달리 말해 토서 $P$는 자명하다.
\end{proof}

\begin{remark}
\label{examples-remark-torsors-over-regular}
보조정리 \ref{examples-lemma-torsors-over-two-dimensional-regular}에서
$2$차원이라는 조건을 없앨 수 있으며, 이 조건은 $W$에 대한 것이다.
즉 $W$가 뇌터 정칙 정역이면
(1) $P(W)\to P(K)$가 전단사이며, 이는 모든 $G$-토서 $P$에 대해 성립하고,
(2) $H^1(W,G)\to H^1(K,G)$가 단사이다.
증명의 개요는 다음과 같다. 먼저 (2)는 (1)에서 분명히 따른다.
(1)에서 단사성은 자명하고, 전사성은 (현재 증명에서처럼) 유리 단면의
스킴론적 폐포가 단면임을 증명하는 문제이다. 이 문제는 $W$ 위에서
fppf-국소적이므로, \'etale 밑변환 뒤에는 $P=G$라 가정해도 된다
(\'etale 덮개는 정칙성을 보존함에 유의하라). 이제
$G=\underline{\mathrm{Pic}}^0_{G'/W}$임을 관찰하라. 여기서 $G'$은 쌍대
아벨 스킴이다. $G'$이 정칙 스킴이므로 일반적으로 정의된 가역층들이
확장된다는 사실에서 결론이 따른다. 사실 ``$W$가 정칙이다''를
``매끄러운 $W$-스킴들이 국소 유일 인수분해적이다''로 약화할 수 있다.
Laurent Moret-Bailly는 Raynaud에게서 Picard 기법을 배웠으며,
Raynaud의 연구 어딘가에 참고문헌이 있을 것이다.
\end{remark}

\begin{lemma}
\label{examples-lemma-torsors-over-field-torsion}
$G$를 체 $K$ 위의 매끄러운 가환 군 대수공간이라 하자. 그러면
$H^1_{fppf}(\Spec(K), G)$는 꼬임 군이다.
\end{lemma}

\begin{proof}
모든 $G$-토서 $P$를 생각하자. 이 토서는 $\Spec(K)$ 위에 있고
$K$ 위에서 매끄러운데, 이는 $G$의 꼴이기 때문이다. 따라서 $P$는
어떤 유한 분리 확대 $L/K$ 위의 점을 갖는다.
$[L : K] = n$이라 하자. $[n](P)$로 $G$-토서를 나타내자. 그 특성류는
$n$배의 것이며, 이때 배를 취하는 원래 특성류는 $P$의 특성류이고
$H^1_{fppf}(\Spec(K), G)$ 안에서 계산한다. 표준 사상
$$
P \times_{\Spec(K)} \ldots \times_{\Spec(K)} P \to [n](P)
$$
이 있으며 이는 $K$ 위의 대수공간의 사상이다. $G$가 아벨이므로 이 사상은 대칭적이다. 따라서 이 사상은 몫
$$
(P \times_{\Spec(K)} \ldots \times_{\Spec(K)} P)/S_n
$$
을 통해 인수분해된다. 한편 사상 $\Spec(L) \to P$는 사상
$$
(\Spec(L) \times_{\Spec(K)} \ldots \times_{\Spec(K)} \Spec(L))/S_n
\longrightarrow (P \times_{\Spec(K)} \ldots \times_{\Spec(K)} P)/S_n
$$
을 정의하며, 왼쪽의 스킴이 $K$-유리점을 가짐은 독자가 확인할 수 있다.
따라서 $[n](P)$가 자명한 토서임을 알 수 있다.
\end{proof}

\noindent
$\mathcal{X} = \underline{\Mor}_S(X, [S/A])$가 대수적 스택이 아님을
증명하려면 『Artin의 공리』의 보조정리 \ref{artin-lemma-effective}에
의해 다음을 보이는 것으로 충분하다.

\begin{lemma}
\label{examples-lemma-not-essentially-surjective}
표준 사상 $\mathcal{X}(S) \to \lim \mathcal{X}(S_n)$은 본질적으로
전사가 아니다.
\end{lemma}

\begin{proof}[증명 개요]
정의들을 풀어 쓰면
$H^1(X, A_0) \to \lim H^1(X_n, A_0)$가 전사가 아님을 확인하는 것으로
충분하다. $X$가 정칙이고 사영이므로 보조정리
\ref{examples-lemma-torsors-over-field-torsion}과
\ref{examples-lemma-torsors-over-two-dimensional-regular}에 의해 각 $A_0$-토서는
$X$ 위에서 꼬임이다. 특히 군 $H^1(X, A_0)$는 꼬임 군이다.
따라서 다음을 보이면 충분하다.
(a) 군 $H^1(X_0, A_0)$는 꼬임 군이 아니고,
(b) 사상 $H^1(X_{n + 1}, A_0) \to H^1(X_n, A_0)$는 전사이며, 이는
모든 $n$에 대해 성립한다.

\medskip\noindent
(a)에 대하여. 꼬임이 아닌 $A_0$-토서 $P_0$를 구성하자. 이는 마디 곡선
$X_0$ 위에 있으며, 자명한 $A_0$-토서들을 $X_0$의 각 성분 위에서 취하고
마디 위에서 꼬임이 아닌 $A_0$의 점들을 동형으로 사용하여 이어 붙인다.
더 정확히 말해
$x \in X_0$를 유리 곡선들로 이루어진 고리에 나타나는 마디라 하자.
$X'_0 \to X_0$를 $X_0$의 정규화라 하되 $X_0 \setminus \{x\}$에서
취하자. 두 점 $x', x'' \in X'_0$를 취하자. 이들은 $x_0$로 사상된다. 그러면
$A_0 \times_{\Spec(k)} X'_0$를 취하고,
$A_0 \times {x'}$와 $A_0 \times \{x''\}$를 동일시한다. 이때 평행이동
$A_0 \to A_0$을 사용하는데, 이는 꼬임이 아닌 점 $a_0 \in A_0(k)$에
의한 것이다(그러한 꼬임이
아닌 점은 $k$가 대수적으로 닫혀 있고 유한체의 대수적 폐포가 아니므로
존재한다. 이는 실제로 증명이 자명하지 않다). 이어 붙인 것이 대수공간임을
(사실 스킴임을) 보일 수 있고, 이것이 꼬임이 아닌 $A_0$-토서임도 보일
수 있으며 그 밑은 $X_0$이다. 이것이 꼬임이 아닌 이유는
$[n](P_0)$가 단면을 가진다면 그 단면이 사상
$s : X'_0 \to A_0$을 만들고, 이 사상은
$[n](a_0) = s(x') - s(x'')$를 만족하기 때문이다. 이 등식은
$A_0(k)$의 군법칙에 관한 것이다. 그러나 고리의 기약 성분들이
유리 곡선이므로 단면 $s$는 그 위에서 상수이다
(『공간의 준군에 관한 추가 내용』의 보조정리
\ref{spaces-more-groupoids-lemma-no-nonconstant-morphism-from-P1-to-group}).
따라서 $s(x') = s(x'')$이고 모순을 얻는다.

\medskip\noindent
(b)에 대하여. 변형 이론에 따르면 $A_0$-토서
$P_n \to X_n$을 $A_0$-토서 $P_{n + 1} \to X_{n + 1}$로 변형하는 데
대한 장애는 $H^2(X_0, \omega)$에 놓인다. 여기서 $\omega$는 $X_0$ 위의
적절한 벡터 다발이다. $X_0$가 곡선이므로 후자는 영이며, 이로써
주장이 증명된다.
\end{proof}

\begin{proposition}
\label{examples-proposition-nonalghomstack}
스택 $\mathcal{X} = \underline{\Mor}_S(X, [S/A])$는 대수적이지 않다.
\end{proposition}

\begin{proof}
위의 논의를 보라.
\end{proof}

\begin{remark}
\label{examples-remark-contradict-aoki}
명제 \ref{examples-proposition-nonalghomstack}는
\cite[Theorem 1.1]{AokiHomStacks}과 모순된다. 문제는
$\underline{\Mor}_S(X, [S/A])$의 형식적 대상들이 유효하지 않다는 데
있다. 같은 문제는 정오표 \cite{AokiHomStacksErr}에도 언급되어 있는데,
이는 \cite{AokiHomStacks}에 대한 정오표이다. 안타깝게도 그 정오표는
$\underline{\Mor}_S(\mathcal{Y}, \mathcal{Z})$가 대수적이라고 계속
주장하며, 그 조건은 $\mathcal{Z}$가 분리라는 것이다. 그러나 $[S/A]$가 분리이므로 이 역시 명제
\ref{examples-proposition-nonalghomstack}와 모순된다.
\end{remark}


\section{강한 형식적 유효성을 만족하지 않는 대수적 스택}
\label{examples-section-non-formal-effectiveness}

\noindent
이것은 \cite[Example 4.12]{Bhatt-Algebraize}이다.
$k$를 대수적으로 닫힌 체라 하자. $A$를 $k$ 위의 아벨 다양체라 하자.
$A(k)$가 꼬임 군이 아니라고 가정하자($k$가 유한체의 대수적 폐포가
아니면 이는 항상 성립한다). $\mathcal{X} = [\Spec(k)/A]$라 하자.
어떤 아이디얼 $I \subset k[x, y]$가 존재하여
$$
\mathcal{X}_{\Spec(k[x, y]^\wedge)}
\longrightarrow
\lim \mathcal{X}_{\Spec(k[x, y]/I^n)}
$$
가 본질적으로 전사가 아니라고 주장한다. 실제로 $I$를
$xy(x + y - 1)$로 생성되는 아이디얼이라 하자. 그러면
$X_0 = V(I)$는 세 점에서 삼각형 모양으로 이어 붙인
$\mathbf{A}^1_k$의 세 복사본으로 이루어진다. 따라서 무한 위수 토서
$P_0$를 $A$에 대해 만들 수 있다. 이 토서는 $X_0$ 위에 있다. 이를 위해
$X_0$의 기약 성분들 위에서 자명한 토서를 취하고 꼬임이 아닌 점들에
의한 평행이동을 사용해 이어 붙인다. 보조정리
\ref{examples-lemma-not-essentially-surjective}의
증명에서와 꼭 마찬가지로 $P_0$를 토서 $P_n$으로 올릴 수 있으며,
후자는 $X_n = \Spec(k[x, y]/I^n)$ 위에 있다.
$k[x, y]^\wedge$는 이차원 정칙 정역이므로 모든 토서 $P$는, $A$에 대한
것이고 $\Spec(k[x, y]^\wedge)$ 위에 있을 때, 꼬임이다
(보조정리 \ref{examples-lemma-torsors-over-two-dimensional-regular}과
\ref{examples-lemma-torsors-over-field-torsion}). 따라서 이 토서들의 체계는
표시된 함자의 상에 속하지 않는다.

\begin{lemma}
\label{examples-lemma-non-formal-effectiveness}
$k$를 유한체의 폐포가 아닌 대수적으로 닫힌 체라 하자. $A$를 $k$ 위의
아벨 다양체라 하자. $\mathcal{X} = [\Spec(k)/A]$라 하자.
전사 전이 사상들을 갖고 그 핵들이 국소 멱영인 $k$-대수들의 역계
$R_n$과, 체계 $(\xi_n)$이 존재한다. 후자는 $\mathcal{X}$의 체계이며
체계 $(\Spec(R_n))$ 위에 놓이고, 『Artin의 공리』의 주석
\ref{artin-remark-strong-effectiveness}의 의미에서 유효하지 않다.
\end{lemma}

\begin{proof}
위의 논의를 보라.
\end{proof}








\section{Grothendieck의 존재 정리에 대한 반례}
\label{examples-section-Grothendieck-existence}

\noindent
$k$를 체라 하고 $A = k[[t]]$라 하자. $X$를
$U = \Spec(A[x])$와 $V = \Spec(A[y])$를 다음 동일시
$$
U \setminus \{0_U\} \longrightarrow V \setminus \{0_V\}
$$
로 붙여 얻은 것으로 하자. 이 동일시는 $x$를 $y$로 보내며, 여기서
$0_U \in U$와 $O_V \in V$는 극대 아이디얼 $(x, t)$와 $(y, t)$에
대응하는 점들이다. $A_n = A/(t^n)$라 하고
$X_n = X \times_{\Spec(A)} \Spec(A_n)$라 하자.
$\mathcal{F}_n$을 $A_n[x]$-가군 $A_n[x]/(x) \cong A_n$과
$A_n[y]$-가군 $0$에 자명한 붙이기를 적용하여 얻는 $X_n$ 위의
연접층이라 하자. $\mathcal{I} \subset \mathcal{O}_X$를 $t$로
생성되는 아이디얼 층이라 하자. 그러면 $(\mathcal{F}_n)$은
스킴의 코호몰로지, 절
\ref{coherent-section-existence-proper-support}에서 정의한 범주
$\textit{Coh}_{\text{지지가 그 위에서 고유한 } A}(X, \mathcal{I})$
의 대상이다. 한편 이 대상은 스킴의 코호몰로지, 식
(\ref{coherent-equation-completion-functor-proper-over-A})의 함자의
상에 속하지 않는다. 실제로 상에 속한다면 유한 $A[x]$-가군 $M$,
유한 $A[y]$-가군 $N$, 그리고 동형사상 $M[1/t] \cong N[1/t]$가
존재하여 모든 $n$에 대해 $M/t^nM \cong A_n[x]/(x)$ 및
$N/t^nN = 0$이 성립해야 한다. 이것이 불가능함은 쉽게 알 수 있다.

\begin{lemma}
\label{examples-lemma-counter-Grothendieck-existence}
연접층의 대수화에 대한 다음 반례들이 존재한다.
\begin{enumerate}
\item 스킴의 코호몰로지, 정리
\ref{coherent-theorem-grothendieck-existence}에서 진술한
Grothendieck의 존재 정리에서 $X \to \Spec(A)$가 분리라는 가정을
제거하면 정리는 거짓이다.
\item Quot, 정리 \ref{quot-theorem-coherent-algebraic-general} 및
\ref{quot-theorem-coherent-algebraic}의 연접층 스택
$\Cohstack_{X/B}$에서 $X \to S$가 분리라는 가정을 제거하면,
일반적으로 이 스택은 대수적이지 않다.
\item Quot, 명제 \ref{quot-proposition-quot}의 함자
$\Quotfunctor_{\mathcal{F}/X/B}$에서 $X \to B$가 분리라는 가정을
제거하면, 일반적으로 이 함자는 대수공간이 아니다.
\end{enumerate}
\end{lemma}

\begin{proof}
(1)은 위에서 보았다. 이는 $\textit{Coh}_{X/A}$가 Artin의 공리, 절
\ref{artin-section-axioms}의 공리 [4]를 만족하지 않음을 보인다.
따라서 Artin의 공리, 보조정리 \ref{artin-lemma-effective}에 의해
대수적 스택일 수 없다. 이로써 (2)가 참임을 알 수 있다. (3)을
보려면 모든 $n$에 대해 서로 양립하는 전사
$\mathcal{O}_{X_n} \to \mathcal{F}_n$이 존재함에 주목하라.
따라서 $\Quotfunctor_{\mathcal{O}_X/X/A}$가 공리 [4]를 만족하지
않으며, 앞과 같이 (3)이 참임을 알 수 있다.
\end{proof}

\section{아핀 형식적 대수공간}
\label{examples-section-affine-formal-algebraic-space}

\noindent
$K$를 체라 하고 $(V_i)_{i \in I}$를 $K$ 위의 영이 아닌 벡터공간들로
이루어진 유향 역계라 하자. 전이 사상들은 전사이고 $\lim V_i = 0$이라
가정한다. 절 \ref{examples-section-zero-limit}을 보라. $V_i$가 제곱이 영인
아이디얼이 되도록 $K$-대수로서 $R_i = K \oplus V_i$라 하자. 그러면
$R_i$는 멱영 핵을 갖는 전사 전이 사상들로 이루어진 $K$-대수의
역계이고 $\lim R_i = K$이다. 아핀 형식적 대수공간
$X = \colim \Spec(R_i)$는 McQuillan이 아닌 아핀 형식적 대수공간의
예이다.

\begin{lemma}
\label{examples-lemma-affine-not-mcquillan}
McQuillan이 아닌 아핀 형식적 대수공간이 존재한다.
\end{lemma}

\begin{proof}
위의 논의를 보라.
\end{proof}

\noindent
절 \ref{examples-section-zero-limit}에서와 같은 완전열의 계
$0 \to W_i \to V_i \to K \to 0$를 잡자. 다음과 같이 놓는다.
$A_i = K[V_i]/(ww'; w, w' \in W_i)$.
그러면 멱영 핵을 갖는 서로 양립하는 전사들의 계
$A_i \to K[t]$가 존재하고, 전이 사상 $A_i \to A_j$ 역시 멱영 핵을
갖는 전사이다. $V_i$가 $s \in S_i$로 주어지는 기저를 갖는 자유
$K$-가군임을 상기하자. 그러면 $K$의 표수가 영일 때 $A_i$의 $d$차
성분은 $s^d$, $s \in S_i$로 주어지는 기저를 갖는 자유 $K$-가군이고,
그 각각은 $t^d$로 사상된다. 따라서 $A_i$들의 $d$차 성분으로 이루어진
역계는 벡터공간 $V_i$들의 역계와 동형이다. $\lim V_i = 0$이므로,
적어도 $K$의 표수가 영일 때 $\lim A_i = K$임을 얻는다. 이는
``정칙함수''들이 점들을 분리하지 못하는 아핀 형식적 대수공간의
예를 준다.

\begin{lemma}
\label{examples-lemma-affine-formal-functions-do-not-separate-points}
다음 의미에서 정칙함수들이 점들을 분리하지 못하는 아핀 형식적
대수공간 $X$가 존재한다. 형식적 공간, 정의
\ref{formal-spaces-definition-affine-formal-algebraic-space}에서와 같이
$X = \colim X_\lambda$로 쓰면,
$\lim \Gamma(X_\lambda, \mathcal{O}_{X_\lambda})$는 체이지만
$X_{red}$는 무한히 많은 점을 갖는다.
\end{lemma}

\begin{proof}
위의 논의를 보라.
\end{proof}

\noindent
$K$, $I$, $(V_i)$는 위와 같다고 하자. 다음과 같은 계들을 생각하자.
$$
\Phi = (\Lambda, J_i \subset \Lambda, (M_i) \to (V_i))
$$
여기서 $\Lambda$는 증대된 $K$-대수이고, 각 $i \in I$에 대해
$J_i \subset \Lambda$는 제곱이 영인 아이디얼이며,
$(M_i) \to (V_i)$는 $K$-벡터공간 역계들의 사상으로서 각 $i$에 대해
$M_i \to V_i$가 전사이다. 또한 $M_i$는 $\Lambda$-가군 구조를 가지며,
$i > j$일 때 전이 사상 $M_i \to M_j$는 $\Lambda$-선형이고
$i > j$에 대해 $J_j M_i \subset \Ker(M_i \to M_j)$이다.
주장: 모든 $i > j$에 대해 $M_j = M_i/J_j M_i$인 위와 같은 계가
존재한다.

\medskip\noindent
주장이 참이면, 아핀 형식적 대수공간들의 표현가능 사상
$$
\colim_{i \in I} \Spec(\Lambda/J_i \oplus M_i)
\longrightarrow
\text{Spf}(\lim \Lambda/J_i)
$$
을 얻는데, 그 원천은 McQuillan이 아니지만 표적은 McQuillan이다.
여기서 $\Lambda/J_i \oplus M_i$에는 $M_i$가 제곱이 영인 아이디얼이
되는 통상적인 $\Lambda/J_i$-대수 구조를 준다. 표현가능성은 정확히
$i > j$일 때 $M_i/J_jM_i = M_j$라는 조건으로 옮겨진다. 사상의
원천이 McQuillan이 아닌 이유는 사영
$\lim_{i \in I} M_i \to M_i$가 전사가 아니기 때문이다. 실제로
사상 $\lim V_i \to V_i$가 전사가 아니고 전사 $M_i \to V_i$가
있으므로 이것이 성립한다. 몇몇 세부사항은 생략한다.

\medskip\noindent
주장의 증명. 우선 적어도 하나의 계가 존재함에 주목하자. 즉,
$$
\Phi_0 = (K, J_i = (0), (V_i) \xrightarrow{\text{항등}} (V_i))
$$
이다. 계 $\Phi$가 주어졌다고 하자. 그러면 계의 사상
$\Phi \to \Phi'$가 존재하여(계의 사상은 자명한 방식으로 정의한다)
$\Ker(M_i/J_j M_i \to M_j)$가 $M'_i/J'_j M'_i$에서 영으로 사상됨을
보이겠다. 이것을 보이고 나면 $n \geq 1$에 대해 귀납적으로
$\Phi_n = (\Phi_{n - 1})'$로 두고 $\Phi = \colim \Phi_n$을 취하는
통상적인 방법으로 원하는 성질을 갖는 계를 얻을 수 있다.
세부사항은 생략한다.

\medskip\noindent
$\Phi$가 주어졌을 때 $\Phi'$의 구성. $i > j$이고
$\xi \in \Ker(M_i \to M_j)$인 삼중항 $u = (i, j, \xi)$들의 집합을
$U$라 하자. $s(i, j, \xi) = i$ 및 $t(i, j, \xi) = j$로 주어지는
사상들을 $s, t : U \to I$로 나타내자. 그리고 $u$의 세 번째 성분을
$\xi_u \in M_{s(u)}$로 둔다. 다음을 취한다.
$$
\Lambda' = \Lambda[x_u; u \in U]/(x_u x_{u'}; u, u' \in U)
$$
증대 $\Lambda' \to K$는 $\Lambda$의 증대로 주고 $x_u$를 영으로
보내도록 한다. 다음과 같이 둔다.
$J'_k = J_k \Lambda' + (x_{u,\ t(u) \geq k})$.
또 다음과 같이 둔다.
$$
M'_i = M_i \oplus \bigoplus\nolimits_{s(u) \geq i} K\epsilon_{i, u}
$$
$i > j$일 때 전이 사상 $M'_i \to M'_j$로는 주어진 사상
$M_i \to M_j$를 사용하고 $\epsilon_{i, u}$를 $\epsilon_{j, u}$로
보낸다. 사상 $M'_i \to V_i$는 주어진 사상 $M_i \to V_i$를 유도하고
$\epsilon_{i, u}$를 영으로 보낸다. 마지막으로 $\Lambda'$가
$M'_i$에 다음과 같이 작용하게 한다. $\lambda \in \Lambda$에 대해서는
$M_i$ 위에서 $\Lambda$-가군 구조로 작용하고, $\epsilon_{i, u}$ 위에서는
증대 $\Lambda \to K$를 통해 작용한다. 원소 $x_u$는 모든 $i$에 대해
$M_i$ 위에서 $0$으로 작용한다. 마지막으로 다음과 같이 정의한다.
$$
x_u \epsilon_{i, u} = \text{원소 }\xi_u\text{의 상, 공역 }M_i
$$
그리고 그 밖의 모든 곱 $x_{u'} \epsilon_{i, u}$는 영으로 둔다.
표시한 공식은 $s(u) \geq i$이고 $\xi_u \in M_{s(u)}$이므로 의미가
있다. 확인해야 할 주된 사항은 $i > j$일 때 $J'_j M'_i \subset M'_i$가
$M'_j$에서 영으로 사상된다는 것과 $\Ker(M_i \to M_j)$가
$M'_i/J_j M'_i$에서 영으로 사상된다는 것이다. 마지막 사실은 임의의
$\xi \in \Ker(M_i \to M_j)$에 대해
$\xi = x_{(i, j, \xi)} \epsilon_{i, (i, j, \xi)} \in J'_j M'_i$이기
때문이다. 세부사항은 생략한다.

\begin{lemma}
\label{examples-lemma-representable-morphism-affine-formal-not-mcquillan-top}
아핀 형식적 대수공간들의 표현가능 사상 $f : X \to Y$로서
$Y$는 McQuillan이지만 $X$는 McQuillan이 아닌 것이 존재한다.
\end{lemma}

\begin{proof}
위의 논의를 보라.
\end{proof}

\section{평탄 사상은 유한 표시 평탄 사상들의 유향 쌍대극한이 아니다}
\label{examples-section-flat-not-colimit-flat-finitely-presented}

\noindent
이 절의 목표는 평탄하고 유한 표시인 환 사상들의 여과 쌍대극한이
아닌 평탄 환 사상의 예를 주는 것이다. \cite{gabber-nonexcellent}에서는
$A$가 차원이 $1$이고 잉여 표수가 영인 비우수 국소환이면 완비화로
가는 (평탄) 환 사상 $A \to A^\wedge$가 유한형 평탄 환 사상들의
여과 쌍대극한이 아님을 보인다. 이 절의 예에서는 원천이 우수환이다.
독자에게 다른 예를 제출해 주기를 권한다. 예가 있다면
\href{mailto:stacks.project@gmail.com}{stacks.project@gmail.com}으로
전자우편을 보내기 바란다.

\medskip\noindent
구성을 위해 소수 $p$를 고정하고
$A = \mathbf{F}_p[x_1, \ldots, x_n]$라 하자. $A$의 절대 정수적 폐포
$A^+$를 선택하자. 즉, $A^+$는 그 분수체의 한 대수적 폐포에서
$A$의 정수적 폐포이다. \cite[\S 6.7]{HHBigCM}에서
$A \to A^+$가 평탄임을 보인다.

\medskip\noindent
$n \geq 3$이면 $A$-대수 $A^+$는 유한 표시 평탄 $A$-대수들의
여과 쌍대극한이 아니라고 주장한다.

\medskip\noindent
$n = 3$인 경우의 논증을 개략적으로 제시하고, 더 큰 $n$으로의
일반화는 독자에게 맡긴다. 원점에서 $A$를 엄밀 헨젤화한 것을 $R$라
하고 그 절대 정수적 폐포를 $R^+$라 할 때, 사상 $R \to R^+$에 대한
유사한 명제를 증명하면 충분하다. $R$는 잉여체 $k$가
$\mathbf{F}_p$의 대수적 폐포인 헨젤 정칙 국소환임에 주목하라.

\medskip\noindent
$k$ 위의 보통 아벨 곡면 $X$와 $X$ 위의 매우 풍부한 선다발 $L$을
선택하자. 절단환
$\Gamma_*(X, L) = \bigoplus_n H^0(X,L^n)$은 $L$에 관한 $X$ 위의
아핀 원뿔의 좌표환이다. $L$이 충분히 양이면 이는 정규환이다.
$S$를 원뿔의 꼭짓점에서 $\Gamma_*(X, L)$을 헨젤화한 것이라 하자.
그러면 $S$는 차원이 $3$인 헨젤 뇌터 정규 정역이다. 유한형
$k$-대수 $\Gamma_*(X, L)$의 한 뇌터 정규화를 헨젤화하여 유한 단사
사상 $R \to S$를 얻는다. $R^+$가 $R$의 절대 정수적 폐포이므로
매입 $S \to R^+$도 고정할 수 있다. 따라서 $R^+$는 $S$의 절대 정수적
폐포이기도 하다. $R^+$가 평탄 $R$-대수들의 여과 쌍대극한이 아님을
보이려면 다음을 보이면 충분하다.
\begin{enumerate}
\item $S \to P$가 환 사상이고 $P$가 $R$ 위에서 평탄 유한형이면,
$T$가 $R$ 위에서 유한 평탄인 환 사상 $S \to T$가 존재한다.
\item 임의의 유한 환 사상 $S \to T$에 대해 환 $T$는 $R$-평탄이
아니다.
\end{enumerate}
실제로 $S$는 $R$ 위에서 유한 표시이므로, 만일 이를 유한 표시
평탄 $R$-대수 $P_i$들의 여과 쌍대극한 $R^+ = \colim_i P_i$로 쓸 수
있다면 $i \gg 0$에 대해 $S \to R^+$는 $S \to P_i \to R^+$로
분해될 것이다. 이는 위의 두 주장과 모순이다. 주장 (1)은 $R$가
헨젤이라는 사실과 절단 논증으로부터 따른다. 사상에 관한 더 많은
내용, 보조정리
\ref{more-morphisms-lemma-qf-fp-flat-neighbourhood-dominates-fppf}를 보라.
(2)는 \cite{BhattSmallCMMod}에서 증명되었다. 독자의 편의를 위해
그 논증을 상기한다.

\medskip\noindent
$U \subset \Spec(S)$를 천공 스펙트럼이라 하자. 그러면 자연스러운
사상들 $X \leftarrow U \subset \Spec(S)$가 있다. 첫 번째 사상은
$H^1(U, \mathcal{O}_U) \simeq H^1(X, \mathcal{O}_X)$라는 동일시를
준다. 완전화의 비트 벡터로 옮기고 Artin--Schreier 열을
사용하면\footnote{여기서는 $S$가 표수 $p$인 엄밀 헨젤 국소환이므로
$S \to S$, $f \mapsto f^p - f$가 전사라는 사실을 사용한다.
또한 $S$는 정규 정역이므로
$\Gamma(U, \mathcal{O}_U) = S$이다. 따라서
$H^1_\etale(U, \mathbf{Z}/p)$는 $f \mapsto f^p - f$가 유도하는 사상
$H^1(U, \mathcal{O}_U) \to H^1(U, \mathcal{O}_U)$의 핵이다.},
$H^1_\etale(U, \mathbf{Z}_p) \simeq H^1_\etale(X, \mathbf{Z}_p)$라는
동일시를 얻는다. 특히 이 군은 계수가 $2$인 유한 자유
$\mathbf{Z}_p$-가군이다($X$가 보통이기 때문이다). 모순을 얻기 위해
위 (2)에서와 같은 $R$-평탄 $T$가 존재한다고 가정하자.
$V \subset \Spec(T)$를 $U$의 원상이라 하고 유도된 유한 전사 사상을
$f : V \to U$라 쓰자. $U$가 정규이므로 $U_\etale$ 위에 대각합 사상
$f_*\mathbf{Z}_p \to \mathbf{Z}_p$가 있으며, 이 사상과 당김
$\mathbf{Z}_p \to f_*\mathbf{Z}_p$의 합성은
$d = \deg(f)$를 곱하는 사상이다. 코호몰로지로 옮기고
$H^1_\etale(U, \mathbf{Z}_p)$가 꼬임 가군이 아니라는 사실을 사용하면
$H^1_\etale(V, \mathbf{Z}_p)$가 영이 아님을 얻는다.
$R^1\lim$의 간섭이 없어서
$H^1_\etale(V, \mathbf{Z}_p) \simeq \lim H^1_\etale(V, \mathbf{Z}/p^n)$
이므로 군 $H^1(V_\etale,\mathbf{Z}/p)$는 영이 아니어야 한다.
$T$가 $R$-평탄이므로 $\Gamma(V, \mathcal{O}_V) = T$이고 이는 엄밀
헨젤이다. Artin--Schreier 열에 의해
$H^1(V, \mathcal{O}_V) \neq 0$이다. 이는
$H^2_\mathfrak m(T) \neq 0$과 동치이며, 여기서
$\mathfrak m \subset R$는 극대 아이디얼이다. 그러나 $T$는
$R$-가군으로서 유한 평탄(즉, 유한 자유)이고
$H^2_\mathfrak m(R) = 0$이므로 모순이다. 이 모순으로 (2)가 증명된다.

\begin{lemma}
\label{examples-lemma-weird-flat-map}
가환환 $A$와 평탄 $A$-대수 $B$로서 유한 표시 평탄 $A$-대수들의
여과 쌍대극한으로 쓸 수 없는 것이 존재한다. 실제로 $A$를 유한형
$\mathbf{F}_p$-대수로 선택할 수도 있고, 잉여체의 표수가 $0$인
$1$차원 뇌터 국소환으로 선택할 수도 있다.
\end{lemma}

\begin{proof}
위의 논의를 보라.
\end{proof}

\section{꼬임 가군들로 나눈 가군의 범주}
\label{examples-section-serre-quotient-modulo-torsion-modules}

\noindent
꼬임 군의 범주는 모든 아벨 군의 범주의 Serre 부분범주이다(호몰로지,
정의 \ref{homology-definition-serre-subcategory}). 더 일반적으로 임의의
환 $A$에 대해 꼬임 $A$-가군의 범주는 모든 $A$-가군의 범주의 Serre
부분범주이다. 대수학에 관한 더 많은 내용, 절
\ref{more-algebra-section-abelian-categories-modules}을 보라. $A$가
정역이면 몫범주(호몰로지, 보조정리
\ref{homology-lemma-serre-subcategory-is-kernel})는 분수체 위의
벡터공간의 범주와 동치이다. 이는 다음의 더 일반적인 명제로부터
따른다.

\begin{proposition}
\label{examples-proposition-localization-and-serre-quotients}
$A$를 환이라 하고 $S$를 $A$의 곱셈적 부분집합이라 하자.
$\text{Mod}_A$는 $A$-가군의 범주를 나타내고, $\mathcal{T}$는 모든
원소가 $S$의 어떤 원소에 의해 소멸되는 가군들로 이루어진 그 Serre
부분범주라 하자. 그러면 표준적 동치
$\text{Mod}_A/\mathcal{T} \rightarrow \text{Mod}_{S^{-1}A}$가 있다.
\end{proposition}

\begin{proof}
$M \mapsto M \otimes_A S^{-1}A$로 주어지는 함자
$\text{Mod}_A \to \text{Mod}_{S^{-1}A}$는 완전하고(대수학, 명제
\ref{algebra-proposition-localization-exact}) $\mathcal{T}$의 가군들을
영으로 보낸다. 따라서 호몰로지, 보조정리
\ref{homology-lemma-serre-subcategory-is-kernel}에서 주어진 보편성에
의해 이 함자는 함자
$\text{Mod}_A/\mathcal{T} \to \text{Mod}_{S^{-1}A}$로 내려간다.

\medskip\noindent
거꾸로, $M \otimes_A S^{-1}A = 0$인 임의의 $A$-가군 $M$은
$\mathcal{T}$의 대상이다. 실제로
$M \otimes_A S^{-1}A \cong S^{-1} M$이다(대수학, 보조정리
\ref{algebra-lemma-tensor-localization}). 따라서 호몰로지, 보조정리
\ref{homology-lemma-quotient-by-kernel-exact-functor}에 의해 함자
$\text{Mod}_A/\mathcal{T} \to \text{Mod}_{S^{-1}A}$는 충실하다.

\medskip\noindent
더 나아가 이 매입은 본질적으로 전사이다. $S^{-1}A$-가군 $N$의
원상으로 $N_A$, 즉 $N$을 $A$-가군으로 본 것을 취할 수 있다. 실제로
$x \otimes a/s$를 $(a/s) \cdot x$로 보내는 표준 사상
$N_A \otimes_A S^{-1}A \to N$은 $S^{-1}A$-가군의 동형사상이다.
\end{proof}

\begin{proposition}
\label{examples-proposition-quotient-by-torsion-modules}
$A$를 환이라 하자. $Q(A)$는 그 전분수환(대수학, 예
\ref{algebra-example-localize-at-prime}에서와 같은 것)을 나타낸다고
하자. $\text{Mod}_A$는 $A$-가군의 범주를 나타내고, $\mathcal{T}$는
꼬임 가군들로 이루어진 그 Serre 부분범주라 하자.
$\text{Mod}_{Q(A)}$는 $Q(A)$-가군의 범주를 나타낸다고 하자. 그러면
표준적 동치
$\text{Mod}_A/\mathcal{T} \rightarrow \text{Mod}_{Q(A)}$가 있다.
\end{proposition}

\begin{proof}
명제 \ref{examples-proposition-localization-and-serre-quotients}를 곱셈적
부분집합
$S = \{f \in A \mid f \text{가 영인자가 아닌 환 }A\}$에
적용하면 곧바로 따른다. 가군이 꼬임 가군일 필요충분조건은 그 모든
원소 각각이 $S$의 어떤 원소에 의해 소멸되는 것이기 때문이다.
\end{proof}

\begin{proposition}
\label{examples-proposition-quotient-by-finitely-generated-torsion-modules}
$A$를 뇌터 정역이라 하고 $K$를 그 분수체라 하자.
$\text{Mod}_A^{fg}$는 유한 생성 $A$-가군의 범주를 나타내고,
$\mathcal{T}^{fg}$는 유한 생성 꼬임 가군들로 이루어진 그 Serre
부분범주라 하자. 그러면 $\text{Mod}_A^{fg}/\mathcal{T}^{fg}$는
유한 차원 $K$-벡터공간의 범주와 표준적으로 동치이다.
\end{proposition}

\begin{proof}
명제 \ref{examples-proposition-quotient-by-torsion-modules}에서 주어진 동치는
매입 $\text{Mod}_A^{fg}/\mathcal{T}^{fg} \to
\text{Mod}_A/\mathcal{T}$를 따라 제한되어 동치
$\text{Mod}_A^{fg}/\mathcal{T}^{fg} \to \text{Vect}_K^{fd}$를 준다.
뇌터 가정은 $\text{Mod}_A^{fg}$가 아벨 범주임(대수학에 관한 더 많은
내용, 절 \ref{more-algebra-section-abelian-categories-modules}을 보라)을
보장하고, 표준 함자
$\text{Mod}_A^{fg}/\mathcal{T}^{fg} \to \text{Mod}_A/\mathcal{T}$가
충만함을 보장한다(그렇지 않으면 유한 생성 가군의 꼬임 부분가군이
$\mathcal{T}^{fg}$의 대상이 아닐 수도 있다).
\end{proof}

\begin{proposition}
\label{examples-proposition-quotient-abelian-groups-by-torsion-groups}
아벨 군의 범주를 그 꼬임 군들의 Serre 부분범주로 나눈 몫은
$\mathbf{Q}$-벡터공간의 범주이다.
\end{proposition}

\begin{proof}
주장은 명제 \ref{examples-proposition-quotient-by-torsion-modules}에서 바로
따른다.
\end{proof}

\section{서로 다른 쌍대극한 위상}
\label{examples-section-colimit-topology}

\noindent
이 예는 \cite[예 1.2, 553쪽]{TSH}이다. $n \geq 1$에 대해
$G_n = \mathbf{Q} \times \mathbf{R}^n$라 하고, 이를 덧셈에 관한
위상군으로 보되 통상적인 (유클리드) 위상을 부여하자.
$(x_0, \ldots, x_n)$을 $(x_0, \ldots, x_n, 0)$으로 보내는 닫힌
매입 $G_n \to G_{n + 1}$을 생각하자. 위상
$$
U \subset G\text{가 열림} \Leftrightarrow G_n \cap U\text{가 열림 }\forall n
$$
을 부여한 $G = \colim G_n$은 위상군이 아니라고 주장한다.

\medskip\noindent
이를 보기 위해 다음 집합을 생각하자.
$$
U = \{(x_0, x_1, x_2, \ldots)\text{ 중 다음을 만족하는 것: }
|x_j| < |\cos(jx_0)| \text{ 모든 } j > 0\}
$$
$\pi$는 유리수가 아니므로 $jx_0$는 결코 $\pi/2$의 정수배가 아니라는
사실을 사용하면 $U \cap G_n$이 열림을 쉽게 보일 수 있다.
$0 \in U$이므로, 위의 위상이 $G$를 위상군으로 만든다면 $0$의 열린
근방 $V \subset G$로서 $V + V \subset U$인 것이 존재할 것이다.
그러면 모든 $j \geq 0$에 대해 $\epsilon_j > 0$이 존재하여
$|x_j| < \epsilon_j$이면
$(0, \ldots, 0, x_j, 0, \ldots) \in V$일 것이다.
$V + V \subset U$이므로 $|x_0| < \epsilon_0$ 및
$|x_j| < \epsilon_j$일 때
$$
(x_0, 0, \ldots, 0, x_j, 0, \ldots) \in U
$$
여야 한다. 그러나 $j \epsilon_0 > \pi/2$가 되도록 $j$를 충분히
크게 잡으면 $|\cos(jx_0)|$가 $\epsilon_j$보다 작도록
$x_0 \in \mathbf{Q}$를 선택할 수 있고, 따라서
$|\cos(jx_0)| < |x_j| < \epsilon_j$인 $x_j$가 존재한다. 이 모순으로
주장이 증명된다.

\begin{lemma}
\label{examples-lemma-colimit-topology}
(아벨) 위상군들의 계 $G_1 \to G_2 \to G_3 \to \ldots$로서 위상공간의
범주에서 취한 $\colim G_n$과 위상군의 범주에서 취한 $\colim G_n$이
서로 다른 것이 존재한다.
\end{lemma}

\begin{proof}
위의 논의를 보라.
\end{proof}

\section{보편적으로 몫위상적이지만 V-덮개가 아닌 사상}
\label{examples-section-universally-submersive-not-V}

\noindent
$A$를 값매김환이라 하자. $\mathfrak p \subset A$를 극소 소 아이디얼도
극대 아이디얼도 아닌 소 아이디얼이라 하자. (염두에 두기 좋은 경우는
$A$가 세 개의 소 아이디얼을 갖고 $\mathfrak p$가 그 ``가운데''에
있는 경우이다.) 아핀 스킴들의 사상
$$
\Spec(A_\mathfrak p) \amalg \Spec(A/\mathfrak p) \longrightarrow \Spec(A)
$$
을 생각하자. 이것이 보편적으로 몫위상적이라고 주장한다. 이를
증명하기 위해 환 사상 $A \to B$로 주어진 아핀 스킴의 사상
$\Spec(B) \to \Spec(A)$를 잡자. 그러면 다음 사상이 몫위상적임을
보여야 한다.
$$
\Spec(B_\mathfrak p) \amalg \Spec(B/\mathfrak pB) \to \Spec(B)
$$
우선 이 사상은 전사이다. 다음으로 $T \subset \Spec(B)$가 부분집합이고
$T_1 = \Spec(B_\mathfrak p) \cap T$ 및
$T_2 = \Spec(B/\mathfrak p B) \cap T$가 닫혀 있다고 하자.
$T_1$과 $T_2$가 모두 $B$-대수의 스펙트럼이므로 $T$가 어떤
$B$-대수의 스펙트럼의 상임을 알 수 있다. 따라서 $T$가 닫혔음을
보이려면 $T$가 특수화에 대해 안정적임을 보이면 충분하다. 대수학,
보조정리 \ref{algebra-lemma-image-stable-specialization-closed}를 보라.
이를 위해 $p \leadsto q$가 $\Spec(B)$의 점들의 특수화이고
$p \in T$라고 하자. $A'$를 값매김환이라 하고, $\Spec(A')$의 일반점
$\eta$가 $p$로 사상되고 $\Spec(A')$의 닫힌점 $s$가 $q$로 사상되는 사상
$\Spec(A') \to \Spec(B)$를 잡자. 스킴, 보조정리
\ref{schemes-lemma-points-specialize}를 보라. 합성
$\gamma : \Spec(A') \to \Spec(A)$의 상은 정확히
$\gamma(\eta) \leadsto \xi \leadsto \gamma(s)$인 점
$\xi \in \Spec(A)$들의 집합임에 주목하라(세부사항은 생략한다).
$\mathfrak p \not \in \Im(\gamma)$이면
$p, q \in \Spec(B_\mathfrak p)$이거나
$p, q \in \Spec(B/\mathfrak pB)$임을 알 수 있다. 이 경우 $T_1$, 각각\ $T_2$가
닫혔다는 사실로부터 $q \in T_1$, 각각\ $q \in T_2$를 얻고,
따라서 $q \in T$이다. 마지막으로
$\mathfrak p \in \Im(\gamma)$라고 하고, 이를
$\mathfrak p = \gamma(r)$로 쓰자. 그러면 특수화
$p \leadsto r$ 및 $r \leadsto q$가 있다. 이 경우
$p, r \in \Spec(B_\mathfrak p)$이고
$r, q \in \Spec(B/\mathfrak pB)$이다. 따라서 먼저
$r \in T_1 \subset T$를 얻고, $r$가 $\mathfrak p$로 사상되므로
$r \in T_2$를 얻은 뒤, 원하는 대로 $q \in T_2 \subset T$를 얻는다.

\medskip\noindent
한편 한 원소로 된 족
$$
\{\Spec(A_\mathfrak p) \amalg \Spec(A/\mathfrak p) \longrightarrow \Spec(A)\}
$$
은 V-덮개가 아니라고 주장한다. 위상, 정의
\ref{topologies-definition-V-covering}을 보라. 실제로 이것이
V-덮개라면 값매김환의 확대 $A \subset B$가 존재하여
$\Spec(B) \to \Spec(A)$가
$\Spec(A_\mathfrak p) \amalg \Spec(A/\mathfrak p)$를 통해 분해될
것이다. 그러면 $\Spec(A')$이 연결되지 않아야 하는데, 이는 불합리하다.

\begin{lemma}
\label{examples-lemma-universally-submersive-not-V}
아핀 스킴들의 사상 $X \to Y$로서 보편적으로 몫위상적이지만
$\{X \to Y\}$가 V-덮개가 아닌 것이 존재한다.
\end{lemma}

\begin{proof}
위의 논의를 보라.
\end{proof}

\section{정수환의 스펙트럼은 준콤팩트하지 않다}
\label{examples-section-canonical}

\noindent
물론 이 절의 제목은 정수환의 스펙트럼을 위상공간으로 본 것을
가리키지 않는다. 모든 스펙트럼은 위상공간으로서 준콤팩트하기
때문이다(대수학, 보조정리 \ref{algebra-lemma-quasi-compact}).
그 대신 이 제목은 스킴의 범주 위의 표준 위상에서 정수환의
스펙트럼을 본 것과 사이트에서 준콤팩트 대상의 정의를 가리킨다
(사이트, 정의 \ref{sites-definition-quasi-compact}).

\medskip\noindent
소수들의 집합 $P$ 위의 비주요 초필터를 $U$라 하자.
부분집합 $T \subset P$에 대해 그 여집합을 $T^c = P \setminus T$로
나타낸다. $A \in U$에 대해 $S_A \subset \mathbf{Z}$를
$p \in A$들로 생성되는 곱셈적 부분집합이라 하자. 다음과 같이 둔다.
$$
\mathbf{Z}_A = S_A^{-1}\mathbf{Z}
$$
$P$를 $\Spec(\mathbf{Z})$의 닫힌점들의 집합으로 생각하면
$\Spec(\mathbf{Z}_A) = \{(0)\} \cup A^c \subset \Spec(\mathbf{Z})$
임에 주목하라. $A, B \in U$이면 $A \cap B \in U$이고
$A \cup B \in U$이며 다음이 성립한다.
$$
\mathbf{Z}_{A \cap B} =
\mathbf{Z}_A \times_{\mathbf{Z}_{A \cup B}} \mathbf{Z}_B
$$
(환들의 섬유곱이다.) 특히 임의의 정수 $n$ 및 원소들
$A_1, \ldots, A_n \in U$에 대해 사상
$$
\Spec(\mathbf{Z}_{A_1}) \amalg \ldots \amalg \Spec(\mathbf{Z}_{A_n})
\longrightarrow \Spec(\mathbf{Z})
$$
은 어떤 $p$에 대해 $\Spec(\mathbf{Z}[1/p])$를 통해 분해된다
(즉, 임의의 $p \in A_1 \cap \ldots \cap A_n$에 대해 그러하다).
따라서 평탄 사상들의 족
$\{\Spec(\mathbf{Z}_A) \to \Spec(\mathbf{Z})\}_{A \in U}$는 공동으로
전사이지만 그 어떤 유한 부분족도 공동으로 전사가 아니다.

\medskip\noindent
$\mathbf{Z}$-가군 $M$에 대해 다음과 같이 둔다.
$$
M_A = S_A^{-1}M = M \otimes_{\mathbf{Z}} \mathbf{Z}_A
$$
주장 I: 모든 $\mathbf{Z}$-가군 $M$에 대해 다음이 성립한다.
$$
M =
\text{동등자}\left(
\xymatrix{
\prod\nolimits_{A \in U} M_A \ar@<1ex>[r] \ar@<-1ex>[r] &
\prod\nolimits_{A, B \in U} M_{A \cup B}
}
\right)
$$
먼저 $M$이 꼬임이 없다고 가정하자. 그러면 모든 $A \in U$에 대해
$M_A \subset M_P$이다. 따라서 다음을 증명해야 함을 알 수 있다.
$$
M = \bigcap\nolimits_{A \in U} M_A\text{, 교집합은 }M_P = M \otimes \mathbf{Q}
$$
실제로 $U$가 비주요이므로 임의의 소수 $p$에 대해
$\{p\}^c \in U$이다. 또한 $M_{\{p\}^c} = M_{(p)}$는 소 아이디얼
$(p)$에서의 국소화와 같다. 따라서 이미
$M_{(2)} \cap M_{(3)} = M$이므로 위의 등식은 자명하다.
다음으로 $M$이 꼬임이라고 가정하자. 그러면
$$
M = \bigoplus\nolimits_{p \in P} M[p^\infty]
$$
이고, 이에 따라 다음이 성립한다.
$$
M_A = \bigoplus\nolimits_{p \not \in A} M[p^\infty]
$$
이는 $A$의 소수들에서 국소화하고 있기 때문이다.
$(x_A) \in \prod M_A$가 동등자에 속한다고 하자.
$x_p = x_{\{p\}^c} \in M[p^{\infty}]$로 나타내자. 동등자 성질은
다음을 뜻한다.
$$
x_A = (x_p)_{p \not \in A}
$$
특히 이 성질은 $p \not \in A$인 것들 중 유한 개를 제외한 모든
소수에 대해 $x_p$가 영임을 뜻한다. 꼬임인 경우의 증명을 끝내려면
유한 개의 소수 $p$를 제외한 모든 소수에 대해 $x_p$가 영임을 보이면
충분하다. 그렇지 않다면
$\{p \in P \mid x_p \not = 0\} = T \amalg T'$로 쓰되, 오른쪽은 두
무한 집합의 서로소 합집합으로 쓰자. 그러면 $U$가 초필터이므로
$T \not \in U$ 또는 $T' \not \in U$이다(실제로 둘 다 $U$에 속하면
$U$가 $T \cap T' = \emptyset$을 포함하게 되는데, 이는 허용되지
않는다). $T, T'$ 가운데 $T \not \in U$라고 하자. 그러면
$T = A^c$이고 이는 위에서
말한 유한성과 모순이다. 마지막으로 $M$이 일반적인 가군이라고 하자.
그러면 짧은 완전열
$$
0 \to M_{tors} \to M \to M/M_{tors} \to 0
$$
과 다음 큰 도식을 살펴본다.
$$
\xymatrix{
M_{tors} \ar[r] \ar[d] &
\prod\nolimits_{A \in U} M_{tors, A} \ar@<1ex>[r] \ar@<-1ex>[r] \ar[d] &
\prod\nolimits_{A, B \in U} M_{tors, A \cup B} \ar[d] \\
M \ar[r] \ar[d] &
\prod\nolimits_{A \in U} M_A \ar@<1ex>[r] \ar@<-1ex>[r] \ar[d] &
\prod\nolimits_{A, B \in U} M_{A \cup B} \ar[d] \\
M/M_{tors} \ar[r] &
\prod\nolimits_{A \in U} (M/M_{tors})_A \ar@<1ex>[r] \ar@<-1ex>[r] &
\prod\nolimits_{A, B \in U} (M/M_{tors})_{A \cup B} \\
}
$$
열들의 완전성과 꼬임 가군 $M_{tors}$ 및 꼬임 없는 가군
$M/M_{tors}$에 대해 얻은 결과를 사용하여 도식 추적을 하면 $M$에
대한 주장 I이 증명된다. 이는 다음 보조정리의 현상에 대한 예를 준다.

\begin{lemma}
\label{examples-lemma-non-fpqc-descent}
환 $A$와 평탄 환 사상들의 무한족
$\{A \to A_i\}_{i \in I}$로서 모든 $A$-가군 $M$에 대해
$$
M =
\text{동등자}\left(
\xymatrix{
\prod\nolimits_{i \in I} M \otimes_A A_i \ar@<1ex>[r] \ar@<-1ex>[r] &
\prod\nolimits_{i, j \in I} M \otimes_A A_i \otimes_A A_j
}
\right)
$$
가 성립하지만 같은 사실이 성립하는 유한 부분족은 없는 것이
존재한다.
\end{lemma}

\begin{proof}
위의 논의를 보라.
\end{proof}

\noindent
소수들의 집합 $P$ 위의 비주요 초필터 $U$를 계속 사용한다.
$R$를 환이라 하자. $A \in U$에 대해
$R_A = S_A^{-1}R = R \otimes \mathbf{Z}_A$로 나타내자.
주장 II: 닫힌 부분집합 $T_A \subset \Spec(R_A)$, $A \in U$가
주어지고 모든 $A, B \in U$에 대해
$$
(\Spec(R_{A \cup B}) \to \Spec(R_A))^{-1}T_A =
(\Spec(R_{A \cup B}) \to \Spec(R_B))^{-1}T_B
$$
가 성립하면, 모든 $A \in U$에 대해
$T_A = (\Spec(R_A) \to \Spec(R))^{-1}(T)$인 닫힌 부분집합
$T \subset \Spec(R)$가 존재한다. $A \in U$에 대해 $T_A$를 잘라내는
근기 아이디얼을 $I_A \subset R_A$라 하자. 그러면 붙이기 조건은
모든 $A, B \in U$에 대해 $R_{A \cup B}$ 안에서
$S_{A \cup B}^{-1}I_A = S_{A \cup B}^{-1}I_B$임을 함의한다
(국소화는 근기 아이디얼이라는 성질을 보존하기 때문이다).
$I' \subset R$를 $I_P \subset R_P = R \otimes \mathbf{Q}$ 안으로
사상되는 원소들의 집합이라 하자. 그러면
$A \in U$에 대해 다음을 알 수 있다.
\begin{enumerate}
\item $I_A \subset I'_A = S_A^{-1}I'$이고,
\item $M_A = I'_A/I_A$는 꼬임 가군이다.
\end{enumerate}
물론 $A, B \in U$에 대해 표준적 동일시
$S_{A \cup B}^{-1}M_A = S_{A \cup B}^{-1}M_B$를 얻는다.
꼬임 가군 $M_A$들을 그 $p$-일차 성분들로 분해하면, 위에서 주어진
표준적 동일시들과 양립하면서
$$
M_A = \bigoplus\nolimits_{p \not \in A} M_p
$$
가 되도록 하는 $p$-멱 꼬임 $R$-가군 $M_p$들이 존재함을 독자가
쉽게 보일 수 있다. $M = \bigoplus_{p \in P} M_p$로 두면 위의 표준적
동일시들과 양립하는 표준 동형사상 $M_A = S_A^{-1}M$을 얻는다.
그러면 모든 $A \in U$에 대해 사상 $I_A \to M_A$를 복원하는
$R$-가군의 표준 사상
$$
I' \longrightarrow M
$$
을 얻는다. 이는 위에서 말한 모든 양립성과, $M$이
$\prod_{A \in U} M_A$에서 $\prod_{A, B \in U} M_{A \cup B}$로 가는
두 사상의 동등자라는 앞서 증명한 주장으로부터 따른다.
$I = \Ker(I' \to M)$라 하자. 그러면 $I$는 아이디얼이고
$T = V(I)$는 모든 $A \in U$에 대해 닫힌 부분집합 $T_A$를 복원하는
닫힌 부분집합이다. 이로써 주장 II가 증명된다.

\begin{lemma}
\label{examples-lemma-Z-not-quasi-compact}
스킴 $\Spec(\mathbf{Z})$는 스킴의 범주 위의 표준 위상에서
준콤팩트하지 않다.
\end{lemma}

\begin{proof}
위의 표기법을 사용하여 사상들의 족
$$
\mathcal{W} = \{\Spec(\mathbf{Z}_A) \to \Spec(\mathbf{Z})\}_{A \in U}
$$
을 생각하자. 내림, 보조정리
\ref{descent-lemma-universal-effective-epimorphism}와 위에서 증명한 두
주장에 의해 이는 보편 유효 전사이다. 섬유곱을 갖는 임의의 범주에서
보편 유효 전사들은 $\mathcal{C}$에 표준 위상을 정의하는 사이트의
구조를 준다(쉽게 해결할 수 있는 몇몇 집합론적 문제를 제외하고).
따라서 $\mathcal{W}$는 표준 위상에 대한 덮개이다. 한편 위에서 임의의
유한 부분족
$$
\{\Spec(\mathbf{Z}_{A_i}) \to \Spec(\mathbf{Z})\}_{i = 1, \ldots, n},\quad
n \in \mathbf{N}, A_1, \ldots, A_n \in U
$$
이 어떤 $p$에 대해 $\Spec(\mathbf{Z}[1/p])$를 통해 분해됨을 보았다.
따라서 이 유한족은 보편 유효 전사일 수 없고, 더 일반적으로 어떤
보편 유효 전사 $\{g_j : T_j \to \Spec(\mathbf{Z})\}$도
$\{\Spec(\mathbf{Z}_{A_i}) \to \Spec(\mathbf{Z})\}_{i = 1, \ldots, n}$을
세분할 수 없다. 사이트, 정의 \ref{sites-definition-quasi-compact}에
의해 이는 $\Spec(\mathbf{Z})$가 표준 위상에서 준콤팩트하지 않음을
뜻한다. 우리의 준콤팩트성 개념이 통상적인 토포스론적 정의와
일치한다는 것은 사이트, 보조정리
\ref{sites-lemma-quasi-compact}을 보라.
\end{proof}
